
% =======================================================================
% APPENDIX C: Detailergebnisse der Monte-Carlo Simulation (Orthogonalität)
% =======================================================================
\section{Orthogonalitäts-Analysen der Design-Modifikationen}
\label{app:c}

In diesem Anhang sind ergänzende Visualisierungen und Analysen zur Bewertung der geometrischen Design-Modifikationen bzw. Orthogonalitäts-Abweichungen aufgeführt.

\subsection{Orthogonalitäts-Analysen aus Abschnitt~\ref{subsubsec:quantiplore}}

\begin{table}[htbp]
    \centering
    \caption{Simulations\-ergebnisse: Schätzvarianz (\ac{MSE}) der einzelnen Koeffizienten bei $\sym{snr}=2{,}4$.}
    \label{tab:results_ortho_mse}
    \renewcommand{\arraystretch}{1.2}
    \small
    \begin{tabular*}{\textwidth}{@{\extracolsep{\fill}}l rrrrrr@{}}
        \toprule
        \textbf{Szenario (Störfall)} & \textbf{$\hat{\sym{beta}}_{0}$} & \textbf{$\hat{\sym{beta}}_{1}$} & \textbf{$\hat{\sym{beta}}_{2}$} & \textbf{$\hat{\sym{beta}}_{12}$} & \textbf{$\hat{\sym{beta}}_{11}$} & \textbf{$\hat{\sym{beta}}_{22}$} \\
        \midrule
        Referenz (Orthogonal) & $3{,}09$ & $3{,}10$ & $3{,}17$ & $6{,}17$ & $3{,}13$ & $3{,}11$ \\
        Fall I: Shift ($-20\,\%$) & $3{,}05$ & $3{,}57$ & $3{,}15$ & $6{,}20$ & $4{,}15$ & $3{,}17$ \\
        Fall I: Shift ($+20\,\%$) & $3{,}07$ & $2{,}94$ & $3{,}14$ & $6{,}37$ & $2{,}21$ & $3{,}06$ \\
        Fall I: \ac{FC} & $3{,}06$ & $3{,}82$ & $3{,}10$ & $6{,}30$ & $4{,}61$ & $3{,}21$ \\
        Fall II: Korrelation (leicht) & $3{,}06$ & $3{,}17$ & $3{,}07$ & $6{,}22$ & $3{,}13$ & $3{,}12$ \\
        Fall II: Korrelation (stark) & $3{,}10$ & $3{,}59$ & $3{,}01$ & $5{,}21$ & $3{,}83$ & $3{,}44$ \\
        Fall III: Scatter ($\pm 10\,\%$) & $3{,}04$ & $3{,}18$ & $2{,}98$ & $6{,}20$ & $3{,}27$ & $2{,}62$ \\
        Fall IV: \ac{SP} & $3{,}18$ & $5{,}29$ & $3{,}07$ & $6{,}17$ & $5{,}46$ & $3{,}47$ \\
        Fall IV: \ac{ZP} & $3{,}53$ & $3{,}09$ & $3{,}13$ & $6{,}29$ & $3{,}24$ & $3{,}25$ \\
        \bottomrule
    \end{tabular*}
\end{table}


\begin{table}[htbp]
    \centering
    \caption{Simulations\-ergebnisse: Durchschnittliche Trennschärfe-Änderung ($\Delta \sym{power}$) der Koeffizienten. Ein negatives Vorzeichen indiziert einen Verlust, ein positives Vorzeichen einen Gewinn an $\sym{power}$ verglichen zur Referenz (Basis \ac{CCD})}
    \label{tab:results_ortho_power}
    \renewcommand{\arraystretch}{1.2}
    \small
    \begin{tabular*}{\textwidth}{@{\extracolsep{\fill}}l rrrrrr@{}}
        \toprule
        \textbf{Szenario (Störfall)} & \textbf{$\sym{power}_{\sym{beta}_{0}}$} & \textbf{$\sym{power}_{\sym{beta}_{1}}$} & \textbf{$\sym{power}_{\sym{beta}_{2}}$} & \textbf{$\sym{power}_{\sym{beta}_{12}}$} & \textbf{$\sym{power}_{\sym{beta}_{11}}$} & \textbf{$\sym{power}_{\sym{beta}_{22}}$} \\
        \midrule
        Referenz (Orthogonal) & $-$ & $-$ & $-$ & $-$ & $-$ & $-$ \\
        Fall I: Shift ($-20\,\%$) & $\pm 0{,}0\,\%$ & $-0{,}9\,\%$ & $-0{,}0\,\%$ & $\pm 0{,}0\,\%$ & $-4{,}5\,\%$ & $-0{,}1\,\%$ \\
        Fall I: Shift ($+20\,\%$) & $-0{,}0\,\%$ & $+0{,}3\,\%$ & $-0{,}0\,\%$ & $-0{,}0\,\%$ & $+4{,}6\,\%$ & $-0{,}1\,\%$ \\
        Fall I: \ac{FC} & $\pm 0{,}0\,\%$ & $-1{,}4\,\%$ & $-0{,}1\,\%$ & $+0{,}1\,\%$ & $-6{,}7\,\%$ & $-0{,}2\,\%$ \\
        Fall II: Korrelation (leicht) & $\pm 0{,}0\,\%$ & $-0{,}1\,\%$ & $-0{,}0\,\%$ & $+0{,}1\,\%$ & $+0{,}1\,\%$ & $-0{,}0\,\%$ \\
        Fall II: Korrelation (stark) & $\pm 0{,}0\,\%$ & $-1{,}2\,\%$ & $+0{,}3\,\%$ & $+3{,}3\,\%$ & $-3{,}1\,\%$ & $-1{,}5\,\%$ \\
        Fall III: Scatter ($\pm 10\,\%$) & $\pm 0{,}0\,\%$ & $-0{,}1\,\%$ & $+0{,}4\,\%$ & $+0{,}2\,\%$ & $-1{,}2\,\%$ & $+2{,}8\,\%$ \\
        Fall IV: \ac{SP} & $\pm 0{,}0\,\%$ & $-4{,}4\,\%$ & $-0{,}1\,\%$ & $-0{,}3\,\%$ & $-10{,}6\,\%$ & $-1{,}5\,\%$ \\
        Fall IV: \ac{ZP} & $-0{,}0\,\%$ & $-0{,}3\,\%$ & $-0{,}2\,\%$ & $-0{,}8\,\%$ & $-1{,}3\,\%$ & $-0{,}9\,\%$ \\
        \bottomrule
    \end{tabular*}
\end{table}


\begin{table}[htbp]
    \centering
    \caption{Simulations\-ergebnisse: Relativer systematischer Schätzfehler (Bias) der einzelnen Koeffizienten in Prozent bezogen auf ihren Nominalwert bei $\sym{snr}=1{,}20$.}
    \label{tab:results_ortho_bias}
    \renewcommand{\arraystretch}{1.2}
    \small
    \begin{tabular*}{\textwidth}{@{\extracolsep{\fill}}l rrrrrr@{}}
        \toprule
        \textbf{Szenario (Störfall)} & \textbf{$\hat{\sym{beta}}_{0}$} & \textbf{$\hat{\sym{beta}}_{1}$} & \textbf{$\hat{\sym{beta}}_{2}$} & \textbf{$\hat{\sym{beta}}_{12}$} & \textbf{$\hat{\sym{beta}}_{11}$} & \textbf{$\hat{\sym{beta}}_{22}$} \\
        \midrule
        Referenz (Orthogonal) & $+0{,}05\,\%$ & $-0{,}01\,\%$ & $-0{,}06\,\%$ & $-0{,}45\,\%$ & $+0{,}17\,\%$ & $+0{,}05\,\%$ \\
        Fall I: Shift ($-20\,\%$) & $+0{,}04\,\%$ & $-0{,}00\,\%$ & $-0{,}09\,\%$ & $-0{,}06\,\%$ & $-0{,}02\,\%$ & $+0{,}31\,\%$ \\
        Fall I: Shift ($+20\,\%$) & $+0{,}01\,\%$ & $-0{,}11\,\%$ & $+0{,}05\,\%$ & $-0{,}24\,\%$ & $+0{,}08\,\%$ & $-0{,}00\,\%$ \\
        Fall I: \ac{FC} & $+0{,}08\,\%$ & $+0{,}13\,\%$ & $-0{,}02\,\%$ & $-0{,}25\,\%$ & $+0{,}13\,\%$ & $-0{,}13\,\%$ \\
        Fall II: Korrelation (leicht) & $-0{,}08\,\%$ & $+0{,}13\,\%$ & $+0{,}08\,\%$ & $-0{,}19\,\%$ & $-0{,}03\,\%$ & $-0{,}23\,\%$ \\
        Fall II: Korrelation (stark) & $-0{,}08\,\%$ & $+0{,}01\,\%$ & $-0{,}27\,\%$ & $+0{,}55\,\%$ & $-0{,}71\,\%$ & $+0{,}17\,\%$ \\
        Fall III: Scatter ($\pm 10\,\%$) & $+0{,}02\,\%$ & $+0{,}09\,\%$ & $+0{,}19\,\%$ & $+0{,}14\,\%$ & $+0{,}10\,\%$ & $+0{,}09\,\%$ \\
        Fall IV: \ac{SP} & $-0{,}03\,\%$ & $+0{,}14\,\%$ & $+0{,}01\,\%$ & $-0{,}56\,\%$ & $-0{,}57\,\%$ & $-0{,}07\,\%$ \\
        Fall IV: \ac{ZP} & $-0{,}05\,\%$ & $-0{,}13\,\%$ & $-0{,}20\,\%$ & $+0{,}01\,\%$ & $+0{,}13\,\%$ & $+0{,}07\,\%$ \\
        \bottomrule
    \end{tabular*}
\end{table}



\subsubsection{Algorithmus zur Berechnung der Exposure-Area}
\begin{algorithm}[htbp]
    \caption{Algorithmische Identifikation der Anker-Exposition (\textit{Exposure-Area})}
    \label{alg:exposure_area}
    \begin{algorithmic}[1]
        \Require Basis-Design $\sym{X}_{\text{base}}$, polynomialer Prognosevektor $f(\sym{x}_{\text{pr}})$,
        diskreter Suchraum $\mathcal{X}$, Anzahl umzuwidmender \acp{AP} $n_{\text{AP}}$, relative Toleranzmarge $q$

        \State $\nu_{\text{min}} \gets \infty$ \Comment{Initialisierung der minimalen Prädiktionsvarianz}
        \State $\mathcal{C} \gets \text{Kombinationen}(\mathcal{X}, n_{\text{AP}})$ \Comment{Menge aller möglichen \ac{AP}-Kombinationen}
        \State $\sym{V} \gets$ Array der Länge $|\mathcal{C}|$ zur Speicherung der Varianzen

        \vspace{0.5em}
        \Statex \textit{// Phase 1: Exhaustive Evaluierung der \ac{SPV} im Prognosepunkt}
        \For{\textbf{each} $C_i \in \mathcal{C}$}
        \State $\sym{X}_{\text{aug}} \gets \sym{X}_{\text{base}} \cup C_i$ \Comment{Augmentierung des Designs}
        \State $\sym{M}_i \gets \sym{X}_{\text{aug}}^\top \sym{X}_{\text{aug}}$ \Comment{Informationsmatrix}
        \State $\nu_i \gets N \cdot f(\sym{x}_{\text{pr}})^\top \sym{M}_i^{-1} f(\sym{x}_{\text{pr}})$ \Comment{Berechnung der \ac{SPV} gem. Gl. X.X}
        \State $\sym{V}[i] \gets \nu_i$
        \If{$\nu_i < \nu_{\text{min}}$}
        \State $\nu_{\text{min}} \gets \nu_i$ \Comment{Registrierung des absoluten Minimums}
        \EndIf
        \EndFor

        \vspace{0.5em}
        \Statex \textit{// Phase 2: Filterung und Konstruktion des Toleranzraums}
        \State $\nu_{\text{thresh}} \gets \nu_{\text{min}} \cdot (1 + q)$ \Comment{Definition der oberen Toleranzschranke}
        \State $\mathcal{X}_{\text{exposure}} \gets \emptyset$

        \For{\textbf{each} $C_i \in \mathcal{C}$}
        \If{$\sym{V}[i] \le \nu_{\text{thresh}}$}
        \State $\mathcal{X}_{\text{exposure}} \gets \mathcal{X}_{\text{exposure}} \cup C_i$ \Comment{Akkumulation valider Stützstellen}
        \EndIf
        \EndFor

        \Ensure Toleranzraum $\mathcal{X}_{\text{exposure}}$ (Menge aller zulässigen geometrischen Verortungen)
    \end{algorithmic}
\end{algorithm}


% \begin{figure}[htbp]
%     \centering
%     \setlength\figurewidth{0.9\textwidth}
%     \setlength\figureheight{4cm}
%     % This file was created by matlab2tikz.
%
\definecolor{mycolor1}{rgb}{0.12941,0.12941,0.12941}%
%
\begin{tikzpicture}

\begin{axis}[%
width=0.951\figurewidth,
height=\figureheight,
at={(0\figurewidth,0\figureheight)},
scale only axis,
bar shift auto,
xmin=0.54,
xmax=6.46,
xtick={1.00,2.00,3.00,4.00,5.00,6.00},
xticklabels={{$\sym{beta}_0$},{$\sym{beta}_1$},{$\sym{beta}_2$},{$\sym{beta}_{12}$},{$\sym{beta}_{11}$},{$\sym{beta}_{22}$}},
ymin=-10.00,
ymax=25.00,
ylabel style={font=\color{mycolor1}},
ylabel={Schätzwert $\hat{\sym{thet}}_{\sym{idx_i}}$},
axis background/.style={fill=white},
xmajorgrids,
ymajorgrids,
grid style={dotted, opacity=0.25},
legend style={legend cell align=left, align=left}
]
\addplot[ybar, bar width=0.071, fill=white!10!black, draw=mycolor1, area legend] table[row sep=crcr] {%
1.00	24.05\\
2.00	12.01\\
3.00	12.10\\
4.00	-6.00\\
5.00	-6.02\\
6.00	-6.13\\
};
\addplot[forget plot, color=mycolor1] table[row sep=crcr] {%
0.54	0.00\\
6.46	0.00\\
};
\addlegendentry{Basis \ac{CCD}}

\addplot[ybar, bar width=0.071, fill=black!25!darkgray, draw=mycolor1, area legend] table[row sep=crcr] {%
1.00	23.97\\
2.00	12.03\\
3.00	12.00\\
4.00	-6.04\\
5.00	-6.05\\
6.00	-6.02\\
};
\addplot[forget plot, color=mycolor1] table[row sep=crcr] {%
0.54	0.00\\
6.46	0.00\\
};
\addlegendentry{Fall I ($-20\,\%$)}

\addplot[ybar, bar width=0.071, fill=black!45!gray, draw=mycolor1, area legend] table[row sep=crcr] {%
1.00	24.01\\
2.00	11.99\\
3.00	12.02\\
4.00	-6.06\\
5.00	-6.01\\
6.00	-5.97\\
};
\addplot[forget plot, color=mycolor1] table[row sep=crcr] {%
0.54	0.00\\
6.46	0.00\\
};
\addlegendentry{Fall I ($+20\,\%$)}

\addplot[ybar, bar width=0.071, fill=white!15!darkgray, draw=mycolor1, area legend] table[row sep=crcr] {%
1.00	23.99\\
2.00	12.01\\
3.00	12.05\\
4.00	-5.95\\
5.00	-5.88\\
6.00	-6.00\\
};
\addplot[forget plot, color=mycolor1] table[row sep=crcr] {%
0.54	0.00\\
6.46	0.00\\
};
\addlegendentry{Fall I (\ac{FC})}

\addplot[ybar, bar width=0.071, fill=white!45!black, draw=mycolor1, area legend] table[row sep=crcr] {%
1.00	24.03\\
2.00	12.02\\
3.00	11.97\\
4.00	-5.96\\
5.00	-6.02\\
6.00	-6.06\\
};
\addplot[forget plot, color=mycolor1] table[row sep=crcr] {%
0.54	0.00\\
6.46	0.00\\
};
\addlegendentry{Fall II (leicht)}

\addplot[ybar, bar width=0.071, fill=gray!85!lightgray, draw=mycolor1, area legend] table[row sep=crcr] {%
1.00	23.91\\
2.00	11.94\\
3.00	12.00\\
4.00	-6.03\\
5.00	-5.96\\
6.00	-6.02\\
};
\addplot[forget plot, color=mycolor1] table[row sep=crcr] {%
0.54	0.00\\
6.46	0.00\\
};
\addlegendentry{Fall II (stark)}

\addplot[ybar, bar width=0.071, fill=white!50!darkgray, draw=mycolor1, area legend] table[row sep=crcr] {%
1.00	24.07\\
2.00	11.92\\
3.00	12.03\\
4.00	-6.02\\
5.00	-6.01\\
6.00	-6.06\\
};
\addplot[forget plot, color=mycolor1] table[row sep=crcr] {%
0.54	0.00\\
6.46	0.00\\
};
\addlegendentry{Fall III}

\addplot[ybar, bar width=0.071, fill=black!5!lightgray, draw=mycolor1, area legend] table[row sep=crcr] {%
1.00	23.99\\
2.00	11.96\\
3.00	11.98\\
4.00	-6.02\\
5.00	-6.02\\
6.00	-5.96\\
};
\addplot[forget plot, color=mycolor1] table[row sep=crcr] {%
0.54	0.00\\
6.46	0.00\\
};
\addlegendentry{Fall IV: \ac{SP}}

\addplot[ybar, bar width=0.071, fill=white!80!black, draw=mycolor1, area legend] table[row sep=crcr] {%
1.00	24.04\\
2.00	11.92\\
3.00	12.05\\
4.00	-6.01\\
5.00	-5.91\\
6.00	-6.08\\
};
\addplot[forget plot, color=mycolor1] table[row sep=crcr] {%
0.54	0.00\\
6.46	0.00\\
};
\addlegendentry{Fall IV: \ac{ZP}}

\end{axis}
\end{tikzpicture}%
%     \caption{Absolute Schätzwerte der Modellkoeffizienten inklusive der $95\,\%$-Konfidenzintervalle für $\sym{snr}=1$.}
%     \label{fig:app.c_betadev}
% \end{figure}


\begin{figure}[htbp]
    \centering
    \setlength\figurewidth{0.9\textwidth}
    \setlength\figureheight{16cm}
    % This file was created by matlab2tikz.
%
\definecolor{mycolor1}{rgb}{0.12941,0.12941,0.12941}%
%
\begin{tikzpicture}

\begin{axis}[%
width=0.391\figurewidth,
height=0.265\figureheight,
at={(0\figurewidth,0.735\figureheight)},
scale only axis,
xmin=1.00,
xmax=10000.00,
xtick={2000.00,4000.00,6000.00,8000.00,10000.00},
xticklabels={\empty},
ymin=0.00,
ymax=105.00,
ylabel style={font=\color{mycolor1}},
ylabel={\sym{power} [\%]},
axis background/.style={fill=white},
title style={font=\bfseries\color{mycolor1}},
title={\textbf{$\sym{beta}_0$}},
xmajorgrids,
ymajorgrids,
grid style={dotted, opacity=0.25},
legend style={at={(0.03,0.03)}, anchor=south west, legend cell align=left, align=left}
]
\addplot [color=white!10!black, line width=1.2pt]
  table[row sep=crcr]{%
1.00	100.00\\
2.00	100.00\\
3.00	100.00\\
4.00	100.00\\
5.00	100.00\\
6.00	100.00\\
7.00	100.00\\
8.00	100.00\\
9.00	100.00\\
10.00	100.00\\
11.00	100.00\\
12.00	100.00\\
13.00	100.00\\
14.00	100.00\\
15.00	100.00\\
16.00	100.00\\
17.00	100.00\\
18.00	100.00\\
19.00	100.00\\
21.00	100.00\\
22.00	100.00\\
23.00	100.00\\
25.00	100.00\\
26.00	100.00\\
28.00	100.00\\
30.00	100.00\\
32.00	100.00\\
34.00	100.00\\
36.00	100.00\\
38.00	100.00\\
41.00	100.00\\
43.00	100.00\\
46.00	100.00\\
49.00	100.00\\
52.00	100.00\\
56.00	100.00\\
59.00	100.00\\
63.00	100.00\\
67.00	100.00\\
71.00	100.00\\
76.00	100.00\\
81.00	100.00\\
86.00	100.00\\
91.00	100.00\\
97.00	100.00\\
103.00	100.00\\
110.00	100.00\\
117.00	100.00\\
124.00	100.00\\
132.00	100.00\\
140.00	100.00\\
149.00	100.00\\
159.00	100.00\\
169.00	100.00\\
180.00	100.00\\
191.00	100.00\\
204.00	100.00\\
217.00	100.00\\
230.00	100.00\\
245.00	100.00\\
261.00	100.00\\
277.00	100.00\\
295.00	100.00\\
314.00	100.00\\
334.00	100.00\\
355.00	100.00\\
378.00	100.00\\
402.00	100.00\\
427.00	100.00\\
455.00	100.00\\
484.00	100.00\\
515.00	100.00\\
547.00	100.00\\
582.00	100.00\\
619.00	100.00\\
659.00	100.00\\
701.00	100.00\\
746.00	100.00\\
793.00	100.00\\
844.00	100.00\\
897.00	100.00\\
955.00	100.00\\
1016.00	100.00\\
1080.00	100.00\\
1149.00	100.00\\
1222.00	100.00\\
1300.00	100.00\\
1383.00	100.00\\
1472.00	100.00\\
1565.00	100.00\\
1665.00	100.00\\
1771.00	100.00\\
1884.00	100.00\\
2005.00	100.00\\
2132.00	100.00\\
2268.00	100.00\\
2413.00	100.00\\
2567.00	100.00\\
2730.00	100.00\\
2905.00	100.00\\
3090.00	100.00\\
3287.00	100.00\\
3496.00	100.00\\
3719.00	100.00\\
3957.00	100.00\\
4209.00	100.00\\
4477.00	100.00\\
4763.00	100.00\\
5066.00	100.00\\
5389.00	100.00\\
5733.00	100.00\\
6099.00	100.00\\
6488.00	100.00\\
6901.00	100.00\\
7341.00	100.00\\
7809.00	100.00\\
8307.00	100.00\\
8837.00	100.00\\
9401.00	100.00\\
10000.00	100.00\\
};
\addlegendentry{Basis \ac{CCD}}

\addplot [color=black!25!darkgray, dashed, line width=1.2pt]
  table[row sep=crcr]{%
1.00	100.00\\
2.00	100.00\\
3.00	100.00\\
4.00	100.00\\
5.00	100.00\\
6.00	100.00\\
7.00	100.00\\
8.00	100.00\\
9.00	100.00\\
10.00	100.00\\
11.00	100.00\\
12.00	100.00\\
13.00	100.00\\
14.00	100.00\\
15.00	100.00\\
16.00	100.00\\
17.00	100.00\\
18.00	100.00\\
19.00	100.00\\
21.00	100.00\\
22.00	100.00\\
23.00	100.00\\
25.00	100.00\\
26.00	100.00\\
28.00	100.00\\
30.00	100.00\\
32.00	100.00\\
34.00	100.00\\
36.00	100.00\\
38.00	100.00\\
41.00	100.00\\
43.00	100.00\\
46.00	100.00\\
49.00	100.00\\
52.00	100.00\\
56.00	100.00\\
59.00	100.00\\
63.00	100.00\\
67.00	100.00\\
71.00	100.00\\
76.00	100.00\\
81.00	100.00\\
86.00	100.00\\
91.00	100.00\\
97.00	100.00\\
103.00	100.00\\
110.00	100.00\\
117.00	100.00\\
124.00	100.00\\
132.00	100.00\\
140.00	100.00\\
149.00	100.00\\
159.00	100.00\\
169.00	100.00\\
180.00	100.00\\
191.00	100.00\\
204.00	100.00\\
217.00	100.00\\
230.00	100.00\\
245.00	100.00\\
261.00	100.00\\
277.00	100.00\\
295.00	100.00\\
314.00	100.00\\
334.00	100.00\\
355.00	100.00\\
378.00	100.00\\
402.00	100.00\\
427.00	100.00\\
455.00	100.00\\
484.00	100.00\\
515.00	100.00\\
547.00	100.00\\
582.00	100.00\\
619.00	100.00\\
659.00	100.00\\
701.00	100.00\\
746.00	100.00\\
793.00	100.00\\
844.00	100.00\\
897.00	100.00\\
955.00	100.00\\
1016.00	100.00\\
1080.00	100.00\\
1149.00	100.00\\
1222.00	100.00\\
1300.00	100.00\\
1383.00	100.00\\
1472.00	100.00\\
1565.00	100.00\\
1665.00	100.00\\
1771.00	100.00\\
1884.00	100.00\\
2005.00	100.00\\
2132.00	100.00\\
2268.00	100.00\\
2413.00	100.00\\
2567.00	100.00\\
2730.00	100.00\\
2905.00	100.00\\
3090.00	100.00\\
3287.00	100.00\\
3496.00	100.00\\
3719.00	100.00\\
3957.00	100.00\\
4209.00	100.00\\
4477.00	100.00\\
4763.00	100.00\\
5066.00	100.00\\
5389.00	100.00\\
5733.00	100.00\\
6099.00	100.00\\
6488.00	100.00\\
6901.00	100.00\\
7341.00	100.00\\
7809.00	100.00\\
8307.00	100.00\\
8837.00	100.00\\
9401.00	100.00\\
10000.00	100.00\\
};
\addlegendentry{Fall I ($-20\,\%$)}

\addplot [color=black!45!gray, dotted, line width=1.2pt]
  table[row sep=crcr]{%
1.00	100.00\\
2.00	100.00\\
3.00	100.00\\
4.00	100.00\\
5.00	100.00\\
6.00	100.00\\
7.00	100.00\\
8.00	100.00\\
9.00	100.00\\
10.00	100.00\\
11.00	100.00\\
12.00	100.00\\
13.00	100.00\\
14.00	100.00\\
15.00	100.00\\
16.00	100.00\\
17.00	100.00\\
18.00	100.00\\
19.00	100.00\\
21.00	100.00\\
22.00	100.00\\
23.00	100.00\\
25.00	100.00\\
26.00	100.00\\
28.00	100.00\\
30.00	100.00\\
32.00	100.00\\
34.00	100.00\\
36.00	100.00\\
38.00	100.00\\
41.00	100.00\\
43.00	100.00\\
46.00	100.00\\
49.00	100.00\\
52.00	100.00\\
56.00	100.00\\
59.00	100.00\\
63.00	100.00\\
67.00	100.00\\
71.00	100.00\\
76.00	100.00\\
81.00	100.00\\
86.00	100.00\\
91.00	100.00\\
97.00	100.00\\
103.00	100.00\\
110.00	100.00\\
117.00	100.00\\
124.00	100.00\\
132.00	100.00\\
140.00	100.00\\
149.00	100.00\\
159.00	100.00\\
169.00	100.00\\
180.00	100.00\\
191.00	100.00\\
204.00	100.00\\
217.00	100.00\\
230.00	100.00\\
245.00	100.00\\
261.00	100.00\\
277.00	100.00\\
295.00	100.00\\
314.00	100.00\\
334.00	100.00\\
355.00	100.00\\
378.00	100.00\\
402.00	100.00\\
427.00	100.00\\
455.00	100.00\\
484.00	100.00\\
515.00	100.00\\
547.00	100.00\\
582.00	100.00\\
619.00	100.00\\
659.00	100.00\\
701.00	100.00\\
746.00	100.00\\
793.00	100.00\\
844.00	100.00\\
897.00	100.00\\
955.00	100.00\\
1016.00	100.00\\
1080.00	100.00\\
1149.00	100.00\\
1222.00	100.00\\
1300.00	100.00\\
1383.00	100.00\\
1472.00	100.00\\
1565.00	100.00\\
1665.00	100.00\\
1771.00	100.00\\
1884.00	100.00\\
2005.00	100.00\\
2132.00	100.00\\
2268.00	100.00\\
2413.00	100.00\\
2567.00	100.00\\
2730.00	100.00\\
2905.00	100.00\\
3090.00	100.00\\
3287.00	100.00\\
3496.00	100.00\\
3719.00	100.00\\
3957.00	100.00\\
4209.00	100.00\\
4477.00	100.00\\
4763.00	100.00\\
5066.00	100.00\\
5389.00	100.00\\
5733.00	100.00\\
6099.00	100.00\\
6488.00	100.00\\
6901.00	100.00\\
7341.00	99.99\\
7809.00	99.99\\
8307.00	99.99\\
8837.00	99.99\\
9401.00	99.99\\
10000.00	99.99\\
};
\addlegendentry{Fall I ($+20\,\%$)}

\addplot [color=white!15!darkgray, dashdotted, line width=1.2pt]
  table[row sep=crcr]{%
1.00	100.00\\
2.00	100.00\\
3.00	100.00\\
4.00	100.00\\
5.00	100.00\\
6.00	100.00\\
7.00	100.00\\
8.00	100.00\\
9.00	100.00\\
10.00	100.00\\
11.00	100.00\\
12.00	100.00\\
13.00	100.00\\
14.00	100.00\\
15.00	100.00\\
16.00	100.00\\
17.00	100.00\\
18.00	100.00\\
19.00	100.00\\
21.00	100.00\\
22.00	100.00\\
23.00	100.00\\
25.00	100.00\\
26.00	100.00\\
28.00	100.00\\
30.00	100.00\\
32.00	100.00\\
34.00	100.00\\
36.00	100.00\\
38.00	100.00\\
41.00	100.00\\
43.00	100.00\\
46.00	100.00\\
49.00	100.00\\
52.00	100.00\\
56.00	100.00\\
59.00	100.00\\
63.00	100.00\\
67.00	100.00\\
71.00	100.00\\
76.00	100.00\\
81.00	100.00\\
86.00	100.00\\
91.00	100.00\\
97.00	100.00\\
103.00	100.00\\
110.00	100.00\\
117.00	100.00\\
124.00	100.00\\
132.00	100.00\\
140.00	100.00\\
149.00	100.00\\
159.00	100.00\\
169.00	100.00\\
180.00	100.00\\
191.00	100.00\\
204.00	100.00\\
217.00	100.00\\
230.00	100.00\\
245.00	100.00\\
261.00	100.00\\
277.00	100.00\\
295.00	100.00\\
314.00	100.00\\
334.00	100.00\\
355.00	100.00\\
378.00	100.00\\
402.00	100.00\\
427.00	100.00\\
455.00	100.00\\
484.00	100.00\\
515.00	100.00\\
547.00	100.00\\
582.00	100.00\\
619.00	100.00\\
659.00	100.00\\
701.00	100.00\\
746.00	100.00\\
793.00	100.00\\
844.00	100.00\\
897.00	100.00\\
955.00	100.00\\
1016.00	100.00\\
1080.00	100.00\\
1149.00	100.00\\
1222.00	100.00\\
1300.00	100.00\\
1383.00	100.00\\
1472.00	100.00\\
1565.00	100.00\\
1665.00	100.00\\
1771.00	100.00\\
1884.00	100.00\\
2005.00	100.00\\
2132.00	100.00\\
2268.00	100.00\\
2413.00	100.00\\
2567.00	100.00\\
2730.00	100.00\\
2905.00	100.00\\
3090.00	100.00\\
3287.00	100.00\\
3496.00	100.00\\
3719.00	100.00\\
3957.00	100.00\\
4209.00	100.00\\
4477.00	100.00\\
4763.00	100.00\\
5066.00	100.00\\
5389.00	100.00\\
5733.00	100.00\\
6099.00	100.00\\
6488.00	100.00\\
6901.00	100.00\\
7341.00	100.00\\
7809.00	100.00\\
8307.00	100.00\\
8837.00	100.00\\
9401.00	100.00\\
10000.00	100.00\\
};
\addlegendentry{Fall I (\ac{FC})}

\addplot [color=white!45!black, line width=1.2pt]
  table[row sep=crcr]{%
1.00	100.00\\
2.00	100.00\\
3.00	100.00\\
4.00	100.00\\
5.00	100.00\\
6.00	100.00\\
7.00	100.00\\
8.00	100.00\\
9.00	100.00\\
10.00	100.00\\
11.00	100.00\\
12.00	100.00\\
13.00	100.00\\
14.00	100.00\\
15.00	100.00\\
16.00	100.00\\
17.00	100.00\\
18.00	100.00\\
19.00	100.00\\
21.00	100.00\\
22.00	100.00\\
23.00	100.00\\
25.00	100.00\\
26.00	100.00\\
28.00	100.00\\
30.00	100.00\\
32.00	100.00\\
34.00	100.00\\
36.00	100.00\\
38.00	100.00\\
41.00	100.00\\
43.00	100.00\\
46.00	100.00\\
49.00	100.00\\
52.00	100.00\\
56.00	100.00\\
59.00	100.00\\
63.00	100.00\\
67.00	100.00\\
71.00	100.00\\
76.00	100.00\\
81.00	100.00\\
86.00	100.00\\
91.00	100.00\\
97.00	100.00\\
103.00	100.00\\
110.00	100.00\\
117.00	100.00\\
124.00	100.00\\
132.00	100.00\\
140.00	100.00\\
149.00	100.00\\
159.00	100.00\\
169.00	100.00\\
180.00	100.00\\
191.00	100.00\\
204.00	100.00\\
217.00	100.00\\
230.00	100.00\\
245.00	100.00\\
261.00	100.00\\
277.00	100.00\\
295.00	100.00\\
314.00	100.00\\
334.00	100.00\\
355.00	100.00\\
378.00	100.00\\
402.00	100.00\\
427.00	100.00\\
455.00	100.00\\
484.00	100.00\\
515.00	100.00\\
547.00	100.00\\
582.00	100.00\\
619.00	100.00\\
659.00	100.00\\
701.00	100.00\\
746.00	100.00\\
793.00	100.00\\
844.00	100.00\\
897.00	100.00\\
955.00	100.00\\
1016.00	100.00\\
1080.00	100.00\\
1149.00	100.00\\
1222.00	100.00\\
1300.00	100.00\\
1383.00	100.00\\
1472.00	100.00\\
1565.00	100.00\\
1665.00	100.00\\
1771.00	100.00\\
1884.00	100.00\\
2005.00	100.00\\
2132.00	100.00\\
2268.00	100.00\\
2413.00	100.00\\
2567.00	100.00\\
2730.00	100.00\\
2905.00	100.00\\
3090.00	100.00\\
3287.00	100.00\\
3496.00	100.00\\
3719.00	100.00\\
3957.00	100.00\\
4209.00	100.00\\
4477.00	100.00\\
4763.00	100.00\\
5066.00	100.00\\
5389.00	100.00\\
5733.00	100.00\\
6099.00	100.00\\
6488.00	100.00\\
6901.00	100.00\\
7341.00	100.00\\
7809.00	99.99\\
8307.00	99.99\\
8837.00	99.99\\
9401.00	99.99\\
10000.00	99.99\\
};
\addlegendentry{Fall II (leicht)}

\addplot [color=gray!85!lightgray, dashed, line width=1.2pt]
  table[row sep=crcr]{%
1.00	100.00\\
2.00	100.00\\
3.00	100.00\\
4.00	100.00\\
5.00	100.00\\
6.00	100.00\\
7.00	100.00\\
8.00	100.00\\
9.00	100.00\\
10.00	100.00\\
11.00	100.00\\
12.00	100.00\\
13.00	100.00\\
14.00	100.00\\
15.00	100.00\\
16.00	100.00\\
17.00	100.00\\
18.00	100.00\\
19.00	100.00\\
21.00	100.00\\
22.00	100.00\\
23.00	100.00\\
25.00	100.00\\
26.00	100.00\\
28.00	100.00\\
30.00	100.00\\
32.00	100.00\\
34.00	100.00\\
36.00	100.00\\
38.00	100.00\\
41.00	100.00\\
43.00	100.00\\
46.00	100.00\\
49.00	100.00\\
52.00	100.00\\
56.00	100.00\\
59.00	100.00\\
63.00	100.00\\
67.00	100.00\\
71.00	100.00\\
76.00	100.00\\
81.00	100.00\\
86.00	100.00\\
91.00	100.00\\
97.00	100.00\\
103.00	100.00\\
110.00	100.00\\
117.00	100.00\\
124.00	100.00\\
132.00	100.00\\
140.00	100.00\\
149.00	100.00\\
159.00	100.00\\
169.00	100.00\\
180.00	100.00\\
191.00	100.00\\
204.00	100.00\\
217.00	100.00\\
230.00	100.00\\
245.00	100.00\\
261.00	100.00\\
277.00	100.00\\
295.00	100.00\\
314.00	100.00\\
334.00	100.00\\
355.00	100.00\\
378.00	100.00\\
402.00	100.00\\
427.00	100.00\\
455.00	100.00\\
484.00	100.00\\
515.00	100.00\\
547.00	100.00\\
582.00	100.00\\
619.00	100.00\\
659.00	100.00\\
701.00	100.00\\
746.00	100.00\\
793.00	100.00\\
844.00	100.00\\
897.00	100.00\\
955.00	100.00\\
1016.00	100.00\\
1080.00	100.00\\
1149.00	100.00\\
1222.00	100.00\\
1300.00	100.00\\
1383.00	100.00\\
1472.00	100.00\\
1565.00	100.00\\
1665.00	100.00\\
1771.00	100.00\\
1884.00	100.00\\
2005.00	100.00\\
2132.00	100.00\\
2268.00	100.00\\
2413.00	100.00\\
2567.00	100.00\\
2730.00	100.00\\
2905.00	100.00\\
3090.00	100.00\\
3287.00	100.00\\
3496.00	100.00\\
3719.00	100.00\\
3957.00	100.00\\
4209.00	100.00\\
4477.00	100.00\\
4763.00	100.00\\
5066.00	100.00\\
5389.00	99.98\\
5733.00	99.98\\
6099.00	99.98\\
6488.00	99.98\\
6901.00	99.99\\
7341.00	99.99\\
7809.00	99.99\\
8307.00	99.99\\
8837.00	99.99\\
9401.00	99.99\\
10000.00	99.99\\
};
\addlegendentry{Fall II (stark)}

\addplot [color=white!50!darkgray, dotted, line width=1.2pt]
  table[row sep=crcr]{%
1.00	100.00\\
2.00	100.00\\
3.00	100.00\\
4.00	100.00\\
5.00	100.00\\
6.00	100.00\\
7.00	100.00\\
8.00	100.00\\
9.00	100.00\\
10.00	100.00\\
11.00	100.00\\
12.00	100.00\\
13.00	100.00\\
14.00	100.00\\
15.00	100.00\\
16.00	100.00\\
17.00	100.00\\
18.00	100.00\\
19.00	100.00\\
21.00	100.00\\
22.00	100.00\\
23.00	100.00\\
25.00	100.00\\
26.00	100.00\\
28.00	100.00\\
30.00	100.00\\
32.00	100.00\\
34.00	100.00\\
36.00	100.00\\
38.00	100.00\\
41.00	100.00\\
43.00	100.00\\
46.00	100.00\\
49.00	100.00\\
52.00	100.00\\
56.00	100.00\\
59.00	100.00\\
63.00	100.00\\
67.00	100.00\\
71.00	100.00\\
76.00	100.00\\
81.00	100.00\\
86.00	100.00\\
91.00	100.00\\
97.00	100.00\\
103.00	100.00\\
110.00	100.00\\
117.00	100.00\\
124.00	100.00\\
132.00	100.00\\
140.00	100.00\\
149.00	100.00\\
159.00	100.00\\
169.00	100.00\\
180.00	100.00\\
191.00	100.00\\
204.00	100.00\\
217.00	100.00\\
230.00	100.00\\
245.00	100.00\\
261.00	100.00\\
277.00	100.00\\
295.00	100.00\\
314.00	100.00\\
334.00	100.00\\
355.00	100.00\\
378.00	100.00\\
402.00	100.00\\
427.00	100.00\\
455.00	100.00\\
484.00	100.00\\
515.00	100.00\\
547.00	100.00\\
582.00	100.00\\
619.00	100.00\\
659.00	100.00\\
701.00	100.00\\
746.00	100.00\\
793.00	100.00\\
844.00	100.00\\
897.00	100.00\\
955.00	100.00\\
1016.00	100.00\\
1080.00	100.00\\
1149.00	100.00\\
1222.00	100.00\\
1300.00	100.00\\
1383.00	100.00\\
1472.00	100.00\\
1565.00	100.00\\
1665.00	100.00\\
1771.00	100.00\\
1884.00	100.00\\
2005.00	100.00\\
2132.00	100.00\\
2268.00	100.00\\
2413.00	100.00\\
2567.00	100.00\\
2730.00	100.00\\
2905.00	100.00\\
3090.00	100.00\\
3287.00	100.00\\
3496.00	100.00\\
3719.00	100.00\\
3957.00	100.00\\
4209.00	100.00\\
4477.00	100.00\\
4763.00	100.00\\
5066.00	100.00\\
5389.00	100.00\\
5733.00	99.98\\
6099.00	99.98\\
6488.00	99.98\\
6901.00	99.99\\
7341.00	99.99\\
7809.00	99.99\\
8307.00	99.99\\
8837.00	99.99\\
9401.00	99.99\\
10000.00	99.99\\
};
\addlegendentry{Fall III}

\addplot [color=black!5!lightgray, dashdotted, line width=1.2pt]
  table[row sep=crcr]{%
1.00	100.00\\
2.00	100.00\\
3.00	100.00\\
4.00	100.00\\
5.00	100.00\\
6.00	100.00\\
7.00	100.00\\
8.00	100.00\\
9.00	100.00\\
10.00	100.00\\
11.00	100.00\\
12.00	100.00\\
13.00	100.00\\
14.00	100.00\\
15.00	100.00\\
16.00	100.00\\
17.00	100.00\\
18.00	100.00\\
19.00	100.00\\
21.00	100.00\\
22.00	100.00\\
23.00	100.00\\
25.00	100.00\\
26.00	100.00\\
28.00	100.00\\
30.00	100.00\\
32.00	100.00\\
34.00	100.00\\
36.00	100.00\\
38.00	100.00\\
41.00	100.00\\
43.00	100.00\\
46.00	100.00\\
49.00	100.00\\
52.00	100.00\\
56.00	100.00\\
59.00	100.00\\
63.00	100.00\\
67.00	100.00\\
71.00	100.00\\
76.00	100.00\\
81.00	100.00\\
86.00	100.00\\
91.00	100.00\\
97.00	100.00\\
103.00	100.00\\
110.00	100.00\\
117.00	100.00\\
124.00	100.00\\
132.00	100.00\\
140.00	100.00\\
149.00	100.00\\
159.00	100.00\\
169.00	100.00\\
180.00	100.00\\
191.00	100.00\\
204.00	100.00\\
217.00	100.00\\
230.00	100.00\\
245.00	100.00\\
261.00	100.00\\
277.00	100.00\\
295.00	100.00\\
314.00	100.00\\
334.00	100.00\\
355.00	100.00\\
378.00	100.00\\
402.00	100.00\\
427.00	100.00\\
455.00	100.00\\
484.00	100.00\\
515.00	100.00\\
547.00	100.00\\
582.00	100.00\\
619.00	100.00\\
659.00	100.00\\
701.00	100.00\\
746.00	100.00\\
793.00	100.00\\
844.00	100.00\\
897.00	100.00\\
955.00	100.00\\
1016.00	100.00\\
1080.00	100.00\\
1149.00	100.00\\
1222.00	100.00\\
1300.00	100.00\\
1383.00	100.00\\
1472.00	100.00\\
1565.00	100.00\\
1665.00	100.00\\
1771.00	100.00\\
1884.00	100.00\\
2005.00	100.00\\
2132.00	100.00\\
2268.00	100.00\\
2413.00	100.00\\
2567.00	100.00\\
2730.00	100.00\\
2905.00	100.00\\
3090.00	100.00\\
3287.00	99.97\\
3496.00	99.97\\
3719.00	99.97\\
3957.00	99.97\\
4209.00	99.98\\
4477.00	99.98\\
4763.00	99.98\\
5066.00	99.98\\
5389.00	99.98\\
5733.00	99.98\\
6099.00	99.98\\
6488.00	99.98\\
6901.00	99.99\\
7341.00	99.99\\
7809.00	99.99\\
8307.00	99.99\\
8837.00	99.99\\
9401.00	99.99\\
10000.00	99.99\\
};
\addlegendentry{Fall IV: \ac{SP}}

\addplot [color=white!80!black, line width=1.2pt]
  table[row sep=crcr]{%
1.00	100.00\\
2.00	100.00\\
3.00	100.00\\
4.00	100.00\\
5.00	100.00\\
6.00	100.00\\
7.00	100.00\\
8.00	100.00\\
9.00	100.00\\
10.00	100.00\\
11.00	100.00\\
12.00	100.00\\
13.00	100.00\\
14.00	100.00\\
15.00	100.00\\
16.00	100.00\\
17.00	100.00\\
18.00	100.00\\
19.00	100.00\\
21.00	100.00\\
22.00	100.00\\
23.00	100.00\\
25.00	100.00\\
26.00	100.00\\
28.00	100.00\\
30.00	100.00\\
32.00	100.00\\
34.00	100.00\\
36.00	100.00\\
38.00	100.00\\
41.00	100.00\\
43.00	100.00\\
46.00	100.00\\
49.00	100.00\\
52.00	100.00\\
56.00	100.00\\
59.00	100.00\\
63.00	100.00\\
67.00	100.00\\
71.00	100.00\\
76.00	100.00\\
81.00	100.00\\
86.00	100.00\\
91.00	100.00\\
97.00	100.00\\
103.00	100.00\\
110.00	100.00\\
117.00	100.00\\
124.00	100.00\\
132.00	100.00\\
140.00	100.00\\
149.00	100.00\\
159.00	100.00\\
169.00	100.00\\
180.00	100.00\\
191.00	100.00\\
204.00	100.00\\
217.00	100.00\\
230.00	100.00\\
245.00	100.00\\
261.00	100.00\\
277.00	100.00\\
295.00	100.00\\
314.00	100.00\\
334.00	100.00\\
355.00	100.00\\
378.00	100.00\\
402.00	100.00\\
427.00	100.00\\
455.00	100.00\\
484.00	100.00\\
515.00	100.00\\
547.00	100.00\\
582.00	100.00\\
619.00	100.00\\
659.00	100.00\\
701.00	100.00\\
746.00	100.00\\
793.00	100.00\\
844.00	100.00\\
897.00	100.00\\
955.00	100.00\\
1016.00	100.00\\
1080.00	100.00\\
1149.00	100.00\\
1222.00	100.00\\
1300.00	100.00\\
1383.00	100.00\\
1472.00	100.00\\
1565.00	100.00\\
1665.00	100.00\\
1771.00	100.00\\
1884.00	100.00\\
2005.00	100.00\\
2132.00	100.00\\
2268.00	100.00\\
2413.00	100.00\\
2567.00	100.00\\
2730.00	99.96\\
2905.00	99.97\\
3090.00	99.97\\
3287.00	99.97\\
3496.00	99.94\\
3719.00	99.95\\
3957.00	99.95\\
4209.00	99.95\\
4477.00	99.96\\
4763.00	99.96\\
5066.00	99.96\\
5389.00	99.96\\
5733.00	99.97\\
6099.00	99.97\\
6488.00	99.97\\
6901.00	99.97\\
7341.00	99.97\\
7809.00	99.97\\
8307.00	99.96\\
8837.00	99.97\\
9401.00	99.97\\
10000.00	99.97\\
};
\addlegendentry{Fall IV: \ac{ZP}}

\end{axis}

\begin{axis}[%
width=0.391\figurewidth,
height=0.265\figureheight,
at={(0.54\figurewidth,0.735\figureheight)},
scale only axis,
xmin=1.00,
xmax=10000.00,
xtick={2000.00,4000.00,6000.00,8000.00,10000.00},
xticklabels={\empty},
ymin=0.00,
ymax=105.00,
ytick={0.00,20.00,40.00,60.00,80.00,100.00},
yticklabels={\empty},
axis background/.style={fill=white},
title style={font=\bfseries\color{mycolor1}},
title={\textbf{$\sym{beta}_1$}},
xmajorgrids,
ymajorgrids,
grid style={dotted, opacity=0.25}
]
\addplot [color=white!10!black, line width=1.2pt, forget plot]
  table[row sep=crcr]{%
1.00	100.00\\
2.00	100.00\\
3.00	100.00\\
4.00	100.00\\
5.00	100.00\\
6.00	100.00\\
7.00	100.00\\
8.00	100.00\\
9.00	100.00\\
10.00	100.00\\
11.00	100.00\\
12.00	100.00\\
13.00	100.00\\
14.00	100.00\\
15.00	100.00\\
16.00	100.00\\
17.00	100.00\\
18.00	100.00\\
19.00	100.00\\
21.00	95.24\\
22.00	95.45\\
23.00	95.65\\
25.00	96.00\\
26.00	96.15\\
28.00	96.43\\
30.00	96.67\\
32.00	96.88\\
34.00	94.12\\
36.00	91.67\\
38.00	92.11\\
41.00	92.68\\
43.00	93.02\\
46.00	93.48\\
49.00	93.88\\
52.00	94.23\\
56.00	92.86\\
59.00	93.22\\
63.00	93.65\\
67.00	92.54\\
71.00	92.96\\
76.00	93.42\\
81.00	93.83\\
86.00	94.19\\
91.00	94.51\\
97.00	94.85\\
103.00	95.15\\
110.00	95.45\\
117.00	94.87\\
124.00	95.16\\
132.00	94.70\\
140.00	95.00\\
149.00	95.30\\
159.00	95.60\\
169.00	95.27\\
180.00	95.00\\
191.00	95.29\\
204.00	95.10\\
217.00	94.93\\
230.00	94.35\\
245.00	94.29\\
261.00	94.25\\
277.00	94.58\\
295.00	94.92\\
314.00	94.59\\
334.00	94.61\\
355.00	94.37\\
378.00	93.92\\
402.00	94.28\\
427.00	93.91\\
455.00	93.63\\
484.00	93.60\\
515.00	93.40\\
547.00	93.60\\
582.00	93.64\\
619.00	93.70\\
659.00	93.78\\
701.00	94.01\\
746.00	93.57\\
793.00	93.57\\
844.00	93.72\\
897.00	93.42\\
955.00	93.61\\
1016.00	93.90\\
1080.00	93.89\\
1149.00	93.91\\
1222.00	93.62\\
1300.00	93.77\\
1383.00	93.85\\
1472.00	93.68\\
1565.00	93.87\\
1665.00	93.75\\
1771.00	93.90\\
1884.00	93.79\\
2005.00	93.62\\
2132.00	93.62\\
2268.00	93.47\\
2413.00	93.54\\
2567.00	93.69\\
2730.00	93.77\\
2905.00	93.67\\
3090.00	93.43\\
3287.00	93.43\\
3496.00	93.62\\
3719.00	93.68\\
3957.00	93.56\\
4209.00	93.47\\
4477.00	93.50\\
4763.00	93.58\\
5066.00	93.51\\
5389.00	93.52\\
5733.00	93.49\\
6099.00	93.54\\
6488.00	93.48\\
6901.00	93.44\\
7341.00	93.49\\
7809.00	93.47\\
8307.00	93.46\\
8837.00	93.57\\
9401.00	93.50\\
10000.00	93.58\\
};
\addplot [color=black!25!darkgray, dashed, line width=1.2pt, forget plot]
  table[row sep=crcr]{%
1.00	100.00\\
2.00	100.00\\
3.00	100.00\\
4.00	100.00\\
5.00	100.00\\
6.00	100.00\\
7.00	100.00\\
8.00	100.00\\
9.00	100.00\\
10.00	100.00\\
11.00	100.00\\
12.00	91.67\\
13.00	92.31\\
14.00	92.86\\
15.00	93.33\\
16.00	93.75\\
17.00	94.12\\
18.00	94.44\\
19.00	94.74\\
21.00	90.48\\
22.00	90.91\\
23.00	91.30\\
25.00	92.00\\
26.00	88.46\\
28.00	85.71\\
30.00	83.33\\
32.00	84.38\\
34.00	85.29\\
36.00	83.33\\
38.00	84.21\\
41.00	82.93\\
43.00	83.72\\
46.00	84.78\\
49.00	85.71\\
52.00	86.54\\
56.00	85.71\\
59.00	84.75\\
63.00	85.71\\
67.00	86.57\\
71.00	87.32\\
76.00	86.84\\
81.00	87.65\\
86.00	88.37\\
91.00	89.01\\
97.00	89.69\\
103.00	90.29\\
110.00	90.91\\
117.00	91.45\\
124.00	91.94\\
132.00	92.42\\
140.00	91.43\\
149.00	91.95\\
159.00	92.45\\
169.00	92.31\\
180.00	92.78\\
191.00	92.67\\
204.00	92.65\\
217.00	92.63\\
230.00	93.04\\
245.00	93.47\\
261.00	93.10\\
277.00	93.50\\
295.00	93.22\\
314.00	93.63\\
334.00	93.41\\
355.00	92.96\\
378.00	93.12\\
402.00	93.53\\
427.00	93.44\\
455.00	93.41\\
484.00	93.39\\
515.00	93.20\\
547.00	93.24\\
582.00	93.47\\
619.00	93.21\\
659.00	93.47\\
701.00	93.44\\
746.00	93.70\\
793.00	93.69\\
844.00	93.72\\
897.00	93.87\\
955.00	94.03\\
1016.00	93.70\\
1080.00	93.70\\
1149.00	93.39\\
1222.00	93.37\\
1300.00	93.31\\
1383.00	93.35\\
1472.00	93.34\\
1565.00	93.35\\
1665.00	93.45\\
1771.00	93.34\\
1884.00	93.58\\
2005.00	93.82\\
2132.00	93.71\\
2268.00	93.69\\
2413.00	93.66\\
2567.00	93.42\\
2730.00	93.48\\
2905.00	93.49\\
3090.00	93.40\\
3287.00	93.46\\
3496.00	93.51\\
3719.00	93.49\\
3957.00	93.33\\
4209.00	93.40\\
4477.00	93.50\\
4763.00	93.51\\
5066.00	93.33\\
5389.00	93.36\\
5733.00	93.34\\
6099.00	93.39\\
6488.00	93.33\\
6901.00	93.31\\
7341.00	93.27\\
7809.00	93.37\\
8307.00	93.34\\
8837.00	93.19\\
9401.00	93.15\\
10000.00	93.23\\
};
\addplot [color=black!45!gray, dotted, line width=1.2pt, forget plot]
  table[row sep=crcr]{%
1.00	100.00\\
2.00	100.00\\
3.00	100.00\\
4.00	100.00\\
5.00	100.00\\
6.00	100.00\\
7.00	100.00\\
8.00	100.00\\
9.00	100.00\\
10.00	100.00\\
11.00	100.00\\
12.00	100.00\\
13.00	100.00\\
14.00	100.00\\
15.00	100.00\\
16.00	93.75\\
17.00	94.12\\
18.00	94.44\\
19.00	94.74\\
21.00	95.24\\
22.00	95.45\\
23.00	95.65\\
25.00	96.00\\
26.00	96.15\\
28.00	96.43\\
30.00	96.67\\
32.00	96.88\\
34.00	97.06\\
36.00	97.22\\
38.00	97.37\\
41.00	97.56\\
43.00	97.67\\
46.00	97.83\\
49.00	97.96\\
52.00	98.08\\
56.00	98.21\\
59.00	98.31\\
63.00	98.41\\
67.00	98.51\\
71.00	98.59\\
76.00	97.37\\
81.00	97.53\\
86.00	96.51\\
91.00	96.70\\
97.00	95.88\\
103.00	96.12\\
110.00	95.45\\
117.00	95.73\\
124.00	95.97\\
132.00	95.45\\
140.00	95.00\\
149.00	95.30\\
159.00	95.60\\
169.00	95.27\\
180.00	95.00\\
191.00	93.72\\
204.00	94.12\\
217.00	94.47\\
230.00	94.78\\
245.00	95.10\\
261.00	95.02\\
277.00	95.31\\
295.00	95.59\\
314.00	95.86\\
334.00	96.11\\
355.00	96.06\\
378.00	95.77\\
402.00	95.77\\
427.00	95.55\\
455.00	95.38\\
484.00	94.83\\
515.00	94.76\\
547.00	94.70\\
582.00	94.67\\
619.00	94.51\\
659.00	94.54\\
701.00	94.29\\
746.00	94.10\\
793.00	94.45\\
844.00	94.19\\
897.00	94.31\\
955.00	93.93\\
1016.00	94.00\\
1080.00	94.17\\
1149.00	94.26\\
1222.00	94.35\\
1300.00	94.54\\
1383.00	94.50\\
1472.00	94.29\\
1565.00	94.19\\
1665.00	94.29\\
1771.00	94.18\\
1884.00	94.06\\
2005.00	94.11\\
2132.00	93.90\\
2268.00	94.09\\
2413.00	94.20\\
2567.00	94.20\\
2730.00	94.10\\
2905.00	93.98\\
3090.00	94.01\\
3287.00	94.10\\
3496.00	94.05\\
3719.00	93.98\\
3957.00	93.96\\
4209.00	93.80\\
4477.00	93.84\\
4763.00	93.81\\
5066.00	93.90\\
5389.00	93.91\\
5733.00	93.86\\
6099.00	93.97\\
6488.00	94.02\\
6901.00	94.02\\
7341.00	93.99\\
7809.00	93.94\\
8307.00	93.96\\
8837.00	93.96\\
9401.00	93.90\\
10000.00	93.97\\
};
\addplot [color=white!15!darkgray, dashdotted, line width=1.2pt, forget plot]
  table[row sep=crcr]{%
1.00	100.00\\
2.00	100.00\\
3.00	100.00\\
4.00	100.00\\
5.00	100.00\\
6.00	83.33\\
7.00	85.71\\
8.00	87.50\\
9.00	88.89\\
10.00	90.00\\
11.00	90.91\\
12.00	91.67\\
13.00	92.31\\
14.00	92.86\\
15.00	93.33\\
16.00	93.75\\
17.00	94.12\\
18.00	94.44\\
19.00	94.74\\
21.00	95.24\\
22.00	95.45\\
23.00	95.65\\
25.00	96.00\\
26.00	96.15\\
28.00	96.43\\
30.00	93.33\\
32.00	93.75\\
34.00	91.18\\
36.00	91.67\\
38.00	92.11\\
41.00	92.68\\
43.00	93.02\\
46.00	93.48\\
49.00	93.88\\
52.00	94.23\\
56.00	94.64\\
59.00	94.92\\
63.00	95.24\\
67.00	95.52\\
71.00	95.77\\
76.00	96.05\\
81.00	96.30\\
86.00	96.51\\
91.00	95.60\\
97.00	95.88\\
103.00	95.15\\
110.00	94.55\\
117.00	94.87\\
124.00	94.35\\
132.00	94.70\\
140.00	94.29\\
149.00	94.63\\
159.00	94.34\\
169.00	94.67\\
180.00	93.89\\
191.00	94.24\\
204.00	94.12\\
217.00	93.55\\
230.00	93.48\\
245.00	93.88\\
261.00	94.25\\
277.00	94.58\\
295.00	94.58\\
314.00	94.59\\
334.00	94.61\\
355.00	94.93\\
378.00	95.24\\
402.00	95.27\\
427.00	95.32\\
455.00	95.16\\
484.00	95.25\\
515.00	95.15\\
547.00	95.25\\
582.00	94.85\\
619.00	94.67\\
659.00	94.99\\
701.00	94.86\\
746.00	94.91\\
793.00	94.58\\
844.00	94.67\\
897.00	94.65\\
955.00	94.76\\
1016.00	94.39\\
1080.00	94.26\\
1149.00	94.34\\
1222.00	94.11\\
1300.00	94.31\\
1383.00	94.36\\
1472.00	94.29\\
1565.00	94.44\\
1665.00	94.41\\
1771.00	94.18\\
1884.00	94.11\\
2005.00	94.06\\
2132.00	93.95\\
2268.00	93.87\\
2413.00	93.87\\
2567.00	94.00\\
2730.00	94.21\\
2905.00	94.01\\
3090.00	93.85\\
3287.00	93.76\\
3496.00	93.88\\
3719.00	93.71\\
3957.00	93.48\\
4209.00	93.49\\
4477.00	93.50\\
4763.00	93.43\\
5066.00	93.09\\
5389.00	93.10\\
5733.00	92.92\\
6099.00	92.75\\
6488.00	92.66\\
6901.00	92.83\\
7341.00	92.74\\
7809.00	92.73\\
8307.00	92.58\\
8837.00	92.57\\
9401.00	92.54\\
10000.00	92.40\\
};
\addplot [color=white!45!black, line width=1.2pt, forget plot]
  table[row sep=crcr]{%
1.00	0.00\\
2.00	50.00\\
3.00	66.67\\
4.00	75.00\\
5.00	80.00\\
6.00	83.33\\
7.00	71.43\\
8.00	75.00\\
9.00	77.78\\
10.00	80.00\\
11.00	81.82\\
12.00	83.33\\
13.00	84.62\\
14.00	85.71\\
15.00	86.67\\
16.00	81.25\\
17.00	82.35\\
18.00	83.33\\
19.00	84.21\\
21.00	85.71\\
22.00	86.36\\
23.00	86.96\\
25.00	88.00\\
26.00	88.46\\
28.00	89.29\\
30.00	86.67\\
32.00	87.50\\
34.00	85.29\\
36.00	86.11\\
38.00	86.84\\
41.00	87.80\\
43.00	88.37\\
46.00	86.96\\
49.00	87.76\\
52.00	86.54\\
56.00	87.50\\
59.00	88.14\\
63.00	88.89\\
67.00	89.55\\
71.00	88.73\\
76.00	88.16\\
81.00	88.89\\
86.00	89.53\\
91.00	89.01\\
97.00	88.66\\
103.00	87.38\\
110.00	88.18\\
117.00	88.89\\
124.00	89.52\\
132.00	90.15\\
140.00	90.00\\
149.00	89.93\\
159.00	90.57\\
169.00	91.12\\
180.00	91.67\\
191.00	92.15\\
204.00	91.18\\
217.00	91.71\\
230.00	91.30\\
245.00	91.43\\
261.00	91.95\\
277.00	92.42\\
295.00	92.88\\
314.00	92.68\\
334.00	92.81\\
355.00	92.68\\
378.00	92.86\\
402.00	92.79\\
427.00	92.74\\
455.00	93.19\\
484.00	92.77\\
515.00	92.82\\
547.00	92.87\\
582.00	92.61\\
619.00	92.57\\
659.00	92.72\\
701.00	92.87\\
746.00	92.90\\
793.00	92.94\\
844.00	92.89\\
897.00	93.20\\
955.00	93.30\\
1016.00	93.31\\
1080.00	93.06\\
1149.00	93.21\\
1222.00	93.29\\
1300.00	93.23\\
1383.00	93.35\\
1472.00	93.21\\
1565.00	93.48\\
1665.00	93.27\\
1771.00	93.28\\
1884.00	93.58\\
2005.00	93.67\\
2132.00	93.76\\
2268.00	93.87\\
2413.00	93.87\\
2567.00	93.88\\
2730.00	93.99\\
2905.00	94.04\\
3090.00	93.98\\
3287.00	93.89\\
3496.00	93.62\\
3719.00	93.55\\
3957.00	93.63\\
4209.00	93.51\\
4477.00	93.54\\
4763.00	93.55\\
5066.00	93.60\\
5389.00	93.45\\
5733.00	93.42\\
6099.00	93.47\\
6488.00	93.50\\
6901.00	93.51\\
7341.00	93.56\\
7809.00	93.64\\
8307.00	93.72\\
8837.00	93.67\\
9401.00	93.67\\
10000.00	93.67\\
};
\addplot [color=gray!85!lightgray, dashed, line width=1.2pt, forget plot]
  table[row sep=crcr]{%
1.00	100.00\\
2.00	100.00\\
3.00	100.00\\
4.00	100.00\\
5.00	100.00\\
6.00	100.00\\
7.00	100.00\\
8.00	100.00\\
9.00	100.00\\
10.00	100.00\\
11.00	100.00\\
12.00	100.00\\
13.00	100.00\\
14.00	92.86\\
15.00	93.33\\
16.00	93.75\\
17.00	94.12\\
18.00	94.44\\
19.00	94.74\\
21.00	95.24\\
22.00	95.45\\
23.00	95.65\\
25.00	96.00\\
26.00	96.15\\
28.00	96.43\\
30.00	93.33\\
32.00	93.75\\
34.00	94.12\\
36.00	91.67\\
38.00	92.11\\
41.00	92.68\\
43.00	93.02\\
46.00	93.48\\
49.00	93.88\\
52.00	94.23\\
56.00	94.64\\
59.00	93.22\\
63.00	92.06\\
67.00	91.04\\
71.00	91.55\\
76.00	92.11\\
81.00	91.36\\
86.00	91.86\\
91.00	92.31\\
97.00	92.78\\
103.00	93.20\\
110.00	93.64\\
117.00	94.02\\
124.00	94.35\\
132.00	93.18\\
140.00	92.14\\
149.00	91.95\\
159.00	91.19\\
169.00	91.72\\
180.00	92.22\\
191.00	92.67\\
204.00	92.16\\
217.00	92.63\\
230.00	93.04\\
245.00	93.47\\
261.00	93.49\\
277.00	93.86\\
295.00	93.56\\
314.00	93.95\\
334.00	94.31\\
355.00	93.52\\
378.00	93.65\\
402.00	93.78\\
427.00	93.44\\
455.00	93.19\\
484.00	92.77\\
515.00	92.43\\
547.00	92.50\\
582.00	92.27\\
619.00	92.41\\
659.00	92.56\\
701.00	92.87\\
746.00	92.90\\
793.00	93.06\\
844.00	93.01\\
897.00	93.09\\
955.00	93.19\\
1016.00	93.11\\
1080.00	92.78\\
1149.00	92.95\\
1222.00	93.04\\
1300.00	93.00\\
1383.00	93.28\\
1472.00	93.21\\
1565.00	93.42\\
1665.00	93.45\\
1771.00	93.22\\
1884.00	93.26\\
2005.00	93.17\\
2132.00	93.34\\
2268.00	93.39\\
2413.00	93.20\\
2567.00	93.07\\
2730.00	93.26\\
2905.00	93.18\\
3090.00	93.33\\
3287.00	93.37\\
3496.00	93.22\\
3719.00	93.17\\
3957.00	93.28\\
4209.00	93.18\\
4477.00	92.99\\
4763.00	93.03\\
5066.00	93.05\\
5389.00	93.08\\
5733.00	93.09\\
6099.00	93.02\\
6488.00	92.86\\
6901.00	92.83\\
7341.00	92.81\\
7809.00	92.87\\
8307.00	92.84\\
8837.00	92.85\\
9401.00	92.79\\
10000.00	92.77\\
};
\addplot [color=white!50!darkgray, dotted, line width=1.2pt, forget plot]
  table[row sep=crcr]{%
1.00	100.00\\
2.00	100.00\\
3.00	100.00\\
4.00	100.00\\
5.00	100.00\\
6.00	100.00\\
7.00	100.00\\
8.00	100.00\\
9.00	100.00\\
10.00	100.00\\
11.00	100.00\\
12.00	100.00\\
13.00	100.00\\
14.00	100.00\\
15.00	100.00\\
16.00	100.00\\
17.00	100.00\\
18.00	100.00\\
19.00	94.74\\
21.00	95.24\\
22.00	95.45\\
23.00	95.65\\
25.00	96.00\\
26.00	96.15\\
28.00	92.86\\
30.00	90.00\\
32.00	90.62\\
34.00	91.18\\
36.00	91.67\\
38.00	92.11\\
41.00	90.24\\
43.00	90.70\\
46.00	91.30\\
49.00	91.84\\
52.00	92.31\\
56.00	92.86\\
59.00	93.22\\
63.00	93.65\\
67.00	94.03\\
71.00	92.96\\
76.00	93.42\\
81.00	93.83\\
86.00	94.19\\
91.00	94.51\\
97.00	93.81\\
103.00	93.20\\
110.00	92.73\\
117.00	93.16\\
124.00	93.55\\
132.00	93.18\\
140.00	93.57\\
149.00	93.96\\
159.00	93.71\\
169.00	93.49\\
180.00	93.33\\
191.00	92.67\\
204.00	92.65\\
217.00	92.63\\
230.00	93.04\\
245.00	93.47\\
261.00	93.10\\
277.00	93.50\\
295.00	93.56\\
314.00	93.95\\
334.00	94.31\\
355.00	94.37\\
378.00	94.71\\
402.00	95.02\\
427.00	94.61\\
455.00	94.95\\
484.00	94.63\\
515.00	94.56\\
547.00	94.70\\
582.00	94.50\\
619.00	94.67\\
659.00	94.69\\
701.00	94.72\\
746.00	94.50\\
793.00	94.33\\
844.00	94.31\\
897.00	93.53\\
955.00	93.40\\
1016.00	93.41\\
1080.00	93.33\\
1149.00	93.39\\
1222.00	93.62\\
1300.00	93.38\\
1383.00	93.56\\
1472.00	93.75\\
1565.00	93.74\\
1665.00	93.57\\
1771.00	93.56\\
1884.00	93.63\\
2005.00	93.72\\
2132.00	93.86\\
2268.00	93.56\\
2413.00	93.49\\
2567.00	93.53\\
2730.00	93.66\\
2905.00	93.73\\
3090.00	93.56\\
3287.00	93.55\\
3496.00	93.74\\
3719.00	93.65\\
3957.00	93.73\\
4209.00	93.61\\
4477.00	93.63\\
4763.00	93.66\\
5066.00	93.60\\
5389.00	93.47\\
5733.00	93.42\\
6099.00	93.29\\
6488.00	93.31\\
6901.00	93.38\\
7341.00	93.39\\
7809.00	93.37\\
8307.00	93.44\\
8837.00	93.43\\
9401.00	93.54\\
10000.00	93.54\\
};
\addplot [color=black!5!lightgray, dashdotted, line width=1.2pt, forget plot]
  table[row sep=crcr]{%
1.00	100.00\\
2.00	50.00\\
3.00	66.67\\
4.00	75.00\\
5.00	80.00\\
6.00	83.33\\
7.00	85.71\\
8.00	87.50\\
9.00	88.89\\
10.00	90.00\\
11.00	81.82\\
12.00	83.33\\
13.00	76.92\\
14.00	78.57\\
15.00	73.33\\
16.00	75.00\\
17.00	76.47\\
18.00	77.78\\
19.00	78.95\\
21.00	80.95\\
22.00	81.82\\
23.00	82.61\\
25.00	84.00\\
26.00	84.62\\
28.00	85.71\\
30.00	83.33\\
32.00	84.38\\
34.00	82.35\\
36.00	83.33\\
38.00	84.21\\
41.00	85.37\\
43.00	86.05\\
46.00	86.96\\
49.00	85.71\\
52.00	84.62\\
56.00	85.71\\
59.00	84.75\\
63.00	85.71\\
67.00	85.07\\
71.00	84.51\\
76.00	85.53\\
81.00	83.95\\
86.00	84.88\\
91.00	84.62\\
97.00	85.57\\
103.00	86.41\\
110.00	86.36\\
117.00	86.32\\
124.00	87.10\\
132.00	87.12\\
140.00	87.14\\
149.00	87.25\\
159.00	87.42\\
169.00	86.39\\
180.00	86.11\\
191.00	85.86\\
204.00	86.76\\
217.00	86.64\\
230.00	85.65\\
245.00	85.31\\
261.00	85.44\\
277.00	85.92\\
295.00	86.44\\
314.00	86.62\\
334.00	87.13\\
355.00	87.04\\
378.00	87.83\\
402.00	87.31\\
427.00	87.35\\
455.00	87.47\\
484.00	87.60\\
515.00	87.96\\
547.00	87.93\\
582.00	87.63\\
619.00	87.88\\
659.00	88.16\\
701.00	88.45\\
746.00	88.34\\
793.00	87.77\\
844.00	87.44\\
897.00	87.96\\
955.00	88.06\\
1016.00	88.19\\
1080.00	88.33\\
1149.00	88.60\\
1222.00	88.79\\
1300.00	89.08\\
1383.00	89.23\\
1472.00	89.20\\
1565.00	89.27\\
1665.00	89.37\\
1771.00	89.33\\
1884.00	89.28\\
2005.00	89.13\\
2132.00	89.21\\
2268.00	89.33\\
2413.00	89.47\\
2567.00	89.37\\
2730.00	89.49\\
2905.00	89.47\\
3090.00	89.45\\
3287.00	89.47\\
3496.00	89.42\\
3719.00	89.35\\
3957.00	89.49\\
4209.00	89.55\\
4477.00	89.61\\
4763.00	89.65\\
5066.00	89.54\\
5389.00	89.40\\
5733.00	89.26\\
6099.00	89.13\\
6488.00	89.06\\
6901.00	89.25\\
7341.00	89.16\\
7809.00	88.96\\
8307.00	89.07\\
8837.00	89.08\\
9401.00	89.14\\
10000.00	89.25\\
};
\addplot [color=white!80!black, line width=1.2pt, forget plot]
  table[row sep=crcr]{%
1.00	100.00\\
2.00	100.00\\
3.00	100.00\\
4.00	100.00\\
5.00	100.00\\
6.00	100.00\\
7.00	100.00\\
8.00	87.50\\
9.00	88.89\\
10.00	90.00\\
11.00	81.82\\
12.00	83.33\\
13.00	84.62\\
14.00	78.57\\
15.00	80.00\\
16.00	81.25\\
17.00	82.35\\
18.00	83.33\\
19.00	84.21\\
21.00	85.71\\
22.00	86.36\\
23.00	86.96\\
25.00	88.00\\
26.00	88.46\\
28.00	89.29\\
30.00	90.00\\
32.00	90.62\\
34.00	91.18\\
36.00	91.67\\
38.00	92.11\\
41.00	92.68\\
43.00	90.70\\
46.00	89.13\\
49.00	87.76\\
52.00	88.46\\
56.00	89.29\\
59.00	89.83\\
63.00	88.89\\
67.00	89.55\\
71.00	88.73\\
76.00	89.47\\
81.00	90.12\\
86.00	90.70\\
91.00	91.21\\
97.00	91.75\\
103.00	92.23\\
110.00	91.82\\
117.00	91.45\\
124.00	91.13\\
132.00	91.67\\
140.00	91.43\\
149.00	91.28\\
159.00	91.82\\
169.00	92.31\\
180.00	92.78\\
191.00	93.19\\
204.00	93.14\\
217.00	93.55\\
230.00	93.48\\
245.00	93.88\\
261.00	93.49\\
277.00	92.78\\
295.00	92.88\\
314.00	92.99\\
334.00	93.11\\
355.00	91.83\\
378.00	92.06\\
402.00	91.79\\
427.00	91.80\\
455.00	91.87\\
484.00	92.36\\
515.00	92.82\\
547.00	92.87\\
582.00	93.30\\
619.00	93.38\\
659.00	93.32\\
701.00	93.30\\
746.00	92.90\\
793.00	93.32\\
844.00	93.25\\
897.00	93.42\\
955.00	93.72\\
1016.00	93.31\\
1080.00	93.33\\
1149.00	93.39\\
1222.00	93.78\\
1300.00	93.85\\
1383.00	93.93\\
1472.00	93.89\\
1565.00	93.61\\
1665.00	93.75\\
1771.00	93.73\\
1884.00	93.63\\
2005.00	93.52\\
2132.00	93.76\\
2268.00	93.74\\
2413.00	93.78\\
2567.00	93.77\\
2730.00	93.81\\
2905.00	93.87\\
3090.00	93.85\\
3287.00	93.82\\
3496.00	93.94\\
3719.00	94.06\\
3957.00	94.04\\
4209.00	93.99\\
4477.00	93.88\\
4763.00	93.79\\
5066.00	93.66\\
5389.00	93.64\\
5733.00	93.63\\
6099.00	93.36\\
6488.00	93.31\\
6901.00	93.29\\
7341.00	93.22\\
7809.00	93.28\\
8307.00	93.27\\
8837.00	93.24\\
9401.00	93.17\\
10000.00	93.25\\
};
\end{axis}

\begin{axis}[%
width=0.391\figurewidth,
height=0.265\figureheight,
at={(0\figurewidth,0.368\figureheight)},
scale only axis,
xmin=1.00,
xmax=10000.00,
xtick={2000.00,4000.00,6000.00,8000.00,10000.00},
xticklabels={\empty},
ymin=0.00,
ymax=105.00,
ylabel style={font=\color{mycolor1}},
ylabel={\sym{power} [\%]},
axis background/.style={fill=white},
title style={font=\bfseries\color{mycolor1}},
title={\textbf{$\sym{beta}_2$}},
xmajorgrids,
ymajorgrids,
grid style={dotted, opacity=0.25}
]
\addplot [color=white!10!black, line width=1.2pt, forget plot]
  table[row sep=crcr]{%
1.00	100.00\\
2.00	100.00\\
3.00	100.00\\
4.00	100.00\\
5.00	100.00\\
6.00	100.00\\
7.00	100.00\\
8.00	100.00\\
9.00	100.00\\
10.00	100.00\\
11.00	100.00\\
12.00	100.00\\
13.00	100.00\\
14.00	100.00\\
15.00	100.00\\
16.00	100.00\\
17.00	100.00\\
18.00	100.00\\
19.00	100.00\\
21.00	100.00\\
22.00	100.00\\
23.00	100.00\\
25.00	100.00\\
26.00	100.00\\
28.00	100.00\\
30.00	100.00\\
32.00	100.00\\
34.00	100.00\\
36.00	100.00\\
38.00	100.00\\
41.00	97.56\\
43.00	97.67\\
46.00	97.83\\
49.00	97.96\\
52.00	98.08\\
56.00	98.21\\
59.00	94.92\\
63.00	95.24\\
67.00	95.52\\
71.00	95.77\\
76.00	94.74\\
81.00	93.83\\
86.00	94.19\\
91.00	94.51\\
97.00	94.85\\
103.00	94.17\\
110.00	93.64\\
117.00	94.02\\
124.00	94.35\\
132.00	94.70\\
140.00	95.00\\
149.00	95.30\\
159.00	94.97\\
169.00	94.08\\
180.00	93.89\\
191.00	93.19\\
204.00	93.14\\
217.00	92.17\\
230.00	92.17\\
245.00	92.24\\
261.00	92.72\\
277.00	92.78\\
295.00	93.22\\
314.00	92.99\\
334.00	92.81\\
355.00	92.68\\
378.00	93.12\\
402.00	93.28\\
427.00	93.68\\
455.00	93.85\\
484.00	93.80\\
515.00	94.17\\
547.00	93.97\\
582.00	93.81\\
619.00	94.02\\
659.00	93.93\\
701.00	93.72\\
746.00	93.30\\
793.00	93.06\\
844.00	92.77\\
897.00	92.98\\
955.00	92.98\\
1016.00	93.31\\
1080.00	93.43\\
1149.00	93.47\\
1222.00	93.45\\
1300.00	93.69\\
1383.00	93.78\\
1472.00	94.02\\
1565.00	93.74\\
1665.00	93.75\\
1771.00	93.68\\
1884.00	93.58\\
2005.00	93.67\\
2132.00	93.67\\
2268.00	93.78\\
2413.00	93.70\\
2567.00	93.84\\
2730.00	93.59\\
2905.00	93.49\\
3090.00	93.59\\
3287.00	93.61\\
3496.00	93.74\\
3719.00	93.60\\
3957.00	93.66\\
4209.00	93.87\\
4477.00	93.92\\
4763.00	93.93\\
5066.00	93.82\\
5389.00	93.84\\
5733.00	93.88\\
6099.00	93.80\\
6488.00	93.76\\
6901.00	93.81\\
7341.00	93.77\\
7809.00	93.71\\
8307.00	93.79\\
8837.00	93.80\\
9401.00	93.70\\
10000.00	93.69\\
};
\addplot [color=black!25!darkgray, dashed, line width=1.2pt, forget plot]
  table[row sep=crcr]{%
1.00	100.00\\
2.00	100.00\\
3.00	100.00\\
4.00	100.00\\
5.00	100.00\\
6.00	100.00\\
7.00	100.00\\
8.00	100.00\\
9.00	100.00\\
10.00	100.00\\
11.00	100.00\\
12.00	100.00\\
13.00	100.00\\
14.00	100.00\\
15.00	100.00\\
16.00	100.00\\
17.00	100.00\\
18.00	100.00\\
19.00	100.00\\
21.00	100.00\\
22.00	100.00\\
23.00	100.00\\
25.00	96.00\\
26.00	96.15\\
28.00	96.43\\
30.00	96.67\\
32.00	93.75\\
34.00	94.12\\
36.00	94.44\\
38.00	94.74\\
41.00	95.12\\
43.00	95.35\\
46.00	95.65\\
49.00	95.92\\
52.00	96.15\\
56.00	94.64\\
59.00	94.92\\
63.00	95.24\\
67.00	95.52\\
71.00	95.77\\
76.00	96.05\\
81.00	96.30\\
86.00	96.51\\
91.00	96.70\\
97.00	96.91\\
103.00	96.12\\
110.00	96.36\\
117.00	96.58\\
124.00	95.97\\
132.00	95.45\\
140.00	95.00\\
149.00	95.30\\
159.00	95.60\\
169.00	95.27\\
180.00	95.56\\
191.00	95.81\\
204.00	95.59\\
217.00	94.93\\
230.00	94.78\\
245.00	93.47\\
261.00	93.49\\
277.00	93.14\\
295.00	92.88\\
314.00	92.68\\
334.00	92.81\\
355.00	92.96\\
378.00	93.12\\
402.00	93.53\\
427.00	93.21\\
455.00	93.63\\
484.00	93.80\\
515.00	93.40\\
547.00	93.24\\
582.00	93.47\\
619.00	93.70\\
659.00	93.63\\
701.00	93.30\\
746.00	93.16\\
793.00	93.32\\
844.00	93.48\\
897.00	93.87\\
955.00	93.51\\
1016.00	93.50\\
1080.00	93.33\\
1149.00	93.30\\
1222.00	93.37\\
1300.00	93.54\\
1383.00	93.49\\
1472.00	93.48\\
1565.00	93.55\\
1665.00	93.75\\
1771.00	93.45\\
1884.00	93.68\\
2005.00	93.77\\
2132.00	93.90\\
2268.00	94.05\\
2413.00	94.07\\
2567.00	94.04\\
2730.00	94.03\\
2905.00	94.04\\
3090.00	93.92\\
3287.00	93.79\\
3496.00	93.74\\
3719.00	93.76\\
3957.00	93.66\\
4209.00	93.70\\
4477.00	93.70\\
4763.00	93.70\\
5066.00	93.82\\
5389.00	93.84\\
5733.00	93.74\\
6099.00	93.72\\
6488.00	93.76\\
6901.00	93.81\\
7341.00	93.83\\
7809.00	93.87\\
8307.00	93.85\\
8837.00	93.76\\
9401.00	93.77\\
10000.00	93.78\\
};
\addplot [color=black!45!gray, dotted, line width=1.2pt, forget plot]
  table[row sep=crcr]{%
1.00	100.00\\
2.00	100.00\\
3.00	100.00\\
4.00	100.00\\
5.00	100.00\\
6.00	100.00\\
7.00	100.00\\
8.00	100.00\\
9.00	100.00\\
10.00	100.00\\
11.00	100.00\\
12.00	100.00\\
13.00	100.00\\
14.00	100.00\\
15.00	100.00\\
16.00	100.00\\
17.00	100.00\\
18.00	100.00\\
19.00	100.00\\
21.00	100.00\\
22.00	100.00\\
23.00	100.00\\
25.00	96.00\\
26.00	96.15\\
28.00	92.86\\
30.00	93.33\\
32.00	93.75\\
34.00	94.12\\
36.00	94.44\\
38.00	94.74\\
41.00	95.12\\
43.00	95.35\\
46.00	95.65\\
49.00	95.92\\
52.00	94.23\\
56.00	94.64\\
59.00	94.92\\
63.00	93.65\\
67.00	94.03\\
71.00	94.37\\
76.00	94.74\\
81.00	95.06\\
86.00	95.35\\
91.00	93.41\\
97.00	93.81\\
103.00	94.17\\
110.00	94.55\\
117.00	94.02\\
124.00	94.35\\
132.00	94.70\\
140.00	95.00\\
149.00	94.63\\
159.00	94.97\\
169.00	95.27\\
180.00	95.00\\
191.00	94.76\\
204.00	94.61\\
217.00	94.47\\
230.00	94.35\\
245.00	94.69\\
261.00	94.64\\
277.00	94.95\\
295.00	95.25\\
314.00	94.90\\
334.00	94.61\\
355.00	94.65\\
378.00	94.44\\
402.00	94.03\\
427.00	94.15\\
455.00	94.07\\
484.00	93.80\\
515.00	93.98\\
547.00	93.97\\
582.00	93.81\\
619.00	93.70\\
659.00	93.78\\
701.00	93.72\\
746.00	93.57\\
793.00	93.57\\
844.00	93.48\\
897.00	93.42\\
955.00	93.40\\
1016.00	93.60\\
1080.00	93.70\\
1149.00	94.08\\
1222.00	94.11\\
1300.00	94.31\\
1383.00	94.14\\
1472.00	94.16\\
1565.00	94.06\\
1665.00	93.99\\
1771.00	94.01\\
1884.00	94.16\\
2005.00	94.31\\
2132.00	94.37\\
2268.00	94.22\\
2413.00	94.24\\
2567.00	94.00\\
2730.00	94.10\\
2905.00	94.08\\
3090.00	94.21\\
3287.00	94.07\\
3496.00	94.19\\
3719.00	94.17\\
3957.00	94.21\\
4209.00	94.25\\
4477.00	94.13\\
4763.00	94.06\\
5066.00	94.02\\
5389.00	94.10\\
5733.00	94.09\\
6099.00	94.15\\
6488.00	94.22\\
6901.00	94.22\\
7341.00	94.28\\
7809.00	94.30\\
8307.00	94.09\\
8837.00	94.10\\
9401.00	94.15\\
10000.00	94.14\\
};
\addplot [color=white!15!darkgray, dashdotted, line width=1.2pt, forget plot]
  table[row sep=crcr]{%
1.00	100.00\\
2.00	100.00\\
3.00	100.00\\
4.00	100.00\\
5.00	100.00\\
6.00	100.00\\
7.00	100.00\\
8.00	100.00\\
9.00	100.00\\
10.00	100.00\\
11.00	100.00\\
12.00	100.00\\
13.00	100.00\\
14.00	100.00\\
15.00	100.00\\
16.00	93.75\\
17.00	94.12\\
18.00	94.44\\
19.00	94.74\\
21.00	95.24\\
22.00	95.45\\
23.00	95.65\\
25.00	96.00\\
26.00	92.31\\
28.00	92.86\\
30.00	93.33\\
32.00	93.75\\
34.00	94.12\\
36.00	94.44\\
38.00	94.74\\
41.00	95.12\\
43.00	95.35\\
46.00	95.65\\
49.00	95.92\\
52.00	96.15\\
56.00	96.43\\
59.00	94.92\\
63.00	95.24\\
67.00	94.03\\
71.00	94.37\\
76.00	94.74\\
81.00	95.06\\
86.00	95.35\\
91.00	95.60\\
97.00	95.88\\
103.00	96.12\\
110.00	96.36\\
117.00	96.58\\
124.00	96.77\\
132.00	96.21\\
140.00	96.43\\
149.00	96.64\\
159.00	96.86\\
169.00	96.45\\
180.00	96.67\\
191.00	96.34\\
204.00	96.57\\
217.00	96.31\\
230.00	96.52\\
245.00	96.73\\
261.00	96.17\\
277.00	95.67\\
295.00	95.93\\
314.00	95.22\\
334.00	95.21\\
355.00	94.93\\
378.00	94.71\\
402.00	95.02\\
427.00	95.08\\
455.00	95.38\\
484.00	95.45\\
515.00	95.34\\
547.00	95.43\\
582.00	95.19\\
619.00	94.83\\
659.00	94.99\\
701.00	95.01\\
746.00	95.17\\
793.00	94.83\\
844.00	94.67\\
897.00	94.65\\
955.00	94.55\\
1016.00	94.69\\
1080.00	94.72\\
1149.00	94.43\\
1222.00	94.52\\
1300.00	94.46\\
1383.00	94.29\\
1472.00	93.95\\
1565.00	93.93\\
1665.00	93.87\\
1771.00	93.73\\
1884.00	93.74\\
2005.00	93.87\\
2132.00	93.76\\
2268.00	93.87\\
2413.00	93.83\\
2567.00	93.96\\
2730.00	94.03\\
2905.00	94.15\\
3090.00	94.11\\
3287.00	93.95\\
3496.00	94.02\\
3719.00	94.00\\
3957.00	93.78\\
4209.00	93.78\\
4477.00	93.81\\
4763.00	93.76\\
5066.00	93.78\\
5389.00	93.93\\
5733.00	93.89\\
6099.00	93.93\\
6488.00	93.88\\
6901.00	93.86\\
7341.00	93.91\\
7809.00	93.89\\
8307.00	93.93\\
8837.00	93.82\\
9401.00	93.87\\
10000.00	93.85\\
};
\addplot [color=white!45!black, line width=1.2pt, forget plot]
  table[row sep=crcr]{%
1.00	100.00\\
2.00	100.00\\
3.00	100.00\\
4.00	100.00\\
5.00	100.00\\
6.00	100.00\\
7.00	85.71\\
8.00	87.50\\
9.00	77.78\\
10.00	80.00\\
11.00	81.82\\
12.00	83.33\\
13.00	84.62\\
14.00	85.71\\
15.00	86.67\\
16.00	87.50\\
17.00	88.24\\
18.00	88.89\\
19.00	89.47\\
21.00	85.71\\
22.00	86.36\\
23.00	86.96\\
25.00	88.00\\
26.00	88.46\\
28.00	89.29\\
30.00	90.00\\
32.00	87.50\\
34.00	88.24\\
36.00	88.89\\
38.00	89.47\\
41.00	90.24\\
43.00	90.70\\
46.00	91.30\\
49.00	91.84\\
52.00	92.31\\
56.00	91.07\\
59.00	91.53\\
63.00	90.48\\
67.00	91.04\\
71.00	91.55\\
76.00	90.79\\
81.00	91.36\\
86.00	91.86\\
91.00	90.11\\
97.00	90.72\\
103.00	91.26\\
110.00	90.91\\
117.00	91.45\\
124.00	91.94\\
132.00	92.42\\
140.00	90.71\\
149.00	90.60\\
159.00	91.19\\
169.00	90.53\\
180.00	90.56\\
191.00	91.10\\
204.00	91.18\\
217.00	91.71\\
230.00	91.74\\
245.00	92.24\\
261.00	92.34\\
277.00	92.78\\
295.00	92.88\\
314.00	93.31\\
334.00	93.41\\
355.00	93.24\\
378.00	93.12\\
402.00	93.53\\
427.00	93.68\\
455.00	93.41\\
484.00	93.80\\
515.00	93.79\\
547.00	93.42\\
582.00	93.47\\
619.00	93.38\\
659.00	93.17\\
701.00	93.44\\
746.00	93.43\\
793.00	93.57\\
844.00	93.60\\
897.00	93.31\\
955.00	93.09\\
1016.00	92.91\\
1080.00	93.06\\
1149.00	93.30\\
1222.00	93.37\\
1300.00	93.23\\
1383.00	93.35\\
1472.00	93.48\\
1565.00	93.61\\
1665.00	93.57\\
1771.00	93.51\\
1884.00	93.58\\
2005.00	93.57\\
2132.00	93.34\\
2268.00	93.17\\
2413.00	92.83\\
2567.00	92.99\\
2730.00	92.78\\
2905.00	93.01\\
3090.00	93.07\\
3287.00	93.09\\
3496.00	92.96\\
3719.00	93.06\\
3957.00	93.20\\
4209.00	93.32\\
4477.00	93.28\\
4763.00	93.37\\
5066.00	93.43\\
5389.00	93.51\\
5733.00	93.55\\
6099.00	93.43\\
6488.00	93.50\\
6901.00	93.41\\
7341.00	93.46\\
7809.00	93.41\\
8307.00	93.37\\
8837.00	93.33\\
9401.00	93.37\\
10000.00	93.33\\
};
\addplot [color=gray!85!lightgray, dashed, line width=1.2pt, forget plot]
  table[row sep=crcr]{%
1.00	100.00\\
2.00	100.00\\
3.00	100.00\\
4.00	100.00\\
5.00	100.00\\
6.00	100.00\\
7.00	100.00\\
8.00	100.00\\
9.00	100.00\\
10.00	100.00\\
11.00	100.00\\
12.00	100.00\\
13.00	100.00\\
14.00	100.00\\
15.00	100.00\\
16.00	100.00\\
17.00	100.00\\
18.00	100.00\\
19.00	100.00\\
21.00	100.00\\
22.00	100.00\\
23.00	100.00\\
25.00	96.00\\
26.00	96.15\\
28.00	92.86\\
30.00	93.33\\
32.00	93.75\\
34.00	94.12\\
36.00	94.44\\
38.00	94.74\\
41.00	95.12\\
43.00	95.35\\
46.00	95.65\\
49.00	95.92\\
52.00	94.23\\
56.00	94.64\\
59.00	94.92\\
63.00	95.24\\
67.00	95.52\\
71.00	94.37\\
76.00	94.74\\
81.00	93.83\\
86.00	94.19\\
91.00	94.51\\
97.00	94.85\\
103.00	93.20\\
110.00	92.73\\
117.00	93.16\\
124.00	92.74\\
132.00	93.18\\
140.00	93.57\\
149.00	93.96\\
159.00	94.34\\
169.00	94.08\\
180.00	94.44\\
191.00	94.76\\
204.00	94.12\\
217.00	94.01\\
230.00	94.35\\
245.00	93.88\\
261.00	93.87\\
277.00	94.22\\
295.00	94.24\\
314.00	93.95\\
334.00	93.41\\
355.00	93.52\\
378.00	93.39\\
402.00	93.28\\
427.00	93.44\\
455.00	92.97\\
484.00	93.39\\
515.00	93.79\\
547.00	93.78\\
582.00	94.16\\
619.00	94.51\\
659.00	94.39\\
701.00	94.29\\
746.00	94.10\\
793.00	94.20\\
844.00	94.31\\
897.00	94.43\\
955.00	94.35\\
1016.00	94.09\\
1080.00	94.17\\
1149.00	94.17\\
1222.00	94.52\\
1300.00	94.46\\
1383.00	94.36\\
1472.00	94.29\\
1565.00	94.38\\
1665.00	94.29\\
1771.00	94.41\\
1884.00	94.48\\
2005.00	94.61\\
2132.00	94.70\\
2268.00	94.49\\
2413.00	94.53\\
2567.00	94.59\\
2730.00	94.54\\
2905.00	94.29\\
3090.00	94.27\\
3287.00	94.25\\
3496.00	94.34\\
3719.00	94.43\\
3957.00	94.39\\
4209.00	94.23\\
4477.00	94.13\\
4763.00	93.95\\
5066.00	94.02\\
5389.00	94.04\\
5733.00	94.05\\
6099.00	94.03\\
6488.00	93.91\\
6901.00	93.91\\
7341.00	93.98\\
7809.00	94.01\\
8307.00	94.07\\
8837.00	94.05\\
9401.00	93.97\\
10000.00	94.01\\
};
\addplot [color=white!50!darkgray, dotted, line width=1.2pt, forget plot]
  table[row sep=crcr]{%
1.00	0.00\\
2.00	50.00\\
3.00	66.67\\
4.00	75.00\\
5.00	80.00\\
6.00	83.33\\
7.00	85.71\\
8.00	87.50\\
9.00	88.89\\
10.00	90.00\\
11.00	90.91\\
12.00	91.67\\
13.00	92.31\\
14.00	92.86\\
15.00	93.33\\
16.00	93.75\\
17.00	94.12\\
18.00	94.44\\
19.00	94.74\\
21.00	95.24\\
22.00	95.45\\
23.00	95.65\\
25.00	96.00\\
26.00	96.15\\
28.00	96.43\\
30.00	93.33\\
32.00	93.75\\
34.00	94.12\\
36.00	94.44\\
38.00	94.74\\
41.00	92.68\\
43.00	90.70\\
46.00	91.30\\
49.00	91.84\\
52.00	90.38\\
56.00	91.07\\
59.00	89.83\\
63.00	90.48\\
67.00	91.04\\
71.00	91.55\\
76.00	92.11\\
81.00	91.36\\
86.00	91.86\\
91.00	92.31\\
97.00	91.75\\
103.00	92.23\\
110.00	91.82\\
117.00	91.45\\
124.00	91.13\\
132.00	91.67\\
140.00	91.43\\
149.00	91.28\\
159.00	91.82\\
169.00	92.31\\
180.00	92.78\\
191.00	93.19\\
204.00	93.14\\
217.00	92.63\\
230.00	92.17\\
245.00	92.24\\
261.00	92.72\\
277.00	92.78\\
295.00	93.22\\
314.00	93.63\\
334.00	93.71\\
355.00	94.08\\
378.00	94.44\\
402.00	94.78\\
427.00	94.61\\
455.00	94.51\\
484.00	94.42\\
515.00	94.56\\
547.00	94.52\\
582.00	94.50\\
619.00	94.51\\
659.00	94.69\\
701.00	94.58\\
746.00	94.77\\
793.00	94.96\\
844.00	94.91\\
897.00	94.87\\
955.00	94.87\\
1016.00	94.49\\
1080.00	94.63\\
1149.00	94.87\\
1222.00	94.60\\
1300.00	94.77\\
1383.00	94.58\\
1472.00	94.63\\
1565.00	94.89\\
1665.00	94.83\\
1771.00	94.64\\
1884.00	94.75\\
2005.00	94.66\\
2132.00	94.70\\
2268.00	94.71\\
2413.00	94.61\\
2567.00	94.51\\
2730.00	94.47\\
2905.00	94.29\\
3090.00	94.17\\
3287.00	94.31\\
3496.00	94.36\\
3719.00	94.27\\
3957.00	94.31\\
4209.00	94.30\\
4477.00	94.35\\
4763.00	94.27\\
5066.00	94.28\\
5389.00	94.27\\
5733.00	94.28\\
6099.00	94.24\\
6488.00	94.17\\
6901.00	94.16\\
7341.00	94.13\\
7809.00	94.15\\
8307.00	94.02\\
8837.00	94.07\\
9401.00	94.22\\
10000.00	94.19\\
};
\addplot [color=black!5!lightgray, dashdotted, line width=1.2pt, forget plot]
  table[row sep=crcr]{%
1.00	100.00\\
2.00	50.00\\
3.00	33.33\\
4.00	50.00\\
5.00	60.00\\
6.00	66.67\\
7.00	71.43\\
8.00	75.00\\
9.00	77.78\\
10.00	80.00\\
11.00	81.82\\
12.00	83.33\\
13.00	84.62\\
14.00	85.71\\
15.00	86.67\\
16.00	87.50\\
17.00	88.24\\
18.00	88.89\\
19.00	89.47\\
21.00	90.48\\
22.00	90.91\\
23.00	91.30\\
25.00	88.00\\
26.00	88.46\\
28.00	89.29\\
30.00	90.00\\
32.00	90.62\\
34.00	91.18\\
36.00	91.67\\
38.00	89.47\\
41.00	90.24\\
43.00	90.70\\
46.00	91.30\\
49.00	87.76\\
52.00	88.46\\
56.00	89.29\\
59.00	89.83\\
63.00	90.48\\
67.00	91.04\\
71.00	91.55\\
76.00	92.11\\
81.00	92.59\\
86.00	93.02\\
91.00	92.31\\
97.00	92.78\\
103.00	93.20\\
110.00	92.73\\
117.00	91.45\\
124.00	91.13\\
132.00	91.67\\
140.00	92.14\\
149.00	92.62\\
159.00	93.08\\
169.00	93.49\\
180.00	93.33\\
191.00	93.19\\
204.00	93.63\\
217.00	93.09\\
230.00	93.48\\
245.00	93.88\\
261.00	93.49\\
277.00	93.86\\
295.00	93.90\\
314.00	93.95\\
334.00	94.01\\
355.00	94.08\\
378.00	93.65\\
402.00	94.03\\
427.00	94.15\\
455.00	94.29\\
484.00	94.21\\
515.00	93.79\\
547.00	94.15\\
582.00	94.33\\
619.00	94.02\\
659.00	93.93\\
701.00	93.72\\
746.00	93.70\\
793.00	93.82\\
844.00	93.36\\
897.00	93.53\\
955.00	93.51\\
1016.00	93.60\\
1080.00	93.80\\
1149.00	93.65\\
1222.00	93.54\\
1300.00	93.46\\
1383.00	93.28\\
1472.00	93.48\\
1565.00	93.23\\
1665.00	93.39\\
1771.00	93.45\\
1884.00	93.74\\
2005.00	93.77\\
2132.00	93.48\\
2268.00	93.47\\
2413.00	93.54\\
2567.00	93.69\\
2730.00	93.63\\
2905.00	93.67\\
3090.00	93.79\\
3287.00	93.95\\
3496.00	93.76\\
3719.00	93.76\\
3957.00	93.78\\
4209.00	93.80\\
4477.00	93.72\\
4763.00	93.70\\
5066.00	93.70\\
5389.00	93.77\\
5733.00	93.83\\
6099.00	93.72\\
6488.00	93.76\\
6901.00	93.68\\
7341.00	93.73\\
7809.00	93.78\\
8307.00	93.84\\
8837.00	93.78\\
9401.00	93.79\\
10000.00	93.76\\
};
\addplot [color=white!80!black, line width=1.2pt, forget plot]
  table[row sep=crcr]{%
1.00	100.00\\
2.00	100.00\\
3.00	100.00\\
4.00	100.00\\
5.00	100.00\\
6.00	100.00\\
7.00	100.00\\
8.00	100.00\\
9.00	100.00\\
10.00	100.00\\
11.00	100.00\\
12.00	100.00\\
13.00	100.00\\
14.00	100.00\\
15.00	100.00\\
16.00	100.00\\
17.00	100.00\\
18.00	100.00\\
19.00	100.00\\
21.00	100.00\\
22.00	100.00\\
23.00	100.00\\
25.00	100.00\\
26.00	100.00\\
28.00	100.00\\
30.00	100.00\\
32.00	100.00\\
34.00	100.00\\
36.00	100.00\\
38.00	100.00\\
41.00	100.00\\
43.00	100.00\\
46.00	100.00\\
49.00	97.96\\
52.00	98.08\\
56.00	98.21\\
59.00	98.31\\
63.00	98.41\\
67.00	98.51\\
71.00	98.59\\
76.00	98.68\\
81.00	98.77\\
86.00	98.84\\
91.00	98.90\\
97.00	98.97\\
103.00	99.03\\
110.00	98.18\\
117.00	97.44\\
124.00	97.58\\
132.00	96.97\\
140.00	97.14\\
149.00	97.32\\
159.00	97.48\\
169.00	97.63\\
180.00	97.22\\
191.00	97.38\\
204.00	97.55\\
217.00	97.24\\
230.00	96.96\\
245.00	97.14\\
261.00	96.93\\
277.00	96.75\\
295.00	96.61\\
314.00	96.82\\
334.00	97.01\\
355.00	96.62\\
378.00	96.56\\
402.00	96.27\\
427.00	96.49\\
455.00	96.26\\
484.00	95.25\\
515.00	95.53\\
547.00	95.43\\
582.00	94.85\\
619.00	95.15\\
659.00	94.69\\
701.00	94.72\\
746.00	94.10\\
793.00	93.69\\
844.00	93.96\\
897.00	93.87\\
955.00	93.61\\
1016.00	93.70\\
1080.00	93.61\\
1149.00	93.47\\
1222.00	93.29\\
1300.00	93.15\\
1383.00	93.49\\
1472.00	93.41\\
1565.00	93.23\\
1665.00	93.39\\
1771.00	93.51\\
1884.00	93.26\\
2005.00	93.07\\
2132.00	93.11\\
2268.00	93.12\\
2413.00	93.08\\
2567.00	93.03\\
2730.00	93.11\\
2905.00	93.12\\
3090.00	93.24\\
3287.00	93.43\\
3496.00	93.42\\
3719.00	93.49\\
3957.00	93.58\\
4209.00	93.70\\
4477.00	93.57\\
4763.00	93.53\\
5066.00	93.66\\
5389.00	93.75\\
5733.00	93.70\\
6099.00	93.62\\
6488.00	93.68\\
6901.00	93.57\\
7341.00	93.61\\
7809.00	93.58\\
8307.00	93.57\\
8837.00	93.55\\
9401.00	93.62\\
10000.00	93.64\\
};
\end{axis}

\begin{axis}[%
width=0.391\figurewidth,
height=0.265\figureheight,
at={(0.54\figurewidth,0.368\figureheight)},
scale only axis,
xmin=1.00,
xmax=10000.00,
xtick={2000.00,4000.00,6000.00,8000.00,10000.00},
xticklabels={\empty},
ymin=0.00,
ymax=105.00,
ytick={0.00,20.00,40.00,60.00,80.00,100.00},
yticklabels={\empty},
axis background/.style={fill=white},
title style={font=\bfseries\color{mycolor1}},
title={\textbf{$\sym{beta}_{12}$}},
xmajorgrids,
ymajorgrids,
grid style={dotted, opacity=0.25}
]
\addplot [color=white!10!black, line width=1.2pt, forget plot]
  table[row sep=crcr]{%
1.00	100.00\\
2.00	100.00\\
3.00	100.00\\
4.00	75.00\\
5.00	60.00\\
6.00	66.67\\
7.00	57.14\\
8.00	50.00\\
9.00	44.44\\
10.00	50.00\\
11.00	45.45\\
12.00	41.67\\
13.00	38.46\\
14.00	42.86\\
15.00	40.00\\
16.00	37.50\\
17.00	41.18\\
18.00	38.89\\
19.00	42.11\\
21.00	47.62\\
22.00	50.00\\
23.00	52.17\\
25.00	56.00\\
26.00	57.69\\
28.00	60.71\\
30.00	63.33\\
32.00	62.50\\
34.00	58.82\\
36.00	58.33\\
38.00	57.89\\
41.00	56.10\\
43.00	55.81\\
46.00	58.70\\
49.00	59.18\\
52.00	59.62\\
56.00	58.93\\
59.00	57.63\\
63.00	55.56\\
67.00	56.72\\
71.00	54.93\\
76.00	52.63\\
81.00	51.85\\
86.00	54.65\\
91.00	57.14\\
97.00	58.76\\
103.00	58.25\\
110.00	56.36\\
117.00	58.12\\
124.00	58.87\\
132.00	59.09\\
140.00	59.29\\
149.00	59.06\\
159.00	58.49\\
169.00	57.40\\
180.00	57.78\\
191.00	57.59\\
204.00	57.84\\
217.00	57.60\\
230.00	58.70\\
245.00	58.78\\
261.00	58.24\\
277.00	59.21\\
295.00	59.32\\
314.00	59.87\\
334.00	59.88\\
355.00	60.00\\
378.00	59.79\\
402.00	59.95\\
427.00	60.19\\
455.00	59.34\\
484.00	58.88\\
515.00	58.06\\
547.00	59.05\\
582.00	58.59\\
619.00	58.32\\
659.00	57.51\\
701.00	57.49\\
746.00	56.70\\
793.00	56.49\\
844.00	56.04\\
897.00	55.96\\
955.00	55.81\\
1016.00	55.81\\
1080.00	55.83\\
1149.00	55.70\\
1222.00	55.89\\
1300.00	56.23\\
1383.00	56.47\\
1472.00	56.52\\
1565.00	56.36\\
1665.00	56.64\\
1771.00	56.30\\
1884.00	56.58\\
2005.00	56.51\\
2132.00	57.08\\
2268.00	57.19\\
2413.00	56.94\\
2567.00	56.72\\
2730.00	56.63\\
2905.00	56.59\\
3090.00	56.83\\
3287.00	56.89\\
3496.00	56.75\\
3719.00	56.82\\
3957.00	56.84\\
4209.00	56.90\\
4477.00	57.02\\
4763.00	57.06\\
5066.00	56.61\\
5389.00	56.24\\
5733.00	56.36\\
6099.00	56.16\\
6488.00	55.98\\
6901.00	55.99\\
7341.00	56.11\\
7809.00	56.18\\
8307.00	56.24\\
8837.00	56.11\\
9401.00	56.17\\
10000.00	56.21\\
};
\addplot [color=black!25!darkgray, dashed, line width=1.2pt, forget plot]
  table[row sep=crcr]{%
1.00	100.00\\
2.00	50.00\\
3.00	66.67\\
4.00	75.00\\
5.00	60.00\\
6.00	66.67\\
7.00	71.43\\
8.00	75.00\\
9.00	66.67\\
10.00	70.00\\
11.00	72.73\\
12.00	75.00\\
13.00	69.23\\
14.00	71.43\\
15.00	73.33\\
16.00	75.00\\
17.00	76.47\\
18.00	77.78\\
19.00	73.68\\
21.00	76.19\\
22.00	72.73\\
23.00	73.91\\
25.00	72.00\\
26.00	69.23\\
28.00	71.43\\
30.00	70.00\\
32.00	65.62\\
34.00	64.71\\
36.00	63.89\\
38.00	65.79\\
41.00	65.85\\
43.00	65.12\\
46.00	67.39\\
49.00	63.27\\
52.00	63.46\\
56.00	62.50\\
59.00	59.32\\
63.00	57.14\\
67.00	53.73\\
71.00	53.52\\
76.00	51.32\\
81.00	53.09\\
86.00	52.33\\
91.00	52.75\\
97.00	52.58\\
103.00	50.49\\
110.00	50.91\\
117.00	50.43\\
124.00	51.61\\
132.00	50.00\\
140.00	49.29\\
149.00	47.65\\
159.00	46.54\\
169.00	48.52\\
180.00	50.00\\
191.00	50.26\\
204.00	51.96\\
217.00	51.61\\
230.00	52.61\\
245.00	52.65\\
261.00	52.11\\
277.00	53.43\\
295.00	52.54\\
314.00	52.87\\
334.00	54.19\\
355.00	54.08\\
378.00	53.70\\
402.00	54.23\\
427.00	54.33\\
455.00	54.51\\
484.00	54.75\\
515.00	54.76\\
547.00	54.30\\
582.00	55.15\\
619.00	55.57\\
659.00	55.69\\
701.00	55.35\\
746.00	55.50\\
793.00	55.61\\
844.00	54.74\\
897.00	55.18\\
955.00	55.29\\
1016.00	55.12\\
1080.00	55.56\\
1149.00	55.87\\
1222.00	55.89\\
1300.00	55.31\\
1383.00	55.39\\
1472.00	55.64\\
1565.00	55.59\\
1665.00	55.74\\
1771.00	55.62\\
1884.00	56.16\\
2005.00	56.16\\
2132.00	56.29\\
2268.00	56.48\\
2413.00	56.44\\
2567.00	56.49\\
2730.00	56.67\\
2905.00	56.52\\
3090.00	56.44\\
3287.00	56.37\\
3496.00	56.29\\
3719.00	56.36\\
3957.00	56.38\\
4209.00	56.50\\
4477.00	56.31\\
4763.00	56.18\\
5066.00	56.20\\
5389.00	56.36\\
5733.00	56.51\\
6099.00	56.71\\
6488.00	56.58\\
6901.00	56.56\\
7341.00	56.55\\
7809.00	56.68\\
8307.00	56.48\\
8837.00	56.35\\
9401.00	56.29\\
10000.00	56.50\\
};
\addplot [color=black!45!gray, dotted, line width=1.2pt, forget plot]
  table[row sep=crcr]{%
1.00	100.00\\
2.00	100.00\\
3.00	100.00\\
4.00	100.00\\
5.00	100.00\\
6.00	100.00\\
7.00	85.71\\
8.00	87.50\\
9.00	88.89\\
10.00	90.00\\
11.00	81.82\\
12.00	83.33\\
13.00	84.62\\
14.00	78.57\\
15.00	80.00\\
16.00	75.00\\
17.00	76.47\\
18.00	72.22\\
19.00	73.68\\
21.00	71.43\\
22.00	72.73\\
23.00	69.57\\
25.00	72.00\\
26.00	73.08\\
28.00	67.86\\
30.00	66.67\\
32.00	65.62\\
34.00	64.71\\
36.00	63.89\\
38.00	65.79\\
41.00	68.29\\
43.00	69.77\\
46.00	65.22\\
49.00	67.35\\
52.00	69.23\\
56.00	69.64\\
59.00	69.49\\
63.00	68.25\\
67.00	68.66\\
71.00	70.42\\
76.00	71.05\\
81.00	71.60\\
86.00	72.09\\
91.00	71.43\\
97.00	70.10\\
103.00	67.96\\
110.00	66.36\\
117.00	64.10\\
124.00	63.71\\
132.00	63.64\\
140.00	63.57\\
149.00	62.42\\
159.00	62.26\\
169.00	63.91\\
180.00	62.22\\
191.00	62.30\\
204.00	62.75\\
217.00	63.13\\
230.00	62.61\\
245.00	64.08\\
261.00	63.98\\
277.00	63.54\\
295.00	61.02\\
314.00	60.51\\
334.00	61.08\\
355.00	61.13\\
378.00	60.85\\
402.00	60.95\\
427.00	60.42\\
455.00	60.00\\
484.00	59.30\\
515.00	57.86\\
547.00	58.50\\
582.00	58.42\\
619.00	59.29\\
659.00	58.57\\
701.00	58.49\\
746.00	58.85\\
793.00	58.01\\
844.00	58.29\\
897.00	57.97\\
955.00	58.53\\
1016.00	58.07\\
1080.00	58.06\\
1149.00	57.53\\
1222.00	57.69\\
1300.00	57.54\\
1383.00	57.77\\
1472.00	57.81\\
1565.00	57.32\\
1665.00	58.02\\
1771.00	58.50\\
1884.00	58.01\\
2005.00	58.15\\
2132.00	57.83\\
2268.00	57.76\\
2413.00	57.73\\
2567.00	57.50\\
2730.00	57.47\\
2905.00	57.21\\
3090.00	57.15\\
3287.00	56.86\\
3496.00	57.21\\
3719.00	57.19\\
3957.00	57.29\\
4209.00	57.16\\
4477.00	57.27\\
4763.00	57.15\\
5066.00	57.07\\
5389.00	57.51\\
5733.00	57.54\\
6099.00	57.71\\
6488.00	57.40\\
6901.00	57.21\\
7341.00	57.16\\
7809.00	57.36\\
8307.00	57.13\\
8837.00	57.08\\
9401.00	57.13\\
10000.00	57.14\\
};
\addplot [color=white!15!darkgray, dashdotted, line width=1.2pt, forget plot]
  table[row sep=crcr]{%
1.00	0.00\\
2.00	0.00\\
3.00	0.00\\
4.00	25.00\\
5.00	20.00\\
6.00	16.67\\
7.00	28.57\\
8.00	25.00\\
9.00	22.22\\
10.00	20.00\\
11.00	18.18\\
12.00	25.00\\
13.00	23.08\\
14.00	28.57\\
15.00	33.33\\
16.00	31.25\\
17.00	35.29\\
18.00	38.89\\
19.00	42.11\\
21.00	42.86\\
22.00	40.91\\
23.00	39.13\\
25.00	40.00\\
26.00	38.46\\
28.00	35.71\\
30.00	40.00\\
32.00	40.62\\
34.00	41.18\\
36.00	41.67\\
38.00	44.74\\
41.00	46.34\\
43.00	46.51\\
46.00	45.65\\
49.00	46.94\\
52.00	46.15\\
56.00	44.64\\
59.00	47.46\\
63.00	47.62\\
67.00	47.76\\
71.00	49.30\\
76.00	47.37\\
81.00	46.91\\
86.00	46.51\\
91.00	47.25\\
97.00	49.48\\
103.00	49.51\\
110.00	50.91\\
117.00	50.43\\
124.00	51.61\\
132.00	50.76\\
140.00	51.43\\
149.00	50.34\\
159.00	51.57\\
169.00	52.07\\
180.00	50.00\\
191.00	50.26\\
204.00	51.96\\
217.00	52.07\\
230.00	52.17\\
245.00	52.65\\
261.00	50.96\\
277.00	52.35\\
295.00	53.56\\
314.00	52.87\\
334.00	52.40\\
355.00	52.96\\
378.00	53.17\\
402.00	52.74\\
427.00	52.69\\
455.00	52.09\\
484.00	52.07\\
515.00	52.04\\
547.00	51.92\\
582.00	50.86\\
619.00	51.05\\
659.00	51.75\\
701.00	52.64\\
746.00	53.35\\
793.00	53.47\\
844.00	53.44\\
897.00	52.95\\
955.00	52.77\\
1016.00	52.17\\
1080.00	52.69\\
1149.00	52.83\\
1222.00	53.27\\
1300.00	53.54\\
1383.00	54.16\\
1472.00	53.74\\
1565.00	53.74\\
1665.00	53.63\\
1771.00	53.81\\
1884.00	53.72\\
2005.00	53.97\\
2132.00	54.50\\
2268.00	54.45\\
2413.00	54.33\\
2567.00	53.92\\
2730.00	54.25\\
2905.00	54.35\\
3090.00	54.47\\
3287.00	54.61\\
3496.00	54.81\\
3719.00	54.77\\
3957.00	54.44\\
4209.00	54.67\\
4477.00	54.90\\
4763.00	55.15\\
5066.00	55.27\\
5389.00	55.43\\
5733.00	55.29\\
6099.00	55.42\\
6488.00	55.69\\
6901.00	55.76\\
7341.00	55.97\\
7809.00	56.01\\
8307.00	55.94\\
8837.00	56.09\\
9401.00	56.14\\
10000.00	55.98\\
};
\addplot [color=white!45!black, line width=1.2pt, forget plot]
  table[row sep=crcr]{%
1.00	0.00\\
2.00	50.00\\
3.00	66.67\\
4.00	50.00\\
5.00	60.00\\
6.00	66.67\\
7.00	57.14\\
8.00	62.50\\
9.00	66.67\\
10.00	60.00\\
11.00	63.64\\
12.00	66.67\\
13.00	69.23\\
14.00	64.29\\
15.00	60.00\\
16.00	62.50\\
17.00	64.71\\
18.00	61.11\\
19.00	57.89\\
21.00	57.14\\
22.00	54.55\\
23.00	56.52\\
25.00	60.00\\
26.00	61.54\\
28.00	64.29\\
30.00	66.67\\
32.00	62.50\\
34.00	61.76\\
36.00	63.89\\
38.00	63.16\\
41.00	63.41\\
43.00	62.79\\
46.00	60.87\\
49.00	57.14\\
52.00	57.69\\
56.00	58.93\\
59.00	59.32\\
63.00	60.32\\
67.00	58.21\\
71.00	57.75\\
76.00	56.58\\
81.00	55.56\\
86.00	58.14\\
91.00	56.04\\
97.00	56.70\\
103.00	57.28\\
110.00	58.18\\
117.00	55.56\\
124.00	55.65\\
132.00	56.06\\
140.00	55.71\\
149.00	54.36\\
159.00	54.09\\
169.00	54.44\\
180.00	54.44\\
191.00	54.97\\
204.00	55.39\\
217.00	56.22\\
230.00	55.65\\
245.00	55.51\\
261.00	55.94\\
277.00	55.96\\
295.00	55.59\\
314.00	54.78\\
334.00	54.79\\
355.00	55.21\\
378.00	55.56\\
402.00	56.47\\
427.00	57.38\\
455.00	56.70\\
484.00	57.02\\
515.00	56.89\\
547.00	57.40\\
582.00	57.73\\
619.00	58.32\\
659.00	58.27\\
701.00	57.77\\
746.00	56.97\\
793.00	56.87\\
844.00	56.87\\
897.00	56.74\\
955.00	56.75\\
1016.00	56.40\\
1080.00	55.83\\
1149.00	55.79\\
1222.00	56.22\\
1300.00	56.23\\
1383.00	56.54\\
1472.00	55.98\\
1565.00	55.97\\
1665.00	56.04\\
1771.00	56.41\\
1884.00	56.63\\
2005.00	56.66\\
2132.00	56.24\\
2268.00	55.78\\
2413.00	55.66\\
2567.00	55.86\\
2730.00	56.01\\
2905.00	55.49\\
3090.00	55.53\\
3287.00	55.83\\
3496.00	55.75\\
3719.00	55.47\\
3957.00	55.24\\
4209.00	55.33\\
4477.00	55.15\\
4763.00	55.45\\
5066.00	55.35\\
5389.00	55.48\\
5733.00	55.71\\
6099.00	55.99\\
6488.00	56.18\\
6901.00	56.30\\
7341.00	56.19\\
7809.00	56.18\\
8307.00	56.10\\
8837.00	55.97\\
9401.00	56.15\\
10000.00	56.15\\
};
\addplot [color=gray!85!lightgray, dashed, line width=1.2pt, forget plot]
  table[row sep=crcr]{%
1.00	0.00\\
2.00	50.00\\
3.00	66.67\\
4.00	75.00\\
5.00	60.00\\
6.00	66.67\\
7.00	57.14\\
8.00	62.50\\
9.00	55.56\\
10.00	60.00\\
11.00	63.64\\
12.00	66.67\\
13.00	61.54\\
14.00	64.29\\
15.00	66.67\\
16.00	68.75\\
17.00	70.59\\
18.00	72.22\\
19.00	73.68\\
21.00	66.67\\
22.00	68.18\\
23.00	65.22\\
25.00	64.00\\
26.00	65.38\\
28.00	67.86\\
30.00	66.67\\
32.00	68.75\\
34.00	70.59\\
36.00	66.67\\
38.00	65.79\\
41.00	63.41\\
43.00	62.79\\
46.00	63.04\\
49.00	59.18\\
52.00	61.54\\
56.00	62.50\\
59.00	59.32\\
63.00	61.90\\
67.00	61.19\\
71.00	61.97\\
76.00	61.84\\
81.00	61.73\\
86.00	61.63\\
91.00	60.44\\
97.00	59.79\\
103.00	60.19\\
110.00	60.00\\
117.00	60.68\\
124.00	58.87\\
132.00	58.33\\
140.00	60.00\\
149.00	59.06\\
159.00	61.01\\
169.00	59.17\\
180.00	60.00\\
191.00	59.69\\
204.00	58.82\\
217.00	58.53\\
230.00	60.00\\
245.00	60.00\\
261.00	58.62\\
277.00	59.57\\
295.00	60.00\\
314.00	59.87\\
334.00	58.98\\
355.00	57.46\\
378.00	57.67\\
402.00	57.71\\
427.00	58.55\\
455.00	58.46\\
484.00	59.30\\
515.00	59.03\\
547.00	58.87\\
582.00	58.76\\
619.00	59.61\\
659.00	59.18\\
701.00	58.63\\
746.00	58.18\\
793.00	58.01\\
844.00	57.58\\
897.00	57.64\\
955.00	57.59\\
1016.00	57.28\\
1080.00	57.22\\
1149.00	57.18\\
1222.00	57.36\\
1300.00	57.85\\
1383.00	57.56\\
1472.00	57.54\\
1565.00	57.76\\
1665.00	57.54\\
1771.00	57.26\\
1884.00	57.43\\
2005.00	58.30\\
2132.00	58.40\\
2268.00	58.47\\
2413.00	58.97\\
2567.00	59.14\\
2730.00	59.34\\
2905.00	59.83\\
3090.00	59.61\\
3287.00	59.93\\
3496.00	60.07\\
3719.00	60.15\\
3957.00	60.20\\
4209.00	60.16\\
4477.00	59.91\\
4763.00	59.90\\
5066.00	60.03\\
5389.00	60.18\\
5733.00	60.42\\
6099.00	60.22\\
6488.00	60.16\\
6901.00	60.01\\
7341.00	60.05\\
7809.00	60.19\\
8307.00	60.21\\
8837.00	60.34\\
9401.00	60.37\\
10000.00	60.46\\
};
\addplot [color=white!50!darkgray, dotted, line width=1.2pt, forget plot]
  table[row sep=crcr]{%
1.00	100.00\\
2.00	100.00\\
3.00	100.00\\
4.00	100.00\\
5.00	80.00\\
6.00	83.33\\
7.00	85.71\\
8.00	75.00\\
9.00	66.67\\
10.00	60.00\\
11.00	54.55\\
12.00	58.33\\
13.00	53.85\\
14.00	57.14\\
15.00	53.33\\
16.00	50.00\\
17.00	52.94\\
18.00	55.56\\
19.00	52.63\\
21.00	52.38\\
22.00	54.55\\
23.00	52.17\\
25.00	52.00\\
26.00	50.00\\
28.00	50.00\\
30.00	46.67\\
32.00	43.75\\
34.00	47.06\\
36.00	47.22\\
38.00	50.00\\
41.00	51.22\\
43.00	51.16\\
46.00	50.00\\
49.00	48.98\\
52.00	50.00\\
56.00	48.21\\
59.00	49.15\\
63.00	50.79\\
67.00	50.75\\
71.00	50.70\\
76.00	51.32\\
81.00	50.62\\
86.00	51.16\\
91.00	50.55\\
97.00	50.52\\
103.00	53.40\\
110.00	53.64\\
117.00	53.85\\
124.00	55.65\\
132.00	54.55\\
140.00	52.86\\
149.00	51.68\\
159.00	52.83\\
169.00	53.25\\
180.00	53.89\\
191.00	52.36\\
204.00	51.47\\
217.00	51.61\\
230.00	52.17\\
245.00	51.02\\
261.00	50.19\\
277.00	50.18\\
295.00	50.17\\
314.00	50.32\\
334.00	51.80\\
355.00	52.39\\
378.00	52.91\\
402.00	53.98\\
427.00	54.33\\
455.00	54.51\\
484.00	54.96\\
515.00	54.37\\
547.00	54.48\\
582.00	54.47\\
619.00	54.77\\
659.00	54.63\\
701.00	54.78\\
746.00	56.03\\
793.00	55.86\\
844.00	56.16\\
897.00	56.08\\
955.00	56.86\\
1016.00	57.28\\
1080.00	56.67\\
1149.00	56.48\\
1222.00	56.30\\
1300.00	55.92\\
1383.00	55.68\\
1472.00	56.25\\
1565.00	55.85\\
1665.00	55.68\\
1771.00	56.13\\
1884.00	56.21\\
2005.00	55.71\\
2132.00	56.10\\
2268.00	56.26\\
2413.00	56.24\\
2567.00	56.45\\
2730.00	57.11\\
2905.00	56.73\\
3090.00	56.96\\
3287.00	56.95\\
3496.00	56.81\\
3719.00	56.84\\
3957.00	57.01\\
4209.00	56.69\\
4477.00	56.58\\
4763.00	56.81\\
5066.00	56.65\\
5389.00	56.63\\
5733.00	56.51\\
6099.00	56.24\\
6488.00	56.32\\
6901.00	56.38\\
7341.00	56.41\\
7809.00	56.49\\
8307.00	56.41\\
8837.00	56.49\\
9401.00	56.37\\
10000.00	56.38\\
};
\addplot [color=black!5!lightgray, dashdotted, line width=1.2pt, forget plot]
  table[row sep=crcr]{%
1.00	0.00\\
2.00	0.00\\
3.00	33.33\\
4.00	50.00\\
5.00	40.00\\
6.00	33.33\\
7.00	42.86\\
8.00	50.00\\
9.00	44.44\\
10.00	50.00\\
11.00	54.55\\
12.00	50.00\\
13.00	46.15\\
14.00	50.00\\
15.00	53.33\\
16.00	50.00\\
17.00	52.94\\
18.00	55.56\\
19.00	57.89\\
21.00	61.90\\
22.00	59.09\\
23.00	56.52\\
25.00	56.00\\
26.00	53.85\\
28.00	53.57\\
30.00	56.67\\
32.00	56.25\\
34.00	58.82\\
36.00	61.11\\
38.00	60.53\\
41.00	58.54\\
43.00	55.81\\
46.00	54.35\\
49.00	55.10\\
52.00	55.77\\
56.00	53.57\\
59.00	52.54\\
63.00	53.97\\
67.00	56.72\\
71.00	56.34\\
76.00	55.26\\
81.00	55.56\\
86.00	55.81\\
91.00	53.85\\
97.00	52.58\\
103.00	52.43\\
110.00	53.64\\
117.00	54.70\\
124.00	56.45\\
132.00	56.06\\
140.00	57.14\\
149.00	58.39\\
159.00	57.86\\
169.00	56.21\\
180.00	56.11\\
191.00	53.93\\
204.00	53.43\\
217.00	54.84\\
230.00	53.91\\
245.00	54.29\\
261.00	53.26\\
277.00	53.43\\
295.00	54.24\\
314.00	53.50\\
334.00	54.49\\
355.00	55.21\\
378.00	55.03\\
402.00	54.73\\
427.00	54.80\\
455.00	54.73\\
484.00	54.96\\
515.00	54.76\\
547.00	54.66\\
582.00	54.64\\
619.00	54.60\\
659.00	55.08\\
701.00	54.49\\
746.00	54.29\\
793.00	54.60\\
844.00	54.50\\
897.00	55.41\\
955.00	54.97\\
1016.00	54.72\\
1080.00	55.09\\
1149.00	55.00\\
1222.00	55.32\\
1300.00	55.92\\
1383.00	55.75\\
1472.00	55.64\\
1565.00	56.10\\
1665.00	56.46\\
1771.00	56.35\\
1884.00	56.53\\
2005.00	56.31\\
2132.00	56.14\\
2268.00	56.53\\
2413.00	56.36\\
2567.00	56.88\\
2730.00	56.37\\
2905.00	56.76\\
3090.00	56.96\\
3287.00	56.37\\
3496.00	56.64\\
3719.00	56.68\\
3957.00	56.99\\
4209.00	57.09\\
4477.00	57.14\\
4763.00	57.09\\
5066.00	56.91\\
5389.00	56.76\\
5733.00	56.71\\
6099.00	56.67\\
6488.00	56.43\\
6901.00	56.56\\
7341.00	56.56\\
7809.00	56.59\\
8307.00	56.46\\
8837.00	56.48\\
9401.00	56.45\\
10000.00	56.37\\
};
\addplot [color=white!80!black, line width=1.2pt, forget plot]
  table[row sep=crcr]{%
1.00	100.00\\
2.00	100.00\\
3.00	66.67\\
4.00	75.00\\
5.00	80.00\\
6.00	66.67\\
7.00	57.14\\
8.00	50.00\\
9.00	44.44\\
10.00	40.00\\
11.00	45.45\\
12.00	41.67\\
13.00	46.15\\
14.00	42.86\\
15.00	46.67\\
16.00	43.75\\
17.00	47.06\\
18.00	50.00\\
19.00	47.37\\
21.00	47.62\\
22.00	50.00\\
23.00	47.83\\
25.00	48.00\\
26.00	50.00\\
28.00	50.00\\
30.00	53.33\\
32.00	56.25\\
34.00	55.88\\
36.00	52.78\\
38.00	55.26\\
41.00	53.66\\
43.00	53.49\\
46.00	54.35\\
49.00	55.10\\
52.00	55.77\\
56.00	58.93\\
59.00	57.63\\
63.00	55.56\\
67.00	53.73\\
71.00	54.93\\
76.00	55.26\\
81.00	55.56\\
86.00	55.81\\
91.00	57.14\\
97.00	55.67\\
103.00	57.28\\
110.00	57.27\\
117.00	58.12\\
124.00	58.06\\
132.00	58.33\\
140.00	57.86\\
149.00	57.05\\
159.00	57.86\\
169.00	58.58\\
180.00	58.89\\
191.00	59.16\\
204.00	59.80\\
217.00	60.83\\
230.00	60.87\\
245.00	60.00\\
261.00	59.00\\
277.00	59.21\\
295.00	58.31\\
314.00	59.87\\
334.00	58.68\\
355.00	58.87\\
378.00	57.94\\
402.00	58.21\\
427.00	58.55\\
455.00	58.24\\
484.00	58.26\\
515.00	59.22\\
547.00	59.23\\
582.00	59.11\\
619.00	59.13\\
659.00	58.42\\
701.00	58.77\\
746.00	58.45\\
793.00	58.76\\
844.00	58.77\\
897.00	58.97\\
955.00	59.69\\
1016.00	59.74\\
1080.00	59.44\\
1149.00	59.62\\
1222.00	59.25\\
1300.00	59.46\\
1383.00	58.86\\
1472.00	59.10\\
1565.00	58.98\\
1665.00	58.92\\
1771.00	59.06\\
1884.00	58.86\\
2005.00	58.40\\
2132.00	58.49\\
2268.00	58.33\\
2413.00	58.02\\
2567.00	58.01\\
2730.00	57.88\\
2905.00	58.14\\
3090.00	57.93\\
3287.00	57.71\\
3496.00	58.07\\
3719.00	57.92\\
3957.00	57.34\\
4209.00	57.23\\
4477.00	56.96\\
4763.00	56.92\\
5066.00	57.13\\
5389.00	57.04\\
5733.00	56.99\\
6099.00	56.65\\
6488.00	56.67\\
6901.00	56.73\\
7341.00	56.68\\
7809.00	56.58\\
8307.00	56.78\\
8837.00	56.61\\
9401.00	56.49\\
10000.00	56.55\\
};
\end{axis}

\begin{axis}[%
width=0.391\figurewidth,
height=0.265\figureheight,
at={(0\figurewidth,0\figureheight)},
scale only axis,
xmin=1.00,
xmax=10000.00,
xlabel style={font=\color{mycolor1}},
xlabel={Iterationen $\sym{n}_{\sym{idx_MC}}$},
ymin=0.00,
ymax=105.00,
ylabel style={font=\color{mycolor1}},
ylabel={\sym{power} [\%]},
axis background/.style={fill=white},
title style={font=\bfseries\color{mycolor1}},
title={\textbf{$\sym{beta}_{11}$}},
xmajorgrids,
ymajorgrids,
grid style={dotted, opacity=0.25}
]
\addplot [color=white!10!black, line width=1.2pt, forget plot]
  table[row sep=crcr]{%
1.00	100.00\\
2.00	50.00\\
3.00	33.33\\
4.00	50.00\\
5.00	60.00\\
6.00	66.67\\
7.00	57.14\\
8.00	62.50\\
9.00	66.67\\
10.00	70.00\\
11.00	63.64\\
12.00	66.67\\
13.00	61.54\\
14.00	64.29\\
15.00	66.67\\
16.00	68.75\\
17.00	70.59\\
18.00	66.67\\
19.00	63.16\\
21.00	61.90\\
22.00	59.09\\
23.00	60.87\\
25.00	60.00\\
26.00	61.54\\
28.00	64.29\\
30.00	63.33\\
32.00	65.62\\
34.00	67.65\\
36.00	69.44\\
38.00	71.05\\
41.00	70.73\\
43.00	67.44\\
46.00	69.57\\
49.00	67.35\\
52.00	65.38\\
56.00	62.50\\
59.00	64.41\\
63.00	65.08\\
67.00	67.16\\
71.00	67.61\\
76.00	69.74\\
81.00	70.37\\
86.00	70.93\\
91.00	71.43\\
97.00	69.07\\
103.00	69.90\\
110.00	70.00\\
117.00	70.09\\
124.00	70.97\\
132.00	69.70\\
140.00	70.00\\
149.00	69.13\\
159.00	69.18\\
169.00	68.64\\
180.00	68.89\\
191.00	69.63\\
204.00	69.61\\
217.00	68.20\\
230.00	68.70\\
245.00	68.57\\
261.00	68.58\\
277.00	69.31\\
295.00	68.81\\
314.00	67.83\\
334.00	68.86\\
355.00	69.01\\
378.00	68.52\\
402.00	68.91\\
427.00	69.56\\
455.00	69.01\\
484.00	68.80\\
515.00	68.54\\
547.00	68.92\\
582.00	68.38\\
619.00	68.01\\
659.00	67.68\\
701.00	67.48\\
746.00	67.69\\
793.00	67.72\\
844.00	67.54\\
897.00	67.00\\
955.00	67.02\\
1016.00	66.54\\
1080.00	66.76\\
1149.00	66.67\\
1222.00	66.94\\
1300.00	66.77\\
1383.00	66.81\\
1472.00	67.46\\
1565.00	67.28\\
1665.00	67.27\\
1771.00	67.02\\
1884.00	67.30\\
2005.00	67.28\\
2132.00	67.45\\
2268.00	67.42\\
2413.00	67.43\\
2567.00	67.24\\
2730.00	67.51\\
2905.00	67.30\\
3090.00	67.38\\
3287.00	68.12\\
3496.00	67.96\\
3719.00	67.89\\
3957.00	67.78\\
4209.00	67.81\\
4477.00	68.10\\
4763.00	68.30\\
5066.00	68.12\\
5389.00	68.31\\
5733.00	68.53\\
6099.00	68.59\\
6488.00	68.53\\
6901.00	68.48\\
7341.00	68.53\\
7809.00	68.47\\
8307.00	68.48\\
8837.00	68.42\\
9401.00	68.54\\
10000.00	68.59\\
};
\addplot [color=black!25!darkgray, dashed, line width=1.2pt, forget plot]
  table[row sep=crcr]{%
1.00	100.00\\
2.00	100.00\\
3.00	66.67\\
4.00	75.00\\
5.00	60.00\\
6.00	66.67\\
7.00	71.43\\
8.00	75.00\\
9.00	77.78\\
10.00	70.00\\
11.00	63.64\\
12.00	66.67\\
13.00	61.54\\
14.00	64.29\\
15.00	66.67\\
16.00	68.75\\
17.00	70.59\\
18.00	72.22\\
19.00	73.68\\
21.00	71.43\\
22.00	72.73\\
23.00	73.91\\
25.00	72.00\\
26.00	69.23\\
28.00	71.43\\
30.00	70.00\\
32.00	71.88\\
34.00	73.53\\
36.00	72.22\\
38.00	73.68\\
41.00	73.17\\
43.00	74.42\\
46.00	71.74\\
49.00	69.39\\
52.00	71.15\\
56.00	69.64\\
59.00	71.19\\
63.00	69.84\\
67.00	70.15\\
71.00	71.83\\
76.00	71.05\\
81.00	71.60\\
86.00	72.09\\
91.00	72.53\\
97.00	71.13\\
103.00	69.90\\
110.00	69.09\\
117.00	68.38\\
124.00	67.74\\
132.00	69.70\\
140.00	67.86\\
149.00	67.11\\
159.00	67.30\\
169.00	67.46\\
180.00	66.67\\
191.00	66.49\\
204.00	66.67\\
217.00	66.82\\
230.00	67.39\\
245.00	67.76\\
261.00	67.05\\
277.00	66.79\\
295.00	66.10\\
314.00	65.61\\
334.00	65.87\\
355.00	65.92\\
378.00	66.14\\
402.00	66.67\\
427.00	66.74\\
455.00	67.69\\
484.00	66.94\\
515.00	67.77\\
547.00	67.82\\
582.00	68.21\\
619.00	67.69\\
659.00	67.37\\
701.00	67.33\\
746.00	67.02\\
793.00	67.21\\
844.00	67.42\\
897.00	67.56\\
955.00	67.23\\
1016.00	67.13\\
1080.00	67.04\\
1149.00	66.58\\
1222.00	66.69\\
1300.00	66.08\\
1383.00	65.58\\
1472.00	65.29\\
1565.00	65.11\\
1665.00	65.41\\
1771.00	65.10\\
1884.00	65.50\\
2005.00	65.19\\
2132.00	65.15\\
2268.00	64.95\\
2413.00	64.90\\
2567.00	64.59\\
2730.00	63.96\\
2905.00	63.92\\
3090.00	64.08\\
3287.00	63.86\\
3496.00	64.04\\
3719.00	64.16\\
3957.00	64.22\\
4209.00	64.31\\
4477.00	64.33\\
4763.00	64.37\\
5066.00	64.35\\
5389.00	64.39\\
5733.00	64.45\\
6099.00	64.40\\
6488.00	64.60\\
6901.00	64.47\\
7341.00	64.39\\
7809.00	64.52\\
8307.00	64.52\\
8837.00	64.58\\
9401.00	64.56\\
10000.00	64.51\\
};
\addplot [color=black!45!gray, dotted, line width=1.2pt, forget plot]
  table[row sep=crcr]{%
1.00	100.00\\
2.00	50.00\\
3.00	66.67\\
4.00	75.00\\
5.00	80.00\\
6.00	83.33\\
7.00	85.71\\
8.00	87.50\\
9.00	88.89\\
10.00	90.00\\
11.00	90.91\\
12.00	91.67\\
13.00	84.62\\
14.00	85.71\\
15.00	86.67\\
16.00	81.25\\
17.00	76.47\\
18.00	77.78\\
19.00	73.68\\
21.00	76.19\\
22.00	72.73\\
23.00	73.91\\
25.00	76.00\\
26.00	76.92\\
28.00	75.00\\
30.00	76.67\\
32.00	75.00\\
34.00	76.47\\
36.00	75.00\\
38.00	71.05\\
41.00	70.73\\
43.00	69.77\\
46.00	71.74\\
49.00	73.47\\
52.00	73.08\\
56.00	73.21\\
59.00	72.88\\
63.00	73.02\\
67.00	73.13\\
71.00	73.24\\
76.00	72.37\\
81.00	70.37\\
86.00	72.09\\
91.00	72.53\\
97.00	73.20\\
103.00	71.84\\
110.00	72.73\\
117.00	72.65\\
124.00	72.58\\
132.00	72.73\\
140.00	72.86\\
149.00	71.81\\
159.00	71.70\\
169.00	71.01\\
180.00	70.56\\
191.00	70.16\\
204.00	72.06\\
217.00	72.35\\
230.00	70.87\\
245.00	71.02\\
261.00	72.41\\
277.00	73.29\\
295.00	73.56\\
314.00	73.57\\
334.00	72.16\\
355.00	70.99\\
378.00	70.11\\
402.00	69.65\\
427.00	70.26\\
455.00	70.77\\
484.00	70.45\\
515.00	70.49\\
547.00	70.38\\
582.00	70.45\\
619.00	70.27\\
659.00	70.41\\
701.00	71.18\\
746.00	70.91\\
793.00	70.49\\
844.00	71.21\\
897.00	71.57\\
955.00	72.25\\
1016.00	72.15\\
1080.00	72.69\\
1149.00	72.67\\
1222.00	72.75\\
1300.00	72.69\\
1383.00	73.10\\
1472.00	73.44\\
1565.00	73.16\\
1665.00	72.97\\
1771.00	72.78\\
1884.00	73.04\\
2005.00	72.52\\
2132.00	72.70\\
2268.00	72.75\\
2413.00	73.23\\
2567.00	73.35\\
2730.00	73.59\\
2905.00	74.11\\
3090.00	73.95\\
3287.00	73.65\\
3496.00	73.66\\
3719.00	73.60\\
3957.00	73.29\\
4209.00	73.51\\
4477.00	73.67\\
4763.00	73.67\\
5066.00	73.85\\
5389.00	74.06\\
5733.00	74.18\\
6099.00	74.39\\
6488.00	74.38\\
6901.00	74.32\\
7341.00	74.30\\
7809.00	74.38\\
8307.00	74.32\\
8837.00	74.29\\
9401.00	74.16\\
10000.00	74.19\\
};
\addplot [color=white!15!darkgray, dashdotted, line width=1.2pt, forget plot]
  table[row sep=crcr]{%
1.00	100.00\\
2.00	50.00\\
3.00	66.67\\
4.00	75.00\\
5.00	80.00\\
6.00	83.33\\
7.00	85.71\\
8.00	87.50\\
9.00	88.89\\
10.00	80.00\\
11.00	81.82\\
12.00	75.00\\
13.00	76.92\\
14.00	78.57\\
15.00	73.33\\
16.00	68.75\\
17.00	70.59\\
18.00	66.67\\
19.00	68.42\\
21.00	61.90\\
22.00	63.64\\
23.00	60.87\\
25.00	64.00\\
26.00	65.38\\
28.00	67.86\\
30.00	66.67\\
32.00	65.62\\
34.00	67.65\\
36.00	63.89\\
38.00	63.16\\
41.00	63.41\\
43.00	65.12\\
46.00	65.22\\
49.00	65.31\\
52.00	65.38\\
56.00	66.07\\
59.00	64.41\\
63.00	65.08\\
67.00	64.18\\
71.00	63.38\\
76.00	63.16\\
81.00	61.73\\
86.00	62.79\\
91.00	63.74\\
97.00	63.92\\
103.00	62.14\\
110.00	61.82\\
117.00	61.54\\
124.00	62.10\\
132.00	63.64\\
140.00	65.00\\
149.00	64.43\\
159.00	64.78\\
169.00	63.31\\
180.00	63.33\\
191.00	62.30\\
204.00	62.25\\
217.00	62.67\\
230.00	63.48\\
245.00	64.08\\
261.00	63.98\\
277.00	64.62\\
295.00	64.75\\
314.00	64.65\\
334.00	65.27\\
355.00	65.35\\
378.00	65.08\\
402.00	64.18\\
427.00	64.17\\
455.00	64.40\\
484.00	64.46\\
515.00	63.69\\
547.00	63.25\\
582.00	63.23\\
619.00	63.17\\
659.00	63.13\\
701.00	63.20\\
746.00	61.80\\
793.00	62.30\\
844.00	61.97\\
897.00	61.76\\
955.00	61.99\\
1016.00	61.81\\
1080.00	61.85\\
1149.00	61.53\\
1222.00	61.87\\
1300.00	61.85\\
1383.00	62.11\\
1472.00	61.55\\
1565.00	61.85\\
1665.00	62.04\\
1771.00	61.89\\
1884.00	61.73\\
2005.00	62.00\\
2132.00	62.29\\
2268.00	62.39\\
2413.00	62.58\\
2567.00	62.37\\
2730.00	62.38\\
2905.00	62.31\\
3090.00	62.39\\
3287.00	62.49\\
3496.00	62.24\\
3719.00	62.41\\
3957.00	62.14\\
4209.00	62.01\\
4477.00	62.41\\
4763.00	62.08\\
5066.00	61.98\\
5389.00	62.18\\
5733.00	61.94\\
6099.00	61.65\\
6488.00	61.68\\
6901.00	61.72\\
7341.00	61.98\\
7809.00	61.92\\
8307.00	62.02\\
8837.00	62.13\\
9401.00	61.99\\
10000.00	61.84\\
};
\addplot [color=white!45!black, line width=1.2pt, forget plot]
  table[row sep=crcr]{%
1.00	0.00\\
2.00	50.00\\
3.00	33.33\\
4.00	50.00\\
5.00	40.00\\
6.00	50.00\\
7.00	57.14\\
8.00	62.50\\
9.00	66.67\\
10.00	70.00\\
11.00	63.64\\
12.00	66.67\\
13.00	69.23\\
14.00	64.29\\
15.00	66.67\\
16.00	62.50\\
17.00	58.82\\
18.00	61.11\\
19.00	57.89\\
21.00	57.14\\
22.00	54.55\\
23.00	52.17\\
25.00	52.00\\
26.00	53.85\\
28.00	57.14\\
30.00	56.67\\
32.00	59.38\\
34.00	61.76\\
36.00	61.11\\
38.00	60.53\\
41.00	58.54\\
43.00	58.14\\
46.00	58.70\\
49.00	57.14\\
52.00	57.69\\
56.00	58.93\\
59.00	59.32\\
63.00	61.90\\
67.00	61.19\\
71.00	63.38\\
76.00	63.16\\
81.00	60.49\\
86.00	60.47\\
91.00	59.34\\
97.00	59.79\\
103.00	59.22\\
110.00	57.27\\
117.00	56.41\\
124.00	56.45\\
132.00	59.09\\
140.00	60.00\\
149.00	60.40\\
159.00	61.64\\
169.00	60.36\\
180.00	62.22\\
191.00	63.35\\
204.00	64.71\\
217.00	64.98\\
230.00	64.78\\
245.00	65.31\\
261.00	66.28\\
277.00	66.06\\
295.00	67.80\\
314.00	68.15\\
334.00	68.26\\
355.00	68.73\\
378.00	68.52\\
402.00	69.65\\
427.00	70.26\\
455.00	70.11\\
484.00	70.04\\
515.00	70.49\\
547.00	70.93\\
582.00	70.79\\
619.00	70.92\\
659.00	70.56\\
701.00	70.47\\
746.00	70.91\\
793.00	70.49\\
844.00	70.85\\
897.00	70.68\\
955.00	70.47\\
1016.00	70.67\\
1080.00	70.83\\
1149.00	71.37\\
1222.00	71.28\\
1300.00	71.08\\
1383.00	70.72\\
1472.00	71.06\\
1565.00	70.73\\
1665.00	70.75\\
1771.00	70.86\\
1884.00	70.70\\
2005.00	70.42\\
2132.00	70.59\\
2268.00	70.72\\
2413.00	70.70\\
2567.00	70.55\\
2730.00	70.51\\
2905.00	70.43\\
3090.00	70.26\\
3287.00	69.88\\
3496.00	69.59\\
3719.00	69.56\\
3957.00	69.55\\
4209.00	69.45\\
4477.00	69.15\\
4763.00	69.22\\
5066.00	69.09\\
5389.00	69.01\\
5733.00	69.16\\
6099.00	69.32\\
6488.00	69.39\\
6901.00	69.37\\
7341.00	69.19\\
7809.00	69.16\\
8307.00	69.10\\
8837.00	69.18\\
9401.00	69.14\\
10000.00	69.26\\
};
\addplot [color=gray!85!lightgray, dashed, line width=1.2pt, forget plot]
  table[row sep=crcr]{%
1.00	0.00\\
2.00	50.00\\
3.00	66.67\\
4.00	50.00\\
5.00	40.00\\
6.00	33.33\\
7.00	42.86\\
8.00	50.00\\
9.00	55.56\\
10.00	60.00\\
11.00	54.55\\
12.00	58.33\\
13.00	61.54\\
14.00	64.29\\
15.00	66.67\\
16.00	62.50\\
17.00	58.82\\
18.00	61.11\\
19.00	63.16\\
21.00	61.90\\
22.00	63.64\\
23.00	60.87\\
25.00	60.00\\
26.00	61.54\\
28.00	64.29\\
30.00	66.67\\
32.00	62.50\\
34.00	61.76\\
36.00	63.89\\
38.00	65.79\\
41.00	63.41\\
43.00	62.79\\
46.00	63.04\\
49.00	59.18\\
52.00	57.69\\
56.00	58.93\\
59.00	59.32\\
63.00	57.14\\
67.00	53.73\\
71.00	53.52\\
76.00	55.26\\
81.00	58.02\\
86.00	59.30\\
91.00	60.44\\
97.00	60.82\\
103.00	60.19\\
110.00	60.91\\
117.00	60.68\\
124.00	61.29\\
132.00	61.36\\
140.00	62.86\\
149.00	63.76\\
159.00	64.78\\
169.00	65.09\\
180.00	65.00\\
191.00	64.92\\
204.00	64.22\\
217.00	64.06\\
230.00	63.48\\
245.00	64.08\\
261.00	64.37\\
277.00	64.26\\
295.00	64.07\\
314.00	64.65\\
334.00	64.97\\
355.00	65.07\\
378.00	65.61\\
402.00	65.67\\
427.00	66.28\\
455.00	66.81\\
484.00	66.12\\
515.00	65.83\\
547.00	65.81\\
582.00	65.81\\
619.00	65.43\\
659.00	64.49\\
701.00	64.76\\
746.00	64.88\\
793.00	64.69\\
844.00	64.10\\
897.00	63.55\\
955.00	62.72\\
1016.00	63.78\\
1080.00	64.35\\
1149.00	64.75\\
1222.00	64.57\\
1300.00	64.38\\
1383.00	64.28\\
1472.00	64.88\\
1565.00	65.24\\
1665.00	65.47\\
1771.00	65.56\\
1884.00	65.50\\
2005.00	65.84\\
2132.00	65.81\\
2268.00	65.65\\
2413.00	65.52\\
2567.00	65.56\\
2730.00	65.42\\
2905.00	65.51\\
3090.00	65.40\\
3287.00	65.38\\
3496.00	65.59\\
3719.00	65.93\\
3957.00	66.09\\
4209.00	66.29\\
4477.00	66.41\\
4763.00	66.45\\
5066.00	66.56\\
5389.00	66.38\\
5733.00	66.28\\
6099.00	65.86\\
6488.00	66.15\\
6901.00	66.08\\
7341.00	66.09\\
7809.00	66.04\\
8307.00	66.08\\
8837.00	66.05\\
9401.00	66.16\\
10000.00	66.21\\
};
\addplot [color=white!50!darkgray, dotted, line width=1.2pt, forget plot]
  table[row sep=crcr]{%
1.00	100.00\\
2.00	50.00\\
3.00	66.67\\
4.00	75.00\\
5.00	60.00\\
6.00	66.67\\
7.00	71.43\\
8.00	62.50\\
9.00	66.67\\
10.00	70.00\\
11.00	72.73\\
12.00	66.67\\
13.00	61.54\\
14.00	64.29\\
15.00	60.00\\
16.00	62.50\\
17.00	64.71\\
18.00	66.67\\
19.00	63.16\\
21.00	61.90\\
22.00	63.64\\
23.00	65.22\\
25.00	64.00\\
26.00	65.38\\
28.00	64.29\\
30.00	63.33\\
32.00	62.50\\
34.00	61.76\\
36.00	61.11\\
38.00	60.53\\
41.00	63.41\\
43.00	65.12\\
46.00	65.22\\
49.00	67.35\\
52.00	69.23\\
56.00	69.64\\
59.00	67.80\\
63.00	65.08\\
67.00	65.67\\
71.00	64.79\\
76.00	64.47\\
81.00	64.20\\
86.00	62.79\\
91.00	62.64\\
97.00	61.86\\
103.00	62.14\\
110.00	60.00\\
117.00	61.54\\
124.00	62.90\\
132.00	62.88\\
140.00	62.14\\
149.00	62.42\\
159.00	63.52\\
169.00	62.72\\
180.00	63.33\\
191.00	63.35\\
204.00	64.22\\
217.00	63.59\\
230.00	63.91\\
245.00	65.31\\
261.00	65.90\\
277.00	65.34\\
295.00	66.10\\
314.00	66.88\\
334.00	67.66\\
355.00	66.76\\
378.00	66.40\\
402.00	66.42\\
427.00	66.04\\
455.00	65.49\\
484.00	65.50\\
515.00	64.08\\
547.00	64.72\\
582.00	65.29\\
619.00	64.94\\
659.00	65.10\\
701.00	64.91\\
746.00	65.28\\
793.00	64.69\\
844.00	64.81\\
897.00	64.77\\
955.00	65.34\\
1016.00	65.26\\
1080.00	64.91\\
1149.00	65.10\\
1222.00	65.38\\
1300.00	65.46\\
1383.00	65.80\\
1472.00	65.90\\
1565.00	66.01\\
1665.00	65.65\\
1771.00	66.01\\
1884.00	66.24\\
2005.00	66.63\\
2132.00	66.74\\
2268.00	66.98\\
2413.00	67.22\\
2567.00	66.97\\
2730.00	67.14\\
2905.00	67.16\\
3090.00	67.48\\
3287.00	67.54\\
3496.00	67.79\\
3719.00	67.87\\
3957.00	67.68\\
4209.00	67.71\\
4477.00	67.84\\
4763.00	68.05\\
5066.00	67.88\\
5389.00	67.88\\
5733.00	68.06\\
6099.00	67.96\\
6488.00	67.91\\
6901.00	67.67\\
7341.00	67.59\\
7809.00	67.60\\
8307.00	67.68\\
8837.00	67.59\\
9401.00	67.66\\
10000.00	67.65\\
};
\addplot [color=black!5!lightgray, dashdotted, line width=1.2pt, forget plot]
  table[row sep=crcr]{%
1.00	100.00\\
2.00	100.00\\
3.00	66.67\\
4.00	75.00\\
5.00	60.00\\
6.00	66.67\\
7.00	71.43\\
8.00	75.00\\
9.00	77.78\\
10.00	70.00\\
11.00	72.73\\
12.00	75.00\\
13.00	76.92\\
14.00	78.57\\
15.00	80.00\\
16.00	81.25\\
17.00	82.35\\
18.00	83.33\\
19.00	84.21\\
21.00	85.71\\
22.00	86.36\\
23.00	86.96\\
25.00	88.00\\
26.00	84.62\\
28.00	82.14\\
30.00	80.00\\
32.00	81.25\\
34.00	76.47\\
36.00	77.78\\
38.00	76.32\\
41.00	78.05\\
43.00	76.74\\
46.00	73.91\\
49.00	69.39\\
52.00	71.15\\
56.00	69.64\\
59.00	67.80\\
63.00	66.67\\
67.00	67.16\\
71.00	66.20\\
76.00	65.79\\
81.00	67.90\\
86.00	68.60\\
91.00	67.03\\
97.00	65.98\\
103.00	65.05\\
110.00	62.73\\
117.00	62.39\\
124.00	62.90\\
132.00	61.36\\
140.00	61.43\\
149.00	61.74\\
159.00	61.64\\
169.00	62.72\\
180.00	62.22\\
191.00	60.21\\
204.00	60.29\\
217.00	59.91\\
230.00	59.57\\
245.00	60.41\\
261.00	61.69\\
277.00	60.29\\
295.00	60.34\\
314.00	59.55\\
334.00	59.88\\
355.00	60.00\\
378.00	60.05\\
402.00	59.95\\
427.00	59.95\\
455.00	60.66\\
484.00	61.16\\
515.00	60.19\\
547.00	60.88\\
582.00	61.51\\
619.00	61.55\\
659.00	60.70\\
701.00	60.91\\
746.00	60.99\\
793.00	61.16\\
844.00	61.14\\
897.00	61.09\\
955.00	61.36\\
1016.00	60.73\\
1080.00	60.19\\
1149.00	60.14\\
1222.00	60.07\\
1300.00	60.08\\
1383.00	60.16\\
1472.00	60.19\\
1565.00	59.68\\
1665.00	59.76\\
1771.00	59.85\\
1884.00	60.19\\
2005.00	59.85\\
2132.00	59.52\\
2268.00	59.35\\
2413.00	59.35\\
2567.00	59.17\\
2730.00	59.12\\
2905.00	59.00\\
3090.00	59.81\\
3287.00	59.45\\
3496.00	59.15\\
3719.00	59.37\\
3957.00	59.29\\
4209.00	59.25\\
4477.00	59.26\\
4763.00	59.08\\
5066.00	59.16\\
5389.00	58.84\\
5733.00	58.80\\
6099.00	58.47\\
6488.00	58.42\\
6901.00	58.57\\
7341.00	58.66\\
7809.00	58.56\\
8307.00	58.63\\
8837.00	58.55\\
9401.00	58.62\\
10000.00	58.46\\
};
\addplot [color=white!80!black, line width=1.2pt, forget plot]
  table[row sep=crcr]{%
1.00	100.00\\
2.00	100.00\\
3.00	100.00\\
4.00	100.00\\
5.00	100.00\\
6.00	100.00\\
7.00	85.71\\
8.00	75.00\\
9.00	77.78\\
10.00	80.00\\
11.00	81.82\\
12.00	83.33\\
13.00	84.62\\
14.00	85.71\\
15.00	80.00\\
16.00	75.00\\
17.00	70.59\\
18.00	66.67\\
19.00	68.42\\
21.00	66.67\\
22.00	68.18\\
23.00	65.22\\
25.00	64.00\\
26.00	61.54\\
28.00	60.71\\
30.00	63.33\\
32.00	65.62\\
34.00	61.76\\
36.00	58.33\\
38.00	57.89\\
41.00	58.54\\
43.00	60.47\\
46.00	58.70\\
49.00	61.22\\
52.00	63.46\\
56.00	64.29\\
59.00	62.71\\
63.00	63.49\\
67.00	62.69\\
71.00	63.38\\
76.00	63.16\\
81.00	62.96\\
86.00	62.79\\
91.00	63.74\\
97.00	63.92\\
103.00	64.08\\
110.00	66.36\\
117.00	67.52\\
124.00	67.74\\
132.00	68.18\\
140.00	69.29\\
149.00	70.47\\
159.00	71.07\\
169.00	70.41\\
180.00	68.89\\
191.00	69.63\\
204.00	69.12\\
217.00	69.59\\
230.00	70.00\\
245.00	71.43\\
261.00	71.26\\
277.00	71.12\\
295.00	69.83\\
314.00	70.70\\
334.00	70.36\\
355.00	69.30\\
378.00	69.31\\
402.00	69.15\\
427.00	70.02\\
455.00	69.67\\
484.00	68.39\\
515.00	68.35\\
547.00	68.01\\
582.00	67.01\\
619.00	66.56\\
659.00	66.62\\
701.00	66.48\\
746.00	66.89\\
793.00	66.33\\
844.00	66.71\\
897.00	67.00\\
955.00	66.91\\
1016.00	66.83\\
1080.00	67.04\\
1149.00	67.28\\
1222.00	67.02\\
1300.00	66.92\\
1383.00	66.67\\
1472.00	66.85\\
1565.00	67.03\\
1665.00	66.85\\
1771.00	66.80\\
1884.00	66.61\\
2005.00	66.88\\
2132.00	67.12\\
2268.00	67.50\\
2413.00	67.51\\
2567.00	67.47\\
2730.00	67.62\\
2905.00	67.81\\
3090.00	67.73\\
3287.00	67.48\\
3496.00	67.51\\
3719.00	67.41\\
3957.00	67.42\\
4209.00	67.38\\
4477.00	67.37\\
4763.00	67.33\\
5066.00	67.45\\
5389.00	67.67\\
5733.00	67.82\\
6099.00	67.63\\
6488.00	67.80\\
6901.00	67.69\\
7341.00	67.88\\
7809.00	67.60\\
8307.00	67.71\\
8837.00	67.49\\
9401.00	67.46\\
10000.00	67.60\\
};
\end{axis}

\begin{axis}[%
width=0.391\figurewidth,
height=0.265\figureheight,
at={(0.54\figurewidth,0\figureheight)},
scale only axis,
xmin=1.00,
xmax=10000.00,
xlabel style={font=\color{mycolor1}},
xlabel={Iterationen $\sym{n}_{\sym{idx_MC}}$},
ymin=0.00,
ymax=105.00,
ytick={0.00,20.00,40.00,60.00,80.00,100.00},
yticklabels={\empty},
axis background/.style={fill=white},
title style={font=\bfseries\color{mycolor1}},
title={\textbf{$\sym{beta}_{22}$}},
xmajorgrids,
ymajorgrids,
grid style={dotted, opacity=0.25}
]
\addplot [color=white!10!black, line width=1.2pt, forget plot]
  table[row sep=crcr]{%
1.00	100.00\\
2.00	100.00\\
3.00	66.67\\
4.00	50.00\\
5.00	60.00\\
6.00	66.67\\
7.00	71.43\\
8.00	75.00\\
9.00	66.67\\
10.00	60.00\\
11.00	63.64\\
12.00	66.67\\
13.00	61.54\\
14.00	64.29\\
15.00	66.67\\
16.00	68.75\\
17.00	70.59\\
18.00	72.22\\
19.00	73.68\\
21.00	71.43\\
22.00	72.73\\
23.00	73.91\\
25.00	76.00\\
26.00	76.92\\
28.00	78.57\\
30.00	73.33\\
32.00	75.00\\
34.00	76.47\\
36.00	75.00\\
38.00	73.68\\
41.00	73.17\\
43.00	72.09\\
46.00	73.91\\
49.00	73.47\\
52.00	71.15\\
56.00	69.64\\
59.00	69.49\\
63.00	66.67\\
67.00	68.66\\
71.00	67.61\\
76.00	67.11\\
81.00	66.67\\
86.00	66.28\\
91.00	67.03\\
97.00	64.95\\
103.00	62.14\\
110.00	61.82\\
117.00	62.39\\
124.00	62.10\\
132.00	62.88\\
140.00	62.86\\
149.00	63.76\\
159.00	64.78\\
169.00	65.09\\
180.00	64.44\\
191.00	64.40\\
204.00	63.73\\
217.00	63.13\\
230.00	63.04\\
245.00	64.08\\
261.00	64.75\\
277.00	66.06\\
295.00	66.44\\
314.00	67.52\\
334.00	68.56\\
355.00	68.73\\
378.00	68.52\\
402.00	68.91\\
427.00	68.85\\
455.00	68.57\\
484.00	69.21\\
515.00	68.93\\
547.00	69.10\\
582.00	69.42\\
619.00	69.79\\
659.00	70.26\\
701.00	69.76\\
746.00	69.71\\
793.00	69.74\\
844.00	70.14\\
897.00	70.57\\
955.00	70.99\\
1016.00	71.06\\
1080.00	71.02\\
1149.00	71.11\\
1222.00	71.44\\
1300.00	71.85\\
1383.00	71.66\\
1472.00	70.92\\
1565.00	70.29\\
1665.00	69.91\\
1771.00	69.96\\
1884.00	69.80\\
2005.00	69.93\\
2132.00	70.08\\
2268.00	70.06\\
2413.00	69.83\\
2567.00	69.42\\
2730.00	69.67\\
2905.00	69.78\\
3090.00	69.84\\
3287.00	69.46\\
3496.00	69.36\\
3719.00	69.16\\
3957.00	69.12\\
4209.00	69.02\\
4477.00	69.09\\
4763.00	69.16\\
5066.00	68.75\\
5389.00	68.94\\
5733.00	69.13\\
6099.00	69.22\\
6488.00	69.30\\
6901.00	69.42\\
7341.00	69.36\\
7809.00	69.51\\
8307.00	69.57\\
8837.00	69.53\\
9401.00	69.65\\
10000.00	69.71\\
};
\addplot [color=black!25!darkgray, dashed, line width=1.2pt, forget plot]
  table[row sep=crcr]{%
1.00	100.00\\
2.00	100.00\\
3.00	100.00\\
4.00	75.00\\
5.00	80.00\\
6.00	66.67\\
7.00	71.43\\
8.00	75.00\\
9.00	77.78\\
10.00	80.00\\
11.00	81.82\\
12.00	83.33\\
13.00	76.92\\
14.00	78.57\\
15.00	80.00\\
16.00	81.25\\
17.00	76.47\\
18.00	72.22\\
19.00	73.68\\
21.00	76.19\\
22.00	72.73\\
23.00	73.91\\
25.00	72.00\\
26.00	73.08\\
28.00	67.86\\
30.00	66.67\\
32.00	65.62\\
34.00	64.71\\
36.00	63.89\\
38.00	65.79\\
41.00	63.41\\
43.00	65.12\\
46.00	65.22\\
49.00	67.35\\
52.00	69.23\\
56.00	71.43\\
59.00	71.19\\
63.00	69.84\\
67.00	70.15\\
71.00	69.01\\
76.00	68.42\\
81.00	69.14\\
86.00	67.44\\
91.00	68.13\\
97.00	67.01\\
103.00	64.08\\
110.00	64.55\\
117.00	64.96\\
124.00	62.90\\
132.00	63.64\\
140.00	64.29\\
149.00	63.76\\
159.00	63.52\\
169.00	63.31\\
180.00	65.00\\
191.00	64.40\\
204.00	65.69\\
217.00	64.98\\
230.00	65.22\\
245.00	65.31\\
261.00	64.75\\
277.00	65.70\\
295.00	67.46\\
314.00	66.88\\
334.00	67.66\\
355.00	67.89\\
378.00	68.25\\
402.00	68.16\\
427.00	69.09\\
455.00	68.79\\
484.00	68.80\\
515.00	69.90\\
547.00	70.20\\
582.00	69.59\\
619.00	69.95\\
659.00	70.11\\
701.00	70.76\\
746.00	70.91\\
793.00	71.12\\
844.00	71.45\\
897.00	71.24\\
955.00	71.31\\
1016.00	71.06\\
1080.00	71.30\\
1149.00	70.58\\
1222.00	70.70\\
1300.00	70.54\\
1383.00	70.50\\
1472.00	70.72\\
1565.00	70.35\\
1665.00	70.27\\
1771.00	70.30\\
1884.00	70.59\\
2005.00	70.52\\
2132.00	70.45\\
2268.00	70.72\\
2413.00	70.62\\
2567.00	70.71\\
2730.00	70.15\\
2905.00	70.12\\
3090.00	69.81\\
3287.00	69.64\\
3496.00	69.79\\
3719.00	69.56\\
3957.00	69.47\\
4209.00	69.54\\
4477.00	69.53\\
4763.00	69.62\\
5066.00	69.62\\
5389.00	69.66\\
5733.00	69.56\\
6099.00	69.49\\
6488.00	69.47\\
6901.00	69.61\\
7341.00	69.49\\
7809.00	69.48\\
8307.00	69.21\\
8837.00	69.16\\
9401.00	69.08\\
10000.00	68.96\\
};
\addplot [color=black!45!gray, dotted, line width=1.2pt, forget plot]
  table[row sep=crcr]{%
1.00	100.00\\
2.00	50.00\\
3.00	33.33\\
4.00	50.00\\
5.00	60.00\\
6.00	66.67\\
7.00	71.43\\
8.00	75.00\\
9.00	66.67\\
10.00	70.00\\
11.00	72.73\\
12.00	75.00\\
13.00	76.92\\
14.00	78.57\\
15.00	80.00\\
16.00	81.25\\
17.00	82.35\\
18.00	83.33\\
19.00	84.21\\
21.00	85.71\\
22.00	86.36\\
23.00	86.96\\
25.00	88.00\\
26.00	88.46\\
28.00	89.29\\
30.00	90.00\\
32.00	87.50\\
34.00	85.29\\
36.00	80.56\\
38.00	81.58\\
41.00	82.93\\
43.00	81.40\\
46.00	80.43\\
49.00	81.63\\
52.00	78.85\\
56.00	78.57\\
59.00	77.97\\
63.00	76.19\\
67.00	77.61\\
71.00	76.06\\
76.00	76.32\\
81.00	74.07\\
86.00	74.42\\
91.00	75.82\\
97.00	76.29\\
103.00	76.70\\
110.00	76.36\\
117.00	76.07\\
124.00	76.61\\
132.00	77.27\\
140.00	77.86\\
149.00	75.84\\
159.00	74.21\\
169.00	73.96\\
180.00	74.44\\
191.00	73.30\\
204.00	74.02\\
217.00	73.73\\
230.00	72.61\\
245.00	72.24\\
261.00	71.26\\
277.00	70.76\\
295.00	72.54\\
314.00	72.29\\
334.00	71.26\\
355.00	71.55\\
378.00	71.69\\
402.00	72.14\\
427.00	72.60\\
455.00	72.75\\
484.00	72.52\\
515.00	72.62\\
547.00	72.21\\
582.00	72.16\\
619.00	71.08\\
659.00	70.71\\
701.00	71.04\\
746.00	70.91\\
793.00	70.24\\
844.00	70.14\\
897.00	70.23\\
955.00	69.53\\
1016.00	69.69\\
1080.00	70.00\\
1149.00	70.41\\
1222.00	70.05\\
1300.00	69.85\\
1383.00	69.49\\
1472.00	69.50\\
1565.00	69.14\\
1665.00	68.83\\
1771.00	68.94\\
1884.00	68.90\\
2005.00	68.38\\
2132.00	67.82\\
2268.00	67.68\\
2413.00	67.55\\
2567.00	67.71\\
2730.00	67.84\\
2905.00	67.68\\
3090.00	67.38\\
3287.00	67.14\\
3496.00	67.25\\
3719.00	67.20\\
3957.00	67.07\\
4209.00	67.09\\
4477.00	67.23\\
4763.00	67.56\\
5066.00	67.61\\
5389.00	67.80\\
5733.00	67.84\\
6099.00	67.91\\
6488.00	68.11\\
6901.00	68.35\\
7341.00	68.42\\
7809.00	68.51\\
8307.00	68.46\\
8837.00	68.54\\
9401.00	68.46\\
10000.00	68.37\\
};
\addplot [color=white!15!darkgray, dashdotted, line width=1.2pt, forget plot]
  table[row sep=crcr]{%
1.00	100.00\\
2.00	100.00\\
3.00	66.67\\
4.00	50.00\\
5.00	40.00\\
6.00	50.00\\
7.00	42.86\\
8.00	50.00\\
9.00	55.56\\
10.00	60.00\\
11.00	63.64\\
12.00	66.67\\
13.00	69.23\\
14.00	71.43\\
15.00	73.33\\
16.00	75.00\\
17.00	76.47\\
18.00	72.22\\
19.00	73.68\\
21.00	71.43\\
22.00	72.73\\
23.00	73.91\\
25.00	76.00\\
26.00	76.92\\
28.00	75.00\\
30.00	70.00\\
32.00	68.75\\
34.00	67.65\\
36.00	69.44\\
38.00	71.05\\
41.00	68.29\\
43.00	69.77\\
46.00	67.39\\
49.00	69.39\\
52.00	69.23\\
56.00	69.64\\
59.00	67.80\\
63.00	68.25\\
67.00	70.15\\
71.00	69.01\\
76.00	71.05\\
81.00	70.37\\
86.00	72.09\\
91.00	69.23\\
97.00	68.04\\
103.00	67.96\\
110.00	66.36\\
117.00	64.10\\
124.00	66.13\\
132.00	67.42\\
140.00	69.29\\
149.00	69.80\\
159.00	67.92\\
169.00	68.64\\
180.00	68.33\\
191.00	69.11\\
204.00	68.63\\
217.00	68.66\\
230.00	68.26\\
245.00	68.57\\
261.00	67.82\\
277.00	67.87\\
295.00	68.14\\
314.00	68.79\\
334.00	69.46\\
355.00	69.58\\
378.00	70.63\\
402.00	71.14\\
427.00	70.49\\
455.00	70.99\\
484.00	70.66\\
515.00	70.87\\
547.00	69.84\\
582.00	70.27\\
619.00	70.44\\
659.00	70.71\\
701.00	70.61\\
746.00	70.38\\
793.00	70.49\\
844.00	70.38\\
897.00	70.01\\
955.00	70.47\\
1016.00	70.37\\
1080.00	70.19\\
1149.00	70.06\\
1222.00	69.97\\
1300.00	69.46\\
1383.00	69.13\\
1472.00	68.82\\
1565.00	69.14\\
1665.00	69.55\\
1771.00	69.28\\
1884.00	68.74\\
2005.00	68.78\\
2132.00	69.04\\
2268.00	69.27\\
2413.00	69.33\\
2567.00	69.22\\
2730.00	69.23\\
2905.00	69.23\\
3090.00	69.26\\
3287.00	69.24\\
3496.00	69.19\\
3719.00	69.32\\
3957.00	69.24\\
4209.00	69.21\\
4477.00	69.15\\
4763.00	69.24\\
5066.00	69.11\\
5389.00	69.20\\
5733.00	68.90\\
6099.00	68.86\\
6488.00	68.85\\
6901.00	68.79\\
7341.00	68.93\\
7809.00	69.02\\
8307.00	68.75\\
8837.00	68.91\\
9401.00	69.04\\
10000.00	69.03\\
};
\addplot [color=white!45!black, line width=1.2pt, forget plot]
  table[row sep=crcr]{%
1.00	100.00\\
2.00	50.00\\
3.00	33.33\\
4.00	25.00\\
5.00	20.00\\
6.00	33.33\\
7.00	42.86\\
8.00	50.00\\
9.00	44.44\\
10.00	50.00\\
11.00	45.45\\
12.00	50.00\\
13.00	53.85\\
14.00	50.00\\
15.00	53.33\\
16.00	56.25\\
17.00	52.94\\
18.00	50.00\\
19.00	52.63\\
21.00	57.14\\
22.00	54.55\\
23.00	52.17\\
25.00	56.00\\
26.00	57.69\\
28.00	57.14\\
30.00	60.00\\
32.00	62.50\\
34.00	61.76\\
36.00	63.89\\
38.00	65.79\\
41.00	65.85\\
43.00	62.79\\
46.00	65.22\\
49.00	65.31\\
52.00	65.38\\
56.00	64.29\\
59.00	62.71\\
63.00	63.49\\
67.00	65.67\\
71.00	67.61\\
76.00	67.11\\
81.00	66.67\\
86.00	66.28\\
91.00	65.93\\
97.00	68.04\\
103.00	67.96\\
110.00	69.09\\
117.00	69.23\\
124.00	68.55\\
132.00	68.94\\
140.00	69.29\\
149.00	69.80\\
159.00	68.55\\
169.00	66.86\\
180.00	67.78\\
191.00	67.54\\
204.00	66.67\\
217.00	66.82\\
230.00	66.52\\
245.00	66.94\\
261.00	67.05\\
277.00	66.06\\
295.00	66.44\\
314.00	67.20\\
334.00	66.17\\
355.00	66.20\\
378.00	66.40\\
402.00	66.67\\
427.00	66.51\\
455.00	66.59\\
484.00	66.74\\
515.00	66.41\\
547.00	66.36\\
582.00	66.15\\
619.00	66.07\\
659.00	66.16\\
701.00	66.76\\
746.00	66.62\\
793.00	66.96\\
844.00	66.59\\
897.00	66.33\\
955.00	66.81\\
1016.00	67.32\\
1080.00	67.04\\
1149.00	67.36\\
1222.00	67.51\\
1300.00	67.77\\
1383.00	67.53\\
1472.00	67.73\\
1565.00	67.67\\
1665.00	67.81\\
1771.00	68.27\\
1884.00	67.89\\
2005.00	67.93\\
2132.00	68.06\\
2268.00	68.47\\
2413.00	68.38\\
2567.00	68.52\\
2730.00	68.13\\
2905.00	68.26\\
3090.00	68.32\\
3287.00	68.42\\
3496.00	68.59\\
3719.00	68.78\\
3957.00	68.61\\
4209.00	68.71\\
4477.00	68.95\\
4763.00	68.86\\
5066.00	68.93\\
5389.00	68.79\\
5733.00	68.67\\
6099.00	68.80\\
6488.00	68.73\\
6901.00	68.63\\
7341.00	68.72\\
7809.00	68.81\\
8307.00	68.77\\
8837.00	68.61\\
9401.00	68.60\\
10000.00	68.69\\
};
\addplot [color=gray!85!lightgray, dashed, line width=1.2pt, forget plot]
  table[row sep=crcr]{%
1.00	100.00\\
2.00	100.00\\
3.00	100.00\\
4.00	100.00\\
5.00	100.00\\
6.00	83.33\\
7.00	71.43\\
8.00	75.00\\
9.00	77.78\\
10.00	70.00\\
11.00	63.64\\
12.00	58.33\\
13.00	61.54\\
14.00	64.29\\
15.00	60.00\\
16.00	62.50\\
17.00	64.71\\
18.00	61.11\\
19.00	63.16\\
21.00	61.90\\
22.00	63.64\\
23.00	65.22\\
25.00	68.00\\
26.00	65.38\\
28.00	64.29\\
30.00	63.33\\
32.00	65.62\\
34.00	64.71\\
36.00	66.67\\
38.00	65.79\\
41.00	68.29\\
43.00	67.44\\
46.00	67.39\\
49.00	65.31\\
52.00	63.46\\
56.00	62.50\\
59.00	62.71\\
63.00	65.08\\
67.00	65.67\\
71.00	66.20\\
76.00	65.79\\
81.00	66.67\\
86.00	63.95\\
91.00	64.84\\
97.00	64.95\\
103.00	66.02\\
110.00	67.27\\
117.00	65.81\\
124.00	65.32\\
132.00	64.39\\
140.00	63.57\\
149.00	63.76\\
159.00	65.41\\
169.00	64.50\\
180.00	66.11\\
191.00	67.02\\
204.00	66.67\\
217.00	67.28\\
230.00	66.52\\
245.00	66.94\\
261.00	66.67\\
277.00	66.06\\
295.00	65.42\\
314.00	64.65\\
334.00	65.57\\
355.00	65.92\\
378.00	65.34\\
402.00	66.67\\
427.00	66.04\\
455.00	65.71\\
484.00	67.15\\
515.00	67.77\\
547.00	67.28\\
582.00	66.84\\
619.00	67.21\\
659.00	67.07\\
701.00	67.33\\
746.00	66.89\\
793.00	67.59\\
844.00	67.42\\
897.00	67.56\\
955.00	67.54\\
1016.00	67.13\\
1080.00	67.13\\
1149.00	67.19\\
1222.00	67.59\\
1300.00	67.62\\
1383.00	67.39\\
1472.00	67.73\\
1565.00	68.63\\
1665.00	68.11\\
1771.00	67.98\\
1884.00	67.68\\
2005.00	67.88\\
2132.00	68.01\\
2268.00	68.21\\
2413.00	68.13\\
2567.00	67.90\\
2730.00	68.28\\
2905.00	68.12\\
3090.00	68.06\\
3287.00	67.93\\
3496.00	67.99\\
3719.00	67.81\\
3957.00	67.93\\
4209.00	67.85\\
4477.00	67.43\\
4763.00	67.54\\
5066.00	67.67\\
5389.00	67.47\\
5733.00	67.22\\
6099.00	67.55\\
6488.00	67.52\\
6901.00	67.60\\
7341.00	67.69\\
7809.00	67.67\\
8307.00	67.69\\
8837.00	67.74\\
9401.00	67.80\\
10000.00	67.88\\
};
\addplot [color=white!50!darkgray, dotted, line width=1.2pt, forget plot]
  table[row sep=crcr]{%
1.00	100.00\\
2.00	100.00\\
3.00	66.67\\
4.00	75.00\\
5.00	80.00\\
6.00	83.33\\
7.00	85.71\\
8.00	87.50\\
9.00	88.89\\
10.00	90.00\\
11.00	90.91\\
12.00	91.67\\
13.00	92.31\\
14.00	92.86\\
15.00	93.33\\
16.00	93.75\\
17.00	94.12\\
18.00	88.89\\
19.00	89.47\\
21.00	85.71\\
22.00	86.36\\
23.00	86.96\\
25.00	88.00\\
26.00	88.46\\
28.00	89.29\\
30.00	90.00\\
32.00	87.50\\
34.00	85.29\\
36.00	86.11\\
38.00	81.58\\
41.00	75.61\\
43.00	74.42\\
46.00	76.09\\
49.00	73.47\\
52.00	73.08\\
56.00	73.21\\
59.00	71.19\\
63.00	73.02\\
67.00	74.63\\
71.00	73.24\\
76.00	71.05\\
81.00	71.60\\
86.00	72.09\\
91.00	72.53\\
97.00	72.16\\
103.00	69.90\\
110.00	68.18\\
117.00	68.38\\
124.00	69.35\\
132.00	70.45\\
140.00	70.71\\
149.00	70.47\\
159.00	70.44\\
169.00	71.01\\
180.00	72.22\\
191.00	72.25\\
204.00	71.57\\
217.00	71.43\\
230.00	72.61\\
245.00	72.65\\
261.00	72.03\\
277.00	71.48\\
295.00	71.86\\
314.00	71.66\\
334.00	71.86\\
355.00	71.27\\
378.00	70.90\\
402.00	70.90\\
427.00	70.96\\
455.00	70.77\\
484.00	70.87\\
515.00	71.26\\
547.00	71.12\\
582.00	70.10\\
619.00	70.76\\
659.00	70.71\\
701.00	70.90\\
746.00	71.72\\
793.00	72.38\\
844.00	72.27\\
897.00	72.24\\
955.00	72.25\\
1016.00	72.24\\
1080.00	72.41\\
1149.00	72.06\\
1222.00	72.42\\
1300.00	72.62\\
1383.00	72.60\\
1472.00	72.15\\
1565.00	72.91\\
1665.00	72.97\\
1771.00	73.01\\
1884.00	73.09\\
2005.00	73.32\\
2132.00	73.78\\
2268.00	73.59\\
2413.00	73.35\\
2567.00	73.16\\
2730.00	72.89\\
2905.00	72.67\\
3090.00	72.78\\
3287.00	72.68\\
3496.00	72.71\\
3719.00	72.30\\
3957.00	72.28\\
4209.00	72.39\\
4477.00	72.35\\
4763.00	72.45\\
5066.00	72.48\\
5389.00	72.50\\
5733.00	72.41\\
6099.00	72.42\\
6488.00	72.53\\
6901.00	72.32\\
7341.00	72.51\\
7809.00	72.56\\
8307.00	72.49\\
8837.00	72.47\\
9401.00	72.51\\
10000.00	72.46\\
};
\addplot [color=black!5!lightgray, dashdotted, line width=1.2pt, forget plot]
  table[row sep=crcr]{%
1.00	100.00\\
2.00	100.00\\
3.00	100.00\\
4.00	75.00\\
5.00	60.00\\
6.00	66.67\\
7.00	71.43\\
8.00	75.00\\
9.00	77.78\\
10.00	80.00\\
11.00	81.82\\
12.00	83.33\\
13.00	84.62\\
14.00	85.71\\
15.00	80.00\\
16.00	75.00\\
17.00	76.47\\
18.00	77.78\\
19.00	73.68\\
21.00	76.19\\
22.00	72.73\\
23.00	73.91\\
25.00	72.00\\
26.00	69.23\\
28.00	64.29\\
30.00	66.67\\
32.00	68.75\\
34.00	70.59\\
36.00	72.22\\
38.00	73.68\\
41.00	73.17\\
43.00	72.09\\
46.00	71.74\\
49.00	71.43\\
52.00	71.15\\
56.00	73.21\\
59.00	72.88\\
63.00	74.60\\
67.00	71.64\\
71.00	71.83\\
76.00	72.37\\
81.00	70.37\\
86.00	70.93\\
91.00	70.33\\
97.00	71.13\\
103.00	68.93\\
110.00	69.09\\
117.00	67.52\\
124.00	66.94\\
132.00	67.42\\
140.00	65.71\\
149.00	66.44\\
159.00	66.04\\
169.00	66.27\\
180.00	65.56\\
191.00	65.97\\
204.00	65.69\\
217.00	66.82\\
230.00	66.52\\
245.00	66.94\\
261.00	67.43\\
277.00	67.15\\
295.00	68.14\\
314.00	67.20\\
334.00	67.07\\
355.00	66.76\\
378.00	65.87\\
402.00	65.17\\
427.00	66.28\\
455.00	67.03\\
484.00	67.36\\
515.00	67.77\\
547.00	67.82\\
582.00	67.53\\
619.00	68.01\\
659.00	67.83\\
701.00	68.19\\
746.00	68.77\\
793.00	69.23\\
844.00	68.25\\
897.00	67.78\\
955.00	69.01\\
1016.00	69.09\\
1080.00	69.35\\
1149.00	69.71\\
1222.00	69.97\\
1300.00	69.69\\
1383.00	69.56\\
1472.00	69.90\\
1565.00	69.39\\
1665.00	69.07\\
1771.00	68.61\\
1884.00	68.31\\
2005.00	67.93\\
2132.00	68.34\\
2268.00	68.12\\
2413.00	67.84\\
2567.00	68.06\\
2730.00	68.21\\
2905.00	68.16\\
3090.00	68.16\\
3287.00	68.39\\
3496.00	68.31\\
3719.00	68.16\\
3957.00	68.01\\
4209.00	67.95\\
4477.00	67.88\\
4763.00	68.09\\
5066.00	67.88\\
5389.00	67.95\\
5733.00	67.89\\
6099.00	67.73\\
6488.00	67.66\\
6901.00	67.82\\
7341.00	67.80\\
7809.00	67.74\\
8307.00	67.79\\
8837.00	67.81\\
9401.00	67.80\\
10000.00	67.80\\
};
\addplot [color=white!80!black, line width=1.2pt, forget plot]
  table[row sep=crcr]{%
1.00	100.00\\
2.00	50.00\\
3.00	66.67\\
4.00	75.00\\
5.00	60.00\\
6.00	66.67\\
7.00	71.43\\
8.00	75.00\\
9.00	77.78\\
10.00	70.00\\
11.00	72.73\\
12.00	75.00\\
13.00	69.23\\
14.00	64.29\\
15.00	66.67\\
16.00	68.75\\
17.00	64.71\\
18.00	66.67\\
19.00	68.42\\
21.00	66.67\\
22.00	68.18\\
23.00	69.57\\
25.00	72.00\\
26.00	69.23\\
28.00	71.43\\
30.00	66.67\\
32.00	68.75\\
34.00	70.59\\
36.00	69.44\\
38.00	68.42\\
41.00	63.41\\
43.00	65.12\\
46.00	65.22\\
49.00	61.22\\
52.00	61.54\\
56.00	62.50\\
59.00	64.41\\
63.00	65.08\\
67.00	65.67\\
71.00	63.38\\
76.00	64.47\\
81.00	66.67\\
86.00	67.44\\
91.00	65.93\\
97.00	65.98\\
103.00	66.99\\
110.00	66.36\\
117.00	66.67\\
124.00	66.94\\
132.00	66.67\\
140.00	67.14\\
149.00	67.79\\
159.00	68.55\\
169.00	69.23\\
180.00	69.44\\
191.00	68.59\\
204.00	70.59\\
217.00	71.43\\
230.00	71.74\\
245.00	71.43\\
261.00	70.50\\
277.00	70.04\\
295.00	70.51\\
314.00	70.38\\
334.00	70.36\\
355.00	70.42\\
378.00	70.90\\
402.00	70.40\\
427.00	70.02\\
455.00	70.11\\
484.00	70.87\\
515.00	70.68\\
547.00	70.75\\
582.00	69.93\\
619.00	69.14\\
659.00	68.59\\
701.00	68.47\\
746.00	68.10\\
793.00	67.84\\
844.00	67.89\\
897.00	68.45\\
955.00	68.48\\
1016.00	69.00\\
1080.00	69.44\\
1149.00	69.89\\
1222.00	69.48\\
1300.00	68.85\\
1383.00	68.62\\
1472.00	68.82\\
1565.00	68.50\\
1665.00	68.17\\
1771.00	68.94\\
1884.00	68.74\\
2005.00	68.28\\
2132.00	68.34\\
2268.00	68.61\\
2413.00	68.34\\
2567.00	68.64\\
2730.00	69.05\\
2905.00	68.85\\
3090.00	68.90\\
3287.00	69.09\\
3496.00	69.28\\
3719.00	69.45\\
3957.00	69.42\\
4209.00	69.64\\
4477.00	69.44\\
4763.00	69.54\\
5066.00	69.54\\
5389.00	69.46\\
5733.00	69.23\\
6099.00	69.09\\
6488.00	69.11\\
6901.00	69.05\\
7341.00	69.13\\
7809.00	69.11\\
8307.00	69.04\\
8837.00	68.65\\
9401.00	68.66\\
10000.00	68.78\\
};
\end{axis}
\end{tikzpicture}%
    \caption{Laufende Konvergenz der Trennschärfe für alle Modellparameter. Die Schätzungen stabilisieren sich bei der $\sym{n}_{\sym{idx_MC}}=10^4$ reproduzierbar.}
    \label{fig:app.c_conv_power}
\end{figure}

\begin{figure}[htbp]
    \centering
    \setlength\figurewidth{0.9\textwidth}
    \setlength\figureheight{16cm}
    % This file was created by matlab2tikz.
%
\definecolor{mycolor1}{rgb}{0.12941,0.12941,0.12941}%
%
\begin{tikzpicture}

\begin{axis}[%
width=0.411\figurewidth,
height=0.265\figureheight,
at={(0\figurewidth,0.735\figureheight)},
scale only axis,
xmin=0.00,
xmax=10000.00,
xtick={0.00,5000.00,10000.00},
xticklabels={\empty},
ymin=0.00,
ymax=40.00,
ylabel style={font=\color{mycolor1}},
ylabel={$\hat{\sym{thet}}_{\sym{idx_i}}$},
axis background/.style={fill=white},
title style={font=\bfseries\color{mycolor1}},
title={\textbf{$\sym{beta}_0$}},
xmajorgrids,
ymajorgrids,
grid style={dotted, opacity=0.25},
legend style={legend cell align=left, align=left}
]
\addplot [color=white!10!black, line width=1.2pt]
  table[row sep=crcr]{%
1.00	18.88\\
2.00	20.02\\
3.00	21.78\\
4.00	22.65\\
5.00	23.34\\
6.00	23.75\\
7.00	23.91\\
8.00	23.46\\
9.00	23.36\\
10.00	22.78\\
11.00	22.65\\
12.00	22.88\\
13.00	22.66\\
14.00	22.48\\
15.00	22.19\\
16.00	22.07\\
17.00	21.44\\
18.00	21.48\\
19.00	21.95\\
21.00	22.68\\
22.00	22.99\\
23.00	23.42\\
25.00	23.76\\
26.00	24.08\\
28.00	23.90\\
30.00	23.90\\
32.00	23.59\\
34.00	23.81\\
36.00	23.64\\
38.00	23.62\\
41.00	23.55\\
43.00	23.25\\
46.00	23.31\\
49.00	23.28\\
52.00	23.60\\
56.00	23.70\\
59.00	23.66\\
63.00	23.76\\
67.00	23.61\\
71.00	23.67\\
76.00	23.95\\
81.00	24.01\\
86.00	24.03\\
91.00	24.13\\
97.00	24.10\\
103.00	23.90\\
117.00	23.90\\
124.00	23.82\\
132.00	23.84\\
149.00	23.64\\
159.00	23.76\\
169.00	23.77\\
180.00	23.68\\
191.00	23.67\\
204.00	23.58\\
217.00	23.69\\
230.00	23.67\\
245.00	23.71\\
261.00	23.52\\
277.00	23.44\\
295.00	23.59\\
314.00	23.72\\
427.00	23.79\\
455.00	23.74\\
515.00	23.96\\
547.00	24.02\\
619.00	24.04\\
659.00	24.03\\
701.00	24.07\\
746.00	24.06\\
793.00	24.12\\
844.00	24.12\\
897.00	24.02\\
955.00	24.07\\
1080.00	24.08\\
1149.00	24.09\\
1383.00	23.98\\
1665.00	23.91\\
1884.00	23.94\\
2268.00	23.91\\
2413.00	23.96\\
2730.00	23.95\\
3719.00	23.94\\
3957.00	23.97\\
4477.00	23.95\\
6488.00	23.99\\
7341.00	24.03\\
10000.00	24.05\\
};
\addlegendentry{Basis \ac{CCD}}

\addplot [color=white!20!black, dashed, line width=1.2pt]
  table[row sep=crcr]{%
1.00	31.50\\
2.00	30.02\\
4.00	26.62\\
5.00	26.13\\
6.00	25.74\\
7.00	25.08\\
8.00	25.53\\
9.00	25.51\\
10.00	24.80\\
11.00	24.26\\
12.00	23.84\\
13.00	24.18\\
14.00	23.91\\
15.00	24.18\\
16.00	24.22\\
17.00	24.08\\
18.00	23.80\\
19.00	23.39\\
21.00	23.71\\
22.00	23.51\\
23.00	23.25\\
25.00	23.65\\
26.00	23.93\\
28.00	23.76\\
30.00	24.01\\
32.00	24.06\\
34.00	24.24\\
36.00	24.33\\
38.00	24.78\\
41.00	24.79\\
43.00	24.85\\
46.00	24.88\\
49.00	24.73\\
52.00	24.86\\
56.00	24.58\\
59.00	24.70\\
63.00	24.64\\
67.00	24.62\\
71.00	24.63\\
76.00	24.58\\
81.00	24.82\\
86.00	24.91\\
97.00	25.25\\
103.00	25.12\\
110.00	24.79\\
117.00	24.79\\
124.00	24.61\\
132.00	24.56\\
140.00	24.55\\
149.00	24.38\\
159.00	24.42\\
191.00	24.26\\
204.00	24.12\\
217.00	24.09\\
230.00	24.15\\
245.00	24.13\\
261.00	24.30\\
277.00	24.31\\
295.00	24.53\\
314.00	24.56\\
334.00	24.62\\
355.00	24.59\\
378.00	24.39\\
402.00	24.38\\
427.00	24.31\\
455.00	24.11\\
515.00	24.06\\
547.00	24.10\\
582.00	24.06\\
619.00	23.95\\
746.00	23.91\\
793.00	23.97\\
897.00	23.93\\
955.00	24.04\\
1149.00	24.08\\
1222.00	24.06\\
1383.00	24.09\\
1472.00	24.15\\
1565.00	24.13\\
1665.00	24.18\\
1771.00	24.19\\
2005.00	24.13\\
2268.00	24.15\\
2413.00	24.08\\
2730.00	24.08\\
2905.00	24.05\\
3090.00	24.06\\
3287.00	24.00\\
3719.00	24.01\\
4209.00	23.99\\
6488.00	23.97\\
6901.00	24.00\\
10000.00	23.97\\
};
\addlegendentry{Fall I ($-20\,\%$)}

\addplot [color=white!30!black, dotted, line width=1.2pt]
  table[row sep=crcr]{%
1.00	26.52\\
2.00	26.28\\
3.00	25.21\\
4.00	25.53\\
5.00	23.99\\
6.00	23.55\\
7.00	24.41\\
8.00	23.39\\
9.00	23.21\\
10.00	23.93\\
11.00	24.45\\
12.00	24.18\\
14.00	24.01\\
15.00	24.76\\
16.00	24.58\\
17.00	24.82\\
18.00	24.97\\
19.00	24.49\\
21.00	24.89\\
22.00	25.24\\
23.00	25.05\\
25.00	25.05\\
26.00	24.71\\
28.00	24.86\\
30.00	24.77\\
32.00	24.85\\
34.00	24.66\\
36.00	24.51\\
38.00	24.53\\
41.00	24.52\\
43.00	24.61\\
49.00	24.63\\
52.00	24.36\\
59.00	24.09\\
63.00	24.25\\
67.00	24.30\\
76.00	24.29\\
81.00	24.36\\
86.00	24.34\\
91.00	24.42\\
97.00	24.44\\
110.00	24.40\\
124.00	24.22\\
132.00	24.31\\
140.00	24.14\\
159.00	23.96\\
169.00	23.93\\
180.00	23.81\\
191.00	23.80\\
204.00	23.86\\
217.00	23.70\\
277.00	23.93\\
295.00	23.94\\
314.00	24.08\\
334.00	24.10\\
355.00	24.17\\
402.00	24.05\\
427.00	24.10\\
455.00	24.07\\
515.00	24.15\\
547.00	24.17\\
582.00	24.06\\
619.00	24.21\\
659.00	24.17\\
701.00	24.20\\
793.00	24.18\\
897.00	24.26\\
1016.00	24.31\\
1149.00	24.24\\
1222.00	24.24\\
1300.00	24.20\\
1383.00	24.26\\
1665.00	24.22\\
1771.00	24.27\\
1884.00	24.27\\
2005.00	24.32\\
2132.00	24.32\\
2413.00	24.23\\
2730.00	24.23\\
3287.00	24.21\\
3496.00	24.23\\
3957.00	24.20\\
4209.00	24.20\\
4763.00	24.14\\
9401.00	24.00\\
10000.00	24.01\\
};
\addlegendentry{Fall I ($+20\,\%$)}

\addplot [color=white!40!black, dashdotted, line width=1.2pt]
  table[row sep=crcr]{%
1.00	26.55\\
2.00	26.91\\
3.00	27.20\\
4.00	27.92\\
5.00	25.37\\
6.00	25.15\\
7.00	24.42\\
8.00	24.20\\
9.00	25.00\\
10.00	24.25\\
11.00	24.14\\
12.00	23.89\\
13.00	24.53\\
14.00	24.99\\
15.00	25.05\\
16.00	25.06\\
18.00	25.23\\
19.00	25.62\\
21.00	25.99\\
22.00	26.14\\
23.00	26.00\\
25.00	25.37\\
26.00	25.13\\
28.00	25.32\\
30.00	25.19\\
32.00	25.39\\
34.00	25.69\\
36.00	25.40\\
38.00	25.16\\
41.00	25.06\\
43.00	24.70\\
46.00	24.47\\
49.00	24.54\\
52.00	24.51\\
56.00	24.36\\
59.00	24.55\\
63.00	24.64\\
67.00	24.50\\
71.00	24.58\\
76.00	24.85\\
81.00	24.67\\
86.00	24.72\\
91.00	24.99\\
97.00	25.09\\
103.00	24.91\\
110.00	24.97\\
117.00	24.93\\
124.00	24.65\\
132.00	24.73\\
140.00	24.92\\
159.00	24.78\\
169.00	24.80\\
180.00	24.74\\
191.00	24.59\\
204.00	24.69\\
217.00	24.62\\
230.00	24.67\\
245.00	24.81\\
261.00	24.82\\
277.00	24.77\\
295.00	24.89\\
314.00	24.72\\
334.00	24.61\\
355.00	24.57\\
378.00	24.49\\
402.00	24.47\\
427.00	24.38\\
484.00	24.36\\
515.00	24.32\\
547.00	24.35\\
619.00	24.35\\
701.00	24.38\\
746.00	24.37\\
793.00	24.43\\
1080.00	24.23\\
1222.00	24.23\\
1383.00	24.20\\
1472.00	24.24\\
1771.00	24.17\\
2005.00	24.12\\
3496.00	24.07\\
3719.00	24.09\\
5066.00	24.01\\
5389.00	24.03\\
5733.00	24.02\\
6099.00	24.05\\
7809.00	23.99\\
10000.00	23.99\\
};
\addlegendentry{Fall I (\ac{FC})}

\addplot [color=gray, line width=1.2pt]
  table[row sep=crcr]{%
1.00	21.02\\
2.00	23.33\\
3.00	23.76\\
4.00	23.91\\
5.00	24.14\\
6.00	23.84\\
7.00	25.59\\
8.00	26.10\\
9.00	25.49\\
10.00	25.83\\
11.00	25.77\\
12.00	25.59\\
13.00	25.77\\
15.00	25.87\\
16.00	25.28\\
18.00	24.93\\
19.00	24.57\\
22.00	24.20\\
25.00	24.03\\
26.00	24.14\\
28.00	23.79\\
30.00	23.92\\
32.00	24.11\\
34.00	24.06\\
38.00	23.32\\
41.00	23.22\\
43.00	23.46\\
49.00	23.63\\
52.00	23.74\\
56.00	23.61\\
59.00	23.63\\
63.00	23.61\\
67.00	23.65\\
71.00	23.61\\
76.00	23.68\\
81.00	23.64\\
86.00	24.01\\
97.00	23.96\\
103.00	23.82\\
110.00	23.90\\
117.00	23.82\\
132.00	23.91\\
140.00	24.02\\
149.00	24.20\\
159.00	24.05\\
169.00	24.18\\
180.00	24.17\\
191.00	24.31\\
204.00	24.26\\
217.00	24.31\\
230.00	24.27\\
245.00	24.13\\
295.00	24.40\\
314.00	24.43\\
355.00	24.39\\
378.00	24.34\\
402.00	24.39\\
427.00	24.38\\
455.00	24.44\\
484.00	24.41\\
515.00	24.30\\
582.00	24.20\\
746.00	24.16\\
844.00	24.31\\
897.00	24.38\\
955.00	24.29\\
1080.00	24.28\\
1149.00	24.27\\
1300.00	24.15\\
1565.00	24.10\\
1665.00	24.05\\
2005.00	24.05\\
2413.00	24.01\\
2730.00	24.00\\
4763.00	23.99\\
6488.00	24.05\\
8307.00	24.02\\
9401.00	24.03\\
10000.00	24.03\\
};
\addlegendentry{Fall II (leicht)}

\addplot [color=white!60!black, dashed, line width=1.2pt]
  table[row sep=crcr]{%
1.00	22.23\\
2.00	27.87\\
3.00	25.21\\
4.00	26.02\\
5.00	25.53\\
6.00	25.12\\
7.00	24.21\\
8.00	23.95\\
9.00	24.07\\
10.00	24.78\\
11.00	26.10\\
12.00	25.86\\
13.00	25.76\\
14.00	25.19\\
15.00	25.48\\
16.00	25.96\\
17.00	25.74\\
18.00	25.74\\
19.00	25.28\\
21.00	25.59\\
22.00	26.14\\
23.00	26.28\\
25.00	25.92\\
26.00	25.92\\
28.00	25.26\\
30.00	25.22\\
32.00	25.34\\
34.00	25.25\\
36.00	25.25\\
38.00	25.33\\
41.00	24.85\\
43.00	24.86\\
46.00	24.76\\
52.00	24.32\\
56.00	24.52\\
59.00	24.62\\
63.00	24.59\\
67.00	24.70\\
71.00	24.34\\
76.00	24.06\\
81.00	24.30\\
91.00	24.36\\
97.00	24.30\\
103.00	24.42\\
110.00	24.52\\
117.00	24.71\\
124.00	24.65\\
132.00	24.42\\
140.00	24.55\\
159.00	24.57\\
169.00	24.51\\
180.00	24.49\\
191.00	24.37\\
204.00	24.48\\
217.00	24.43\\
230.00	24.26\\
277.00	24.53\\
334.00	24.33\\
355.00	24.18\\
378.00	24.16\\
402.00	24.20\\
427.00	24.11\\
455.00	24.07\\
484.00	24.08\\
515.00	23.95\\
619.00	24.00\\
659.00	23.93\\
701.00	23.82\\
746.00	23.80\\
844.00	23.86\\
897.00	23.81\\
1016.00	23.84\\
1383.00	23.89\\
1472.00	23.96\\
1565.00	23.97\\
1665.00	23.93\\
1771.00	24.00\\
2005.00	23.96\\
2132.00	23.94\\
2413.00	23.95\\
2567.00	23.95\\
2730.00	24.00\\
2905.00	24.02\\
3957.00	23.96\\
4209.00	23.96\\
4763.00	23.89\\
5733.00	23.83\\
6099.00	23.87\\
7341.00	23.87\\
7809.00	23.89\\
8837.00	23.87\\
10000.00	23.91\\
};
\addlegendentry{Fall II (stark)}

\addplot [color=white!70!black, dotted, line width=1.2pt]
  table[row sep=crcr]{%
1.00	26.02\\
2.00	30.32\\
3.00	30.53\\
4.00	29.40\\
5.00	29.08\\
6.00	28.64\\
7.00	29.96\\
8.00	28.95\\
9.00	28.85\\
10.00	27.86\\
11.00	27.35\\
12.00	27.31\\
13.00	26.99\\
14.00	26.53\\
15.00	26.57\\
16.00	26.55\\
17.00	26.31\\
18.00	26.01\\
19.00	25.97\\
21.00	25.38\\
22.00	25.77\\
23.00	25.39\\
25.00	25.73\\
26.00	25.93\\
28.00	25.57\\
30.00	25.29\\
32.00	25.44\\
34.00	25.08\\
36.00	25.21\\
38.00	25.09\\
41.00	25.15\\
43.00	25.15\\
46.00	24.73\\
52.00	24.20\\
56.00	24.36\\
59.00	24.60\\
63.00	24.80\\
67.00	24.81\\
71.00	24.76\\
76.00	24.79\\
86.00	24.75\\
91.00	24.38\\
97.00	24.37\\
103.00	24.57\\
110.00	24.74\\
117.00	24.56\\
124.00	24.66\\
132.00	24.88\\
149.00	24.90\\
159.00	24.68\\
169.00	24.51\\
180.00	24.49\\
191.00	24.26\\
204.00	24.51\\
217.00	24.53\\
230.00	24.44\\
245.00	24.50\\
261.00	24.68\\
277.00	24.56\\
295.00	24.50\\
314.00	24.54\\
334.00	24.51\\
355.00	24.41\\
378.00	24.25\\
427.00	24.31\\
455.00	24.20\\
484.00	24.20\\
547.00	24.06\\
746.00	23.94\\
793.00	23.98\\
1016.00	23.95\\
1080.00	24.02\\
1149.00	24.04\\
1222.00	23.97\\
1565.00	23.96\\
2268.00	23.98\\
2905.00	24.01\\
3496.00	24.01\\
5733.00	24.06\\
7341.00	24.05\\
10000.00	24.07\\
};
\addlegendentry{Fall III}

\addplot [color=white!80!black, dashdotted, line width=1.2pt]
  table[row sep=crcr]{%
1.00	22.56\\
2.00	27.87\\
3.00	26.04\\
4.00	24.80\\
5.00	24.08\\
6.00	24.37\\
7.00	24.09\\
8.00	22.96\\
9.00	22.84\\
10.00	23.42\\
11.00	22.63\\
12.00	22.57\\
13.00	22.38\\
14.00	22.49\\
15.00	22.80\\
16.00	22.49\\
17.00	22.59\\
18.00	22.62\\
19.00	22.97\\
21.00	23.08\\
22.00	22.92\\
23.00	23.00\\
26.00	23.42\\
28.00	23.53\\
30.00	23.76\\
32.00	23.85\\
34.00	23.80\\
36.00	23.61\\
38.00	23.14\\
41.00	23.31\\
43.00	23.39\\
46.00	23.43\\
49.00	23.82\\
52.00	23.97\\
59.00	24.29\\
67.00	23.65\\
71.00	23.79\\
76.00	23.73\\
86.00	24.05\\
91.00	24.02\\
97.00	24.16\\
103.00	24.00\\
110.00	24.14\\
117.00	24.18\\
124.00	24.13\\
132.00	23.97\\
140.00	24.02\\
149.00	24.21\\
159.00	24.06\\
169.00	24.05\\
180.00	23.98\\
191.00	23.87\\
204.00	23.86\\
217.00	23.81\\
230.00	23.66\\
245.00	23.71\\
261.00	23.82\\
295.00	23.79\\
314.00	23.82\\
334.00	23.70\\
402.00	23.59\\
427.00	23.54\\
455.00	23.55\\
484.00	23.64\\
659.00	23.76\\
793.00	23.68\\
844.00	23.76\\
897.00	23.76\\
955.00	23.84\\
1080.00	23.80\\
1149.00	23.84\\
1472.00	23.78\\
1665.00	23.83\\
1771.00	23.79\\
1884.00	23.79\\
2005.00	23.82\\
2132.00	23.81\\
2268.00	23.85\\
2413.00	23.85\\
2567.00	23.91\\
2905.00	23.93\\
3090.00	23.91\\
3287.00	23.93\\
3719.00	23.89\\
4477.00	23.91\\
5066.00	23.88\\
5389.00	23.91\\
9401.00	23.99\\
10000.00	23.99\\
};
\addlegendentry{Fall IV: \ac{SP}}

\addplot [color=white!90!black, line width=1.2pt]
  table[row sep=crcr]{%
1.00	26.42\\
2.00	24.68\\
3.00	24.64\\
4.00	24.02\\
5.00	20.83\\
6.00	20.61\\
7.00	21.47\\
8.00	20.56\\
9.00	20.04\\
10.00	19.89\\
11.00	19.26\\
12.00	19.01\\
13.00	18.96\\
14.00	19.10\\
15.00	19.51\\
16.00	20.43\\
17.00	20.77\\
18.00	20.76\\
19.00	20.98\\
21.00	21.14\\
22.00	21.33\\
23.00	21.86\\
25.00	22.61\\
26.00	22.68\\
28.00	23.04\\
30.00	23.06\\
32.00	23.35\\
34.00	23.47\\
36.00	23.90\\
38.00	23.85\\
41.00	24.14\\
43.00	24.00\\
46.00	24.00\\
49.00	23.85\\
52.00	23.88\\
56.00	24.13\\
59.00	24.08\\
71.00	24.34\\
76.00	24.31\\
81.00	24.32\\
86.00	24.24\\
91.00	24.33\\
97.00	24.39\\
103.00	24.29\\
110.00	24.25\\
117.00	24.08\\
124.00	24.01\\
132.00	24.22\\
140.00	24.19\\
149.00	24.28\\
159.00	24.26\\
169.00	24.09\\
191.00	24.15\\
204.00	24.11\\
217.00	24.20\\
230.00	24.19\\
245.00	24.11\\
277.00	24.16\\
314.00	24.23\\
334.00	24.16\\
355.00	24.23\\
378.00	24.27\\
427.00	24.14\\
455.00	24.13\\
515.00	24.18\\
547.00	24.30\\
619.00	24.32\\
659.00	24.36\\
746.00	24.26\\
793.00	24.20\\
844.00	24.24\\
955.00	24.17\\
1149.00	24.05\\
1222.00	24.03\\
1472.00	24.09\\
1665.00	24.06\\
1771.00	24.05\\
1884.00	24.08\\
2905.00	24.03\\
3719.00	24.02\\
3957.00	24.03\\
4209.00	24.00\\
5066.00	23.99\\
5733.00	23.98\\
6901.00	24.02\\
7809.00	24.01\\
8837.00	24.04\\
10000.00	24.04\\
};
\addlegendentry{Fall IV: \ac{ZP}}

\addplot [color=black, line width=1.2pt]
  table[row sep=crcr]{%
0.00	24.00\\
10000.00	24.00\\
};
\addlegendentry{Wahrer Wert}

\end{axis}

\begin{axis}[%
width=0.411\figurewidth,
height=0.265\figureheight,
at={(0.54\figurewidth,0.735\figureheight)},
scale only axis,
xmin=0.00,
xmax=10000.00,
xtick={0.00,5000.00,10000.00},
xticklabels={\empty},
ymin=0.00,
ymax=40.00,
ytick={0.00,10.00,20.00,30.00,40.00},
yticklabels={\empty},
axis background/.style={fill=white},
title style={font=\bfseries\color{mycolor1}},
title={\textbf{$\sym{beta}_1$}},
xmajorgrids,
ymajorgrids,
grid style={dotted, opacity=0.25}
]
\addplot [color=white!10!black, line width=1.2pt, forget plot]
  table[row sep=crcr]{%
1.00	19.31\\
2.00	15.22\\
3.00	12.65\\
4.00	12.13\\
5.00	12.74\\
7.00	12.72\\
8.00	13.47\\
9.00	12.94\\
10.00	13.61\\
11.00	13.14\\
12.00	13.58\\
13.00	13.29\\
14.00	13.06\\
15.00	13.20\\
16.00	12.77\\
17.00	12.81\\
18.00	12.37\\
19.00	12.30\\
22.00	11.89\\
23.00	11.83\\
25.00	12.06\\
26.00	12.06\\
28.00	12.12\\
30.00	12.09\\
32.00	12.52\\
34.00	12.06\\
36.00	11.52\\
38.00	12.08\\
41.00	12.11\\
43.00	12.42\\
46.00	12.43\\
49.00	12.31\\
52.00	12.45\\
56.00	12.38\\
59.00	12.45\\
63.00	12.37\\
71.00	12.08\\
76.00	12.02\\
81.00	12.00\\
86.00	12.02\\
91.00	11.97\\
103.00	11.66\\
110.00	11.79\\
117.00	11.81\\
124.00	11.99\\
132.00	11.95\\
140.00	12.16\\
159.00	12.32\\
169.00	12.23\\
180.00	12.24\\
191.00	12.20\\
204.00	12.08\\
217.00	12.08\\
230.00	11.99\\
245.00	11.99\\
277.00	11.90\\
314.00	11.95\\
378.00	11.71\\
402.00	11.62\\
427.00	11.63\\
455.00	11.54\\
484.00	11.52\\
582.00	11.72\\
619.00	11.73\\
659.00	11.81\\
701.00	11.84\\
793.00	11.79\\
1016.00	11.92\\
1080.00	11.88\\
1149.00	11.90\\
1222.00	11.85\\
1383.00	11.91\\
1472.00	11.87\\
1665.00	11.88\\
2005.00	11.88\\
2132.00	11.91\\
2268.00	11.85\\
3719.00	11.97\\
3957.00	11.97\\
4209.00	11.93\\
5066.00	12.02\\
5733.00	12.00\\
6099.00	12.03\\
6901.00	12.00\\
10000.00	12.01\\
};
\addplot [color=white!20!black, dashed, line width=1.2pt, forget plot]
  table[row sep=crcr]{%
1.00	12.04\\
2.00	15.44\\
3.00	15.32\\
4.00	13.92\\
5.00	13.96\\
6.00	13.27\\
7.00	13.52\\
8.00	12.60\\
9.00	12.53\\
10.00	12.26\\
11.00	12.48\\
12.00	11.69\\
13.00	12.24\\
14.00	12.52\\
15.00	12.28\\
16.00	12.24\\
17.00	12.27\\
18.00	12.21\\
19.00	12.36\\
21.00	11.62\\
22.00	12.02\\
23.00	12.74\\
25.00	12.73\\
26.00	12.44\\
30.00	12.19\\
32.00	11.97\\
34.00	12.00\\
38.00	11.59\\
41.00	11.25\\
43.00	11.12\\
46.00	11.30\\
49.00	11.13\\
52.00	11.29\\
56.00	11.06\\
59.00	10.96\\
63.00	10.87\\
67.00	11.11\\
71.00	11.30\\
76.00	11.39\\
81.00	11.34\\
86.00	11.41\\
91.00	11.52\\
97.00	11.55\\
103.00	11.68\\
110.00	11.73\\
117.00	11.95\\
124.00	12.11\\
140.00	11.99\\
149.00	12.22\\
159.00	12.16\\
169.00	12.21\\
180.00	12.38\\
204.00	12.50\\
217.00	12.44\\
230.00	12.46\\
245.00	12.53\\
261.00	12.54\\
277.00	12.51\\
314.00	12.35\\
378.00	12.34\\
402.00	12.40\\
427.00	12.37\\
455.00	12.41\\
484.00	12.42\\
515.00	12.36\\
582.00	12.40\\
619.00	12.33\\
659.00	12.32\\
701.00	12.27\\
746.00	12.31\\
793.00	12.29\\
844.00	12.32\\
897.00	12.30\\
955.00	12.34\\
1016.00	12.26\\
1080.00	12.31\\
1149.00	12.22\\
1222.00	12.18\\
1300.00	12.19\\
1472.00	12.09\\
1665.00	12.12\\
1771.00	12.09\\
2268.00	12.12\\
2567.00	12.05\\
2730.00	12.07\\
3090.00	12.06\\
3287.00	12.10\\
3957.00	12.02\\
6099.00	12.03\\
7341.00	12.06\\
9401.00	12.01\\
10000.00	12.03\\
};
\addplot [color=white!30!black, dotted, line width=1.2pt, forget plot]
  table[row sep=crcr]{%
1.00	12.61\\
2.00	15.53\\
3.00	16.54\\
4.00	16.43\\
5.00	15.09\\
6.00	14.43\\
7.00	14.30\\
8.00	13.89\\
9.00	13.80\\
10.00	12.95\\
11.00	13.13\\
12.00	13.15\\
13.00	13.38\\
14.00	13.10\\
15.00	13.89\\
16.00	13.08\\
17.00	12.99\\
19.00	13.56\\
21.00	13.93\\
22.00	13.94\\
23.00	13.74\\
25.00	13.56\\
26.00	13.62\\
28.00	13.71\\
30.00	13.42\\
32.00	13.44\\
34.00	13.37\\
36.00	13.35\\
41.00	13.62\\
43.00	13.58\\
46.00	13.46\\
49.00	13.20\\
52.00	13.49\\
56.00	13.64\\
63.00	13.74\\
67.00	13.72\\
71.00	13.45\\
76.00	13.34\\
81.00	13.32\\
91.00	13.11\\
110.00	12.47\\
117.00	12.43\\
124.00	12.30\\
132.00	12.08\\
140.00	12.12\\
149.00	12.40\\
159.00	12.31\\
169.00	12.16\\
180.00	12.08\\
204.00	12.12\\
217.00	12.22\\
230.00	12.16\\
245.00	12.25\\
261.00	12.20\\
277.00	12.28\\
295.00	12.41\\
314.00	12.47\\
355.00	12.40\\
378.00	12.46\\
547.00	12.16\\
582.00	12.15\\
659.00	12.22\\
746.00	12.16\\
793.00	12.21\\
844.00	12.14\\
1080.00	12.02\\
1222.00	12.09\\
1565.00	11.97\\
1665.00	11.99\\
1884.00	11.97\\
2005.00	11.99\\
2132.00	11.94\\
2413.00	12.02\\
2730.00	12.04\\
2905.00	12.00\\
3957.00	11.95\\
4209.00	11.92\\
5066.00	11.93\\
5733.00	11.96\\
10000.00	11.99\\
};
\addplot [color=white!40!black, dashdotted, line width=1.2pt, forget plot]
  table[row sep=crcr]{%
1.00	17.12\\
2.00	12.60\\
3.00	12.73\\
4.00	11.74\\
5.00	12.20\\
6.00	10.63\\
7.00	10.35\\
8.00	11.55\\
9.00	11.86\\
10.00	11.83\\
11.00	11.52\\
12.00	12.40\\
13.00	12.99\\
14.00	13.08\\
15.00	13.02\\
16.00	13.22\\
17.00	12.81\\
18.00	12.75\\
19.00	12.76\\
21.00	13.14\\
22.00	13.19\\
23.00	13.55\\
25.00	13.44\\
26.00	13.14\\
28.00	13.13\\
30.00	12.78\\
32.00	12.61\\
34.00	12.21\\
36.00	12.41\\
38.00	12.09\\
43.00	11.80\\
46.00	11.98\\
49.00	12.10\\
52.00	11.98\\
56.00	12.02\\
59.00	11.86\\
63.00	12.04\\
76.00	12.35\\
81.00	12.29\\
86.00	12.35\\
91.00	12.16\\
97.00	12.57\\
110.00	12.47\\
117.00	12.64\\
132.00	12.86\\
140.00	12.80\\
149.00	12.89\\
159.00	12.81\\
169.00	12.61\\
180.00	12.50\\
191.00	12.68\\
204.00	12.69\\
217.00	12.60\\
230.00	12.54\\
314.00	12.68\\
334.00	12.62\\
378.00	12.44\\
402.00	12.48\\
455.00	12.39\\
515.00	12.44\\
582.00	12.25\\
619.00	12.15\\
659.00	12.26\\
701.00	12.28\\
746.00	12.36\\
793.00	12.37\\
844.00	12.34\\
955.00	12.37\\
1080.00	12.22\\
1222.00	12.20\\
1300.00	12.25\\
1472.00	12.22\\
1665.00	12.28\\
1771.00	12.27\\
1884.00	12.21\\
2268.00	12.22\\
2413.00	12.18\\
2730.00	12.21\\
3090.00	12.11\\
4209.00	12.11\\
4477.00	12.13\\
6488.00	12.01\\
10000.00	12.01\\
};
\addplot [color=gray, line width=1.2pt, forget plot]
  table[row sep=crcr]{%
1.00	2.08\\
2.00	5.66\\
3.00	9.38\\
4.00	10.44\\
5.00	9.62\\
6.00	9.83\\
7.00	9.02\\
8.00	9.26\\
9.00	9.58\\
10.00	9.43\\
11.00	9.68\\
12.00	11.34\\
13.00	11.71\\
14.00	12.86\\
15.00	12.80\\
16.00	12.04\\
17.00	12.39\\
18.00	12.59\\
19.00	12.27\\
21.00	12.76\\
23.00	12.77\\
25.00	12.28\\
26.00	12.29\\
28.00	12.21\\
30.00	11.95\\
32.00	11.85\\
34.00	11.78\\
36.00	11.97\\
38.00	11.79\\
41.00	11.84\\
43.00	11.68\\
49.00	11.05\\
52.00	11.11\\
63.00	11.60\\
67.00	11.55\\
76.00	11.17\\
91.00	11.30\\
97.00	11.20\\
110.00	11.19\\
117.00	11.23\\
124.00	11.43\\
132.00	11.44\\
140.00	11.65\\
149.00	11.61\\
159.00	11.64\\
169.00	11.74\\
180.00	11.82\\
191.00	11.80\\
204.00	11.62\\
230.00	11.67\\
261.00	11.50\\
277.00	11.50\\
295.00	11.58\\
314.00	11.55\\
378.00	11.56\\
402.00	11.63\\
455.00	11.62\\
547.00	11.69\\
582.00	11.67\\
619.00	11.61\\
659.00	11.71\\
701.00	11.77\\
746.00	11.63\\
897.00	11.68\\
1016.00	11.69\\
1222.00	11.77\\
1472.00	11.72\\
1665.00	11.77\\
1771.00	11.79\\
1884.00	11.85\\
2005.00	11.85\\
2268.00	11.90\\
2567.00	11.89\\
2730.00	11.95\\
3090.00	11.93\\
3287.00	11.96\\
3496.00	11.93\\
10000.00	12.02\\
};
\addplot [color=white!60!black, dashed, line width=1.2pt, forget plot]
  table[row sep=crcr]{%
1.00	10.63\\
2.00	14.00\\
3.00	14.75\\
4.00	13.91\\
5.00	13.58\\
6.00	12.32\\
7.00	12.78\\
8.00	14.49\\
9.00	15.15\\
11.00	13.90\\
12.00	13.52\\
13.00	13.27\\
14.00	12.54\\
15.00	12.98\\
16.00	13.14\\
17.00	12.82\\
18.00	12.68\\
19.00	12.67\\
21.00	12.44\\
23.00	12.72\\
25.00	12.50\\
26.00	12.23\\
28.00	12.26\\
30.00	11.94\\
32.00	11.94\\
34.00	11.67\\
36.00	11.72\\
38.00	11.88\\
41.00	11.77\\
43.00	11.57\\
46.00	11.83\\
49.00	11.76\\
52.00	11.78\\
56.00	11.52\\
67.00	11.28\\
71.00	11.53\\
76.00	12.07\\
81.00	11.97\\
86.00	11.98\\
91.00	12.04\\
97.00	12.34\\
103.00	12.57\\
110.00	12.29\\
124.00	12.26\\
132.00	12.07\\
149.00	11.98\\
159.00	11.81\\
169.00	12.03\\
180.00	12.18\\
191.00	12.19\\
204.00	12.24\\
217.00	12.34\\
277.00	12.44\\
295.00	12.38\\
314.00	12.38\\
355.00	12.24\\
378.00	12.24\\
547.00	11.91\\
582.00	11.88\\
619.00	11.90\\
746.00	11.85\\
793.00	11.90\\
844.00	11.90\\
897.00	11.94\\
955.00	11.96\\
1016.00	11.93\\
1080.00	11.87\\
1383.00	11.98\\
2005.00	11.93\\
2132.00	11.95\\
2268.00	11.93\\
2413.00	11.98\\
2567.00	11.97\\
3090.00	12.03\\
3287.00	12.04\\
3496.00	12.00\\
3719.00	12.01\\
4209.00	11.99\\
5066.00	11.94\\
5733.00	11.99\\
10000.00	11.94\\
};
\addplot [color=white!70!black, dotted, line width=1.2pt, forget plot]
  table[row sep=crcr]{%
1.00	17.26\\
2.00	15.25\\
3.00	12.78\\
4.00	10.74\\
5.00	9.76\\
6.00	10.62\\
7.00	10.28\\
8.00	11.36\\
9.00	11.64\\
10.00	12.24\\
11.00	12.32\\
12.00	12.58\\
13.00	12.65\\
14.00	12.04\\
15.00	12.02\\
16.00	12.07\\
17.00	12.42\\
18.00	11.97\\
19.00	11.44\\
21.00	11.48\\
23.00	11.99\\
25.00	11.92\\
26.00	11.60\\
28.00	11.18\\
30.00	10.85\\
32.00	10.79\\
36.00	11.17\\
38.00	11.40\\
41.00	10.93\\
43.00	10.97\\
46.00	10.98\\
49.00	10.84\\
52.00	11.01\\
56.00	11.20\\
63.00	11.16\\
71.00	10.92\\
81.00	11.43\\
86.00	11.44\\
91.00	11.65\\
97.00	11.52\\
103.00	11.66\\
110.00	11.62\\
117.00	11.64\\
132.00	11.57\\
140.00	11.62\\
149.00	11.93\\
159.00	11.82\\
169.00	11.85\\
180.00	11.73\\
217.00	11.92\\
230.00	12.06\\
245.00	12.08\\
261.00	12.05\\
277.00	12.15\\
314.00	12.08\\
334.00	12.09\\
355.00	12.07\\
378.00	12.19\\
402.00	12.25\\
427.00	12.22\\
455.00	12.24\\
484.00	12.16\\
515.00	12.15\\
547.00	12.19\\
619.00	12.17\\
701.00	12.14\\
844.00	12.12\\
897.00	12.00\\
955.00	12.05\\
1016.00	12.06\\
1080.00	12.01\\
1149.00	12.00\\
1222.00	12.04\\
1300.00	11.99\\
1472.00	12.06\\
1665.00	11.93\\
1884.00	11.94\\
2268.00	11.92\\
2413.00	11.88\\
2730.00	11.97\\
3719.00	11.96\\
4209.00	11.96\\
5733.00	11.92\\
6901.00	11.88\\
7809.00	11.87\\
10000.00	11.92\\
};
\addplot [color=white!80!black, dashdotted, line width=1.2pt, forget plot]
  table[row sep=crcr]{%
1.00	12.55\\
2.00	8.84\\
3.00	11.10\\
4.00	11.89\\
5.00	13.61\\
6.00	12.36\\
7.00	12.85\\
8.00	13.10\\
9.00	13.42\\
10.00	13.46\\
11.00	12.54\\
12.00	12.47\\
13.00	11.65\\
14.00	11.50\\
15.00	10.81\\
16.00	11.50\\
17.00	11.31\\
18.00	11.55\\
19.00	11.49\\
21.00	11.47\\
22.00	11.68\\
23.00	11.51\\
25.00	11.83\\
26.00	11.78\\
28.00	11.79\\
30.00	11.59\\
32.00	11.65\\
34.00	11.31\\
36.00	11.85\\
38.00	11.68\\
41.00	11.60\\
43.00	11.67\\
46.00	11.82\\
49.00	11.83\\
52.00	11.68\\
56.00	11.78\\
59.00	11.74\\
63.00	11.93\\
71.00	12.11\\
76.00	12.37\\
81.00	12.07\\
86.00	12.14\\
91.00	12.27\\
97.00	12.30\\
103.00	12.23\\
110.00	12.30\\
117.00	12.26\\
124.00	12.39\\
132.00	12.26\\
159.00	12.29\\
169.00	12.11\\
180.00	12.13\\
191.00	12.05\\
204.00	12.11\\
217.00	12.09\\
245.00	11.93\\
261.00	11.89\\
277.00	12.01\\
295.00	12.02\\
314.00	11.96\\
355.00	11.95\\
378.00	12.16\\
402.00	12.00\\
427.00	12.10\\
455.00	12.11\\
582.00	11.98\\
619.00	12.00\\
701.00	11.97\\
746.00	11.86\\
844.00	11.75\\
897.00	11.88\\
955.00	11.85\\
1080.00	11.83\\
1300.00	11.96\\
1472.00	11.93\\
1665.00	11.98\\
1771.00	11.96\\
1884.00	12.00\\
3957.00	11.96\\
4477.00	11.99\\
5066.00	12.02\\
5733.00	11.97\\
6488.00	11.95\\
7341.00	11.95\\
8307.00	11.94\\
10000.00	11.96\\
};
\addplot [color=white!90!black, line width=1.2pt, forget plot]
  table[row sep=crcr]{%
1.00	7.72\\
2.00	9.15\\
3.00	9.95\\
5.00	10.54\\
7.00	12.75\\
8.00	11.67\\
9.00	11.51\\
10.00	11.56\\
11.00	10.53\\
12.00	10.91\\
13.00	11.00\\
14.00	10.47\\
15.00	10.39\\
16.00	10.64\\
17.00	10.72\\
18.00	11.14\\
19.00	11.28\\
21.00	11.33\\
22.00	11.56\\
23.00	11.72\\
25.00	11.93\\
26.00	11.76\\
28.00	11.79\\
30.00	12.20\\
36.00	11.85\\
38.00	12.08\\
41.00	11.91\\
43.00	11.53\\
46.00	11.26\\
49.00	11.22\\
52.00	11.05\\
56.00	11.07\\
59.00	11.24\\
63.00	11.00\\
67.00	10.88\\
71.00	10.73\\
76.00	10.88\\
81.00	11.10\\
86.00	11.07\\
91.00	11.34\\
97.00	11.30\\
117.00	11.35\\
124.00	11.22\\
132.00	11.33\\
140.00	11.28\\
149.00	11.27\\
159.00	11.46\\
169.00	11.56\\
180.00	11.61\\
191.00	11.76\\
204.00	11.82\\
217.00	11.95\\
230.00	12.01\\
245.00	11.92\\
261.00	11.86\\
277.00	11.75\\
334.00	11.86\\
355.00	11.66\\
378.00	11.67\\
402.00	11.56\\
455.00	11.63\\
547.00	11.85\\
582.00	11.94\\
619.00	11.92\\
701.00	12.00\\
746.00	11.92\\
897.00	11.89\\
955.00	11.93\\
1016.00	11.86\\
1080.00	11.94\\
1149.00	11.94\\
1222.00	12.00\\
1300.00	12.03\\
1665.00	11.88\\
1771.00	11.89\\
1884.00	11.87\\
3090.00	11.97\\
3496.00	12.01\\
3719.00	12.02\\
3957.00	11.99\\
4477.00	11.98\\
5389.00	11.93\\
5733.00	11.97\\
6099.00	11.91\\
8307.00	11.90\\
8837.00	11.91\\
9401.00	11.89\\
10000.00	11.92\\
};
\addplot [color=black, line width=1.2pt, forget plot]
  table[row sep=crcr]{%
0.00	12.00\\
10000.00	12.00\\
};
\end{axis}

\begin{axis}[%
width=0.411\figurewidth,
height=0.265\figureheight,
at={(0\figurewidth,0.368\figureheight)},
scale only axis,
xmin=0.00,
xmax=10000.00,
xtick={0.00,5000.00,10000.00},
xticklabels={\empty},
ymin=0.00,
ymax=40.00,
ylabel style={font=\color{mycolor1}},
ylabel={$\hat{\sym{thet}}_{\sym{idx_i}}$},
axis background/.style={fill=white},
title style={font=\bfseries\color{mycolor1}},
title={\textbf{$\sym{beta}_2$}},
xmajorgrids,
ymajorgrids,
grid style={dotted, opacity=0.25}
]
\addplot [color=white!10!black, line width=1.2pt, forget plot]
  table[row sep=crcr]{%
1.00	16.91\\
2.00	17.51\\
3.00	14.96\\
4.00	16.07\\
5.00	15.16\\
6.00	15.51\\
7.00	14.89\\
8.00	14.17\\
9.00	14.25\\
10.00	14.16\\
11.00	15.16\\
12.00	16.03\\
13.00	15.89\\
14.00	15.69\\
15.00	15.76\\
16.00	16.16\\
17.00	16.41\\
18.00	16.27\\
19.00	16.43\\
21.00	15.97\\
22.00	16.15\\
23.00	16.08\\
25.00	16.08\\
26.00	16.17\\
28.00	15.75\\
30.00	15.21\\
32.00	14.74\\
34.00	14.64\\
36.00	14.59\\
38.00	14.65\\
41.00	13.85\\
43.00	13.83\\
46.00	13.52\\
49.00	13.31\\
52.00	13.22\\
56.00	13.23\\
59.00	12.81\\
63.00	12.75\\
67.00	12.92\\
71.00	12.95\\
76.00	12.73\\
81.00	12.47\\
86.00	12.27\\
91.00	12.31\\
97.00	12.62\\
103.00	12.61\\
110.00	12.52\\
117.00	12.67\\
124.00	12.59\\
132.00	12.72\\
140.00	12.69\\
149.00	12.54\\
159.00	12.55\\
169.00	12.37\\
180.00	12.26\\
204.00	12.17\\
217.00	12.03\\
230.00	12.02\\
261.00	12.17\\
277.00	12.10\\
295.00	12.13\\
314.00	12.23\\
355.00	12.24\\
378.00	12.25\\
402.00	12.21\\
427.00	12.23\\
484.00	12.12\\
515.00	12.19\\
547.00	12.21\\
619.00	12.07\\
659.00	12.12\\
701.00	12.03\\
746.00	12.02\\
793.00	11.96\\
844.00	11.93\\
1080.00	12.10\\
1149.00	12.05\\
1383.00	12.06\\
1565.00	12.07\\
1665.00	12.05\\
1771.00	12.08\\
2005.00	12.05\\
2413.00	12.06\\
2567.00	12.10\\
2905.00	12.06\\
3287.00	12.11\\
3496.00	12.12\\
3719.00	12.09\\
4209.00	12.11\\
4477.00	12.08\\
4763.00	12.10\\
5066.00	12.06\\
5733.00	12.06\\
6099.00	12.06\\
6901.00	12.08\\
10000.00	12.10\\
};
\addplot [color=white!20!black, dashed, line width=1.2pt, forget plot]
  table[row sep=crcr]{%
1.00	5.02\\
2.00	11.14\\
3.00	12.22\\
4.00	13.19\\
5.00	12.84\\
6.00	13.78\\
7.00	13.37\\
8.00	12.92\\
10.00	14.60\\
11.00	14.64\\
12.00	14.54\\
13.00	14.38\\
14.00	13.83\\
15.00	14.00\\
16.00	13.57\\
18.00	12.96\\
19.00	12.80\\
21.00	12.87\\
22.00	13.02\\
23.00	12.98\\
25.00	12.35\\
26.00	12.25\\
28.00	12.34\\
30.00	12.58\\
32.00	12.27\\
36.00	12.54\\
38.00	12.76\\
41.00	12.79\\
43.00	12.61\\
46.00	12.88\\
49.00	12.85\\
52.00	12.93\\
56.00	12.66\\
59.00	12.57\\
63.00	12.30\\
67.00	12.36\\
71.00	12.09\\
76.00	12.17\\
81.00	12.22\\
86.00	12.24\\
97.00	12.38\\
103.00	12.30\\
110.00	12.42\\
117.00	12.48\\
124.00	12.28\\
132.00	12.15\\
140.00	11.99\\
149.00	11.94\\
169.00	11.96\\
180.00	12.15\\
191.00	12.12\\
204.00	12.29\\
217.00	12.35\\
261.00	11.99\\
277.00	12.05\\
295.00	11.97\\
314.00	11.97\\
334.00	11.94\\
378.00	12.04\\
427.00	12.07\\
484.00	12.19\\
547.00	12.16\\
582.00	12.13\\
659.00	12.13\\
746.00	11.99\\
793.00	12.13\\
897.00	12.25\\
1149.00	12.06\\
1222.00	11.98\\
1383.00	12.02\\
1565.00	12.02\\
1665.00	12.05\\
1771.00	12.00\\
1884.00	12.01\\
2413.00	12.18\\
2730.00	12.15\\
3719.00	11.99\\
4209.00	11.99\\
5733.00	11.96\\
10000.00	12.00\\
};
\addplot [color=white!30!black, dotted, line width=1.2pt, forget plot]
  table[row sep=crcr]{%
1.00	13.88\\
2.00	14.14\\
3.00	13.97\\
5.00	12.91\\
6.00	12.63\\
7.00	11.49\\
8.00	12.02\\
9.00	11.65\\
10.00	11.60\\
11.00	12.30\\
12.00	11.94\\
13.00	11.80\\
14.00	11.77\\
15.00	12.14\\
17.00	11.63\\
18.00	11.87\\
19.00	11.52\\
21.00	11.67\\
22.00	11.56\\
23.00	11.77\\
25.00	10.99\\
26.00	10.99\\
28.00	10.55\\
30.00	10.64\\
32.00	11.43\\
34.00	11.73\\
36.00	11.82\\
38.00	11.70\\
41.00	11.95\\
43.00	11.92\\
46.00	11.79\\
49.00	11.84\\
52.00	11.61\\
56.00	11.59\\
59.00	11.65\\
63.00	11.51\\
71.00	12.02\\
76.00	11.94\\
86.00	12.05\\
97.00	11.76\\
103.00	11.66\\
110.00	11.73\\
117.00	11.64\\
124.00	11.62\\
140.00	12.05\\
159.00	12.10\\
169.00	12.11\\
180.00	11.85\\
191.00	11.90\\
204.00	11.81\\
217.00	11.74\\
230.00	11.74\\
245.00	11.81\\
261.00	11.84\\
277.00	11.80\\
295.00	11.66\\
314.00	11.56\\
355.00	11.57\\
378.00	11.69\\
427.00	11.77\\
484.00	11.72\\
515.00	11.74\\
547.00	11.69\\
619.00	11.72\\
659.00	11.74\\
746.00	11.69\\
793.00	11.76\\
844.00	11.72\\
1080.00	11.78\\
1149.00	11.86\\
1300.00	11.93\\
1472.00	11.89\\
1771.00	11.86\\
1884.00	11.93\\
2132.00	11.98\\
2268.00	11.95\\
2413.00	11.98\\
2567.00	11.96\\
2730.00	12.00\\
2905.00	12.00\\
3090.00	12.04\\
3287.00	12.02\\
3496.00	12.05\\
3719.00	12.04\\
3957.00	12.09\\
4477.00	12.06\\
5066.00	12.05\\
5389.00	12.06\\
5733.00	12.02\\
6488.00	12.04\\
8837.00	12.02\\
10000.00	12.02\\
};
\addplot [color=white!40!black, dashdotted, line width=1.2pt, forget plot]
  table[row sep=crcr]{%
1.00	10.38\\
3.00	12.97\\
4.00	11.52\\
5.00	13.13\\
6.00	13.03\\
7.00	12.67\\
8.00	12.17\\
9.00	11.85\\
10.00	11.26\\
11.00	11.70\\
12.00	11.97\\
13.00	11.56\\
15.00	10.96\\
16.00	10.54\\
17.00	11.11\\
18.00	11.85\\
19.00	11.57\\
21.00	11.86\\
22.00	11.76\\
23.00	11.95\\
25.00	11.90\\
26.00	11.62\\
28.00	12.17\\
30.00	12.30\\
32.00	12.03\\
34.00	11.71\\
36.00	11.69\\
38.00	11.82\\
41.00	11.71\\
43.00	11.47\\
46.00	11.40\\
52.00	11.44\\
56.00	11.35\\
59.00	11.20\\
63.00	11.55\\
71.00	11.20\\
76.00	11.28\\
81.00	11.19\\
86.00	11.36\\
91.00	11.48\\
97.00	11.52\\
103.00	11.44\\
110.00	11.56\\
117.00	11.48\\
124.00	11.61\\
132.00	11.63\\
140.00	11.53\\
149.00	11.66\\
169.00	11.77\\
180.00	11.76\\
191.00	11.70\\
204.00	11.68\\
217.00	11.74\\
230.00	11.76\\
245.00	11.83\\
261.00	11.87\\
277.00	11.82\\
295.00	11.87\\
314.00	11.82\\
334.00	11.71\\
378.00	11.78\\
402.00	11.87\\
455.00	11.96\\
484.00	11.91\\
547.00	11.90\\
619.00	11.89\\
659.00	11.97\\
701.00	11.98\\
746.00	12.04\\
793.00	11.97\\
844.00	11.95\\
955.00	12.03\\
1016.00	12.02\\
1149.00	12.05\\
1383.00	11.98\\
1565.00	11.89\\
1771.00	11.87\\
2132.00	11.96\\
2268.00	12.02\\
2413.00	12.00\\
2567.00	12.05\\
2730.00	12.02\\
3090.00	12.03\\
3287.00	12.04\\
3496.00	12.00\\
4209.00	12.04\\
5066.00	12.11\\
10000.00	12.05\\
};
\addplot [color=gray, line width=1.2pt, forget plot]
  table[row sep=crcr]{%
1.00	10.38\\
2.00	9.62\\
3.00	12.22\\
4.00	12.93\\
5.00	14.40\\
6.00	15.00\\
7.00	13.16\\
8.00	12.61\\
9.00	11.44\\
10.00	10.80\\
11.00	11.08\\
12.00	10.60\\
13.00	11.36\\
14.00	11.43\\
15.00	11.38\\
16.00	11.84\\
17.00	12.02\\
18.00	12.14\\
19.00	11.96\\
21.00	11.49\\
23.00	11.67\\
25.00	11.95\\
26.00	11.81\\
28.00	11.96\\
30.00	11.84\\
32.00	11.81\\
34.00	11.88\\
36.00	11.85\\
38.00	12.19\\
41.00	12.50\\
43.00	12.38\\
46.00	12.43\\
49.00	12.74\\
56.00	12.75\\
59.00	12.68\\
67.00	12.41\\
71.00	12.70\\
76.00	12.38\\
81.00	12.40\\
86.00	12.37\\
91.00	12.26\\
103.00	12.49\\
124.00	12.42\\
132.00	12.50\\
140.00	12.23\\
149.00	12.13\\
159.00	12.19\\
169.00	11.95\\
180.00	11.92\\
191.00	12.00\\
204.00	12.01\\
217.00	12.06\\
230.00	12.15\\
245.00	12.29\\
261.00	12.28\\
277.00	12.37\\
295.00	12.24\\
314.00	12.17\\
334.00	12.20\\
355.00	12.16\\
378.00	12.03\\
402.00	12.08\\
427.00	12.01\\
455.00	12.06\\
515.00	12.05\\
547.00	11.98\\
659.00	11.94\\
746.00	11.75\\
844.00	11.69\\
955.00	11.70\\
1016.00	11.65\\
1080.00	11.73\\
1149.00	11.78\\
1222.00	11.78\\
1300.00	11.73\\
1383.00	11.80\\
1771.00	11.89\\
2005.00	11.90\\
2132.00	11.85\\
2268.00	11.86\\
2413.00	11.82\\
3287.00	11.84\\
4209.00	11.92\\
4477.00	11.90\\
5733.00	11.95\\
8307.00	11.93\\
10000.00	11.97\\
};
\addplot [color=white!60!black, dashed, line width=1.2pt, forget plot]
  table[row sep=crcr]{%
1.00	8.85\\
2.00	13.10\\
3.00	14.40\\
4.00	13.36\\
5.00	12.92\\
6.00	12.22\\
7.00	10.98\\
8.00	10.98\\
9.00	11.06\\
10.00	11.53\\
11.00	11.61\\
12.00	11.21\\
13.00	12.06\\
14.00	12.53\\
15.00	12.81\\
16.00	12.18\\
17.00	12.16\\
18.00	11.76\\
19.00	11.66\\
21.00	11.97\\
22.00	11.81\\
23.00	12.06\\
25.00	11.87\\
26.00	11.81\\
28.00	11.29\\
30.00	11.41\\
34.00	11.20\\
36.00	10.93\\
41.00	10.96\\
43.00	10.85\\
46.00	10.91\\
52.00	10.82\\
56.00	11.06\\
59.00	11.19\\
63.00	11.00\\
67.00	11.05\\
71.00	10.94\\
76.00	10.98\\
81.00	10.81\\
86.00	11.09\\
97.00	11.48\\
103.00	11.42\\
110.00	11.42\\
124.00	11.56\\
140.00	11.83\\
149.00	11.79\\
159.00	11.97\\
169.00	12.04\\
180.00	12.21\\
204.00	12.25\\
230.00	12.19\\
245.00	12.09\\
261.00	12.11\\
314.00	11.91\\
334.00	11.98\\
355.00	12.00\\
378.00	11.86\\
427.00	11.73\\
484.00	11.87\\
515.00	12.02\\
547.00	11.98\\
619.00	12.07\\
659.00	11.99\\
793.00	11.92\\
844.00	11.94\\
1016.00	11.90\\
1080.00	11.94\\
1149.00	11.92\\
1222.00	11.98\\
1383.00	11.98\\
1472.00	11.92\\
2005.00	12.06\\
2132.00	12.01\\
2268.00	12.01\\
2413.00	12.06\\
3090.00	12.10\\
3287.00	12.05\\
3957.00	12.04\\
4209.00	11.98\\
5066.00	12.02\\
5389.00	12.04\\
6099.00	12.03\\
6488.00	12.03\\
7341.00	12.06\\
10000.00	12.00\\
};
\addplot [color=white!70!black, dotted, line width=1.2pt, forget plot]
  table[row sep=crcr]{%
1.00	1.51\\
2.00	3.62\\
3.00	5.34\\
4.00	8.05\\
5.00	8.03\\
6.00	9.73\\
7.00	10.98\\
8.00	10.58\\
9.00	9.97\\
10.00	10.63\\
11.00	11.15\\
12.00	11.76\\
13.00	11.78\\
14.00	12.50\\
15.00	12.49\\
16.00	12.66\\
17.00	12.55\\
18.00	12.36\\
19.00	12.38\\
21.00	12.16\\
22.00	11.97\\
23.00	12.05\\
25.00	12.31\\
26.00	12.41\\
28.00	12.39\\
30.00	11.82\\
32.00	11.65\\
34.00	11.58\\
38.00	12.27\\
41.00	12.28\\
43.00	12.02\\
46.00	12.05\\
49.00	12.18\\
56.00	12.02\\
59.00	11.83\\
63.00	12.22\\
67.00	12.22\\
71.00	12.08\\
76.00	12.02\\
81.00	11.84\\
86.00	11.98\\
91.00	11.96\\
97.00	11.75\\
103.00	11.66\\
110.00	11.70\\
117.00	11.65\\
124.00	11.78\\
132.00	11.68\\
140.00	11.72\\
149.00	11.66\\
180.00	11.78\\
191.00	11.66\\
204.00	11.65\\
217.00	11.57\\
230.00	11.55\\
245.00	11.57\\
261.00	11.68\\
277.00	11.66\\
295.00	11.73\\
314.00	11.71\\
355.00	11.75\\
378.00	11.72\\
427.00	11.82\\
455.00	11.80\\
484.00	11.73\\
582.00	11.86\\
619.00	11.84\\
746.00	11.95\\
844.00	11.97\\
897.00	11.89\\
955.00	11.90\\
1016.00	11.84\\
1149.00	11.92\\
1222.00	11.91\\
1300.00	11.96\\
1383.00	11.92\\
1565.00	11.95\\
2005.00	11.95\\
2567.00	11.91\\
2730.00	11.95\\
3090.00	11.92\\
3719.00	11.99\\
3957.00	11.97\\
4477.00	11.98\\
4763.00	11.96\\
5389.00	12.02\\
5733.00	12.00\\
6901.00	12.03\\
10000.00	12.03\\
};
\addplot [color=white!80!black, dashdotted, line width=1.2pt, forget plot]
  table[row sep=crcr]{%
1.00	10.84\\
2.00	7.69\\
3.00	5.96\\
4.00	5.65\\
5.00	7.83\\
6.00	9.86\\
7.00	9.89\\
8.00	10.33\\
9.00	9.71\\
10.00	9.68\\
11.00	10.46\\
12.00	11.14\\
13.00	10.60\\
14.00	10.50\\
16.00	10.50\\
17.00	10.68\\
18.00	11.11\\
19.00	11.11\\
21.00	10.91\\
22.00	10.92\\
23.00	10.89\\
26.00	10.53\\
28.00	10.62\\
30.00	10.75\\
32.00	10.71\\
34.00	10.94\\
36.00	10.76\\
38.00	10.46\\
41.00	10.50\\
43.00	10.62\\
46.00	10.40\\
49.00	10.07\\
52.00	10.31\\
56.00	10.59\\
63.00	10.96\\
67.00	11.38\\
71.00	11.58\\
76.00	11.54\\
81.00	11.80\\
86.00	11.65\\
97.00	11.65\\
103.00	11.87\\
117.00	11.87\\
124.00	12.13\\
149.00	12.22\\
159.00	12.22\\
169.00	12.27\\
191.00	12.28\\
204.00	12.25\\
217.00	12.28\\
245.00	12.25\\
295.00	12.01\\
334.00	11.99\\
355.00	12.04\\
378.00	11.99\\
427.00	12.07\\
455.00	12.05\\
515.00	11.90\\
582.00	11.92\\
619.00	11.83\\
659.00	11.93\\
793.00	11.92\\
844.00	11.82\\
897.00	11.91\\
1016.00	11.88\\
1149.00	11.91\\
1383.00	11.91\\
1472.00	11.98\\
1565.00	11.99\\
1884.00	11.94\\
2005.00	11.94\\
2132.00	11.88\\
2567.00	11.89\\
2905.00	11.87\\
3090.00	11.89\\
3287.00	11.87\\
4477.00	11.95\\
4763.00	11.93\\
5733.00	11.99\\
6488.00	11.97\\
7341.00	11.95\\
8837.00	11.99\\
10000.00	11.98\\
};
\addplot [color=white!90!black, line width=1.2pt, forget plot]
  table[row sep=crcr]{%
1.00	13.68\\
2.00	12.85\\
4.00	12.96\\
5.00	13.25\\
7.00	13.55\\
8.00	14.07\\
9.00	13.56\\
10.00	13.35\\
11.00	13.06\\
12.00	13.48\\
13.00	13.02\\
14.00	13.12\\
16.00	12.49\\
17.00	12.56\\
18.00	12.77\\
19.00	12.31\\
21.00	13.00\\
22.00	12.80\\
23.00	12.94\\
25.00	13.13\\
26.00	13.12\\
28.00	13.06\\
30.00	12.72\\
32.00	12.81\\
34.00	12.76\\
36.00	13.26\\
38.00	13.15\\
41.00	13.17\\
43.00	13.22\\
46.00	13.12\\
49.00	12.84\\
52.00	13.07\\
56.00	13.15\\
59.00	13.27\\
63.00	13.34\\
67.00	13.14\\
71.00	13.02\\
76.00	12.94\\
81.00	13.07\\
86.00	12.85\\
91.00	12.73\\
97.00	12.88\\
103.00	12.80\\
110.00	12.60\\
117.00	12.49\\
124.00	12.47\\
132.00	12.24\\
140.00	12.26\\
149.00	12.39\\
159.00	12.38\\
169.00	12.43\\
180.00	12.42\\
191.00	12.16\\
204.00	12.20\\
217.00	12.31\\
230.00	12.30\\
245.00	12.39\\
295.00	12.48\\
314.00	12.45\\
334.00	12.48\\
355.00	12.47\\
378.00	12.52\\
402.00	12.46\\
455.00	12.37\\
484.00	12.30\\
515.00	12.19\\
582.00	12.19\\
619.00	12.23\\
701.00	12.23\\
746.00	12.11\\
793.00	12.04\\
897.00	12.10\\
955.00	12.06\\
1016.00	12.07\\
1080.00	12.02\\
1149.00	12.04\\
1222.00	11.98\\
1300.00	11.95\\
1472.00	12.08\\
1565.00	12.02\\
1771.00	12.01\\
2005.00	12.01\\
2268.00	11.96\\
2567.00	11.92\\
2730.00	11.98\\
2905.00	11.96\\
3090.00	11.98\\
3496.00	11.99\\
3957.00	12.04\\
4209.00	12.08\\
4763.00	12.11\\
5066.00	12.13\\
6901.00	12.06\\
8837.00	12.09\\
10000.00	12.05\\
};
\addplot [color=black, line width=1.2pt, forget plot]
  table[row sep=crcr]{%
0.00	12.00\\
10000.00	12.00\\
};
\end{axis}

\begin{axis}[%
width=0.411\figurewidth,
height=0.265\figureheight,
at={(0.54\figurewidth,0.368\figureheight)},
scale only axis,
xmin=0.00,
xmax=6.00,
xtick={0.00,2.00,4.00,6.00},
xticklabels={\empty},
ymin=0.00,
ymax=40.00,
ytick={0.00,10.00,20.00,30.00,40.00},
yticklabels={\empty},
axis background/.style={fill=white},
title style={font=\bfseries\color{mycolor1}},
title={\textbf{$\sym{beta}_{12}$}},
xmajorgrids,
ymajorgrids,
grid style={dotted, opacity=0.25}
]
\addplot [color=white!40!black, dashdotted, line width=1.2pt, forget plot]
  table[row sep=crcr]{%
1.00	1.21\\
2.00	2.28\\
3.00	1.31\\
4.00	-0.50\\
5.00	0.34\\
6.00	-0.52\\
};
\addplot [color=gray, line width=1.2pt, forget plot]
  table[row sep=crcr]{%
1.00	0.99\\
1.63	-4.00\\
};
\addplot [color=black, line width=1.2pt, forget plot]
  table[row sep=crcr]{%
0.00	-6.00\\
6.00	-6.00\\
};
\end{axis}

\begin{axis}[%
width=0.411\figurewidth,
height=0.265\figureheight,
at={(0\figurewidth,0\figureheight)},
scale only axis,
xmin=0.00,
xmax=1.00,
xlabel style={font=\color{mycolor1}},
xlabel={Iterationen $\sym{n}_{\sym{idx_MC}}$},
ymin=0.00,
ymax=40.00,
ylabel style={font=\color{mycolor1}},
ylabel={$\hat{\sym{thet}}_{\sym{idx_i}}$},
axis background/.style={fill=white},
title style={font=\bfseries\color{mycolor1}},
title={\textbf{$\sym{beta}_{11}$}},
xmajorgrids,
ymajorgrids,
grid style={dotted, opacity=0.25}
]
\addplot [color=black, line width=1.2pt, forget plot]
  table[row sep=crcr]{%
0.00	-6.00\\
1.00	-6.00\\
};
\end{axis}

\begin{axis}[%
width=0.411\figurewidth,
height=0.265\figureheight,
at={(0.54\figurewidth,0\figureheight)},
scale only axis,
xmin=0.00,
xmax=10000.00,
xlabel style={font=\color{mycolor1}},
xlabel={Iterationen $\sym{n}_{\sym{idx_MC}}$},
ymin=0.00,
ymax=40.00,
ytick={0.00,10.00,20.00,30.00,40.00},
yticklabels={\empty},
axis background/.style={fill=white},
title style={font=\bfseries\color{mycolor1}},
title={\textbf{$\sym{beta}_{22}$}},
xmajorgrids,
ymajorgrids,
grid style={dotted, opacity=0.25}
]
\addplot [color=black, line width=1.2pt, forget plot]
  table[row sep=crcr]{%
0.00	-6.00\\
10000.00	-6.00\\
};
\addplot [color=black, line width=1.2pt, forget plot]
  table[row sep=crcr]{%
0.00	24.00\\
10000.00	24.00\\
};
\addplot [color=black, line width=1.2pt, forget plot]
  table[row sep=crcr]{%
0.00	12.00\\
10000.00	12.00\\
};
\addplot [color=black, line width=1.2pt, forget plot]
  table[row sep=crcr]{%
0.00	12.00\\
10000.00	12.00\\
};
\addplot [color=black, line width=1.2pt, forget plot]
  table[row sep=crcr]{%
0.00	-6.00\\
6.00	-6.00\\
};
\addplot [color=black, line width=1.2pt, forget plot]
  table[row sep=crcr]{%
0.00	-6.00\\
1.00	-6.00\\
};
\end{axis}
\end{tikzpicture}%
    \caption{Konvergenz der \ac{OLS}-Parameterschätzwerte ($\hat{\sym{beta}}_i$). Alle Modelle erweisen sich unabhängig von der geometrischen Verzerrung als erwartungstreu. Die Schätzungen stabilisieren sich bei der $\sym{n}_{\sym{idx_MC}}=10^4$ reproduzierbar.}
    \label{fig:app.c_conv_beta}
\end{figure}

\begin{figure}[htbp]
    \centering
    \setlength\figurewidth{0.9\textwidth}
    \setlength\figureheight{16cm}
    % This file was created by matlab2tikz.
%
\definecolor{mycolor1}{rgb}{0.12941,0.12941,0.12941}%
%
\begin{tikzpicture}

\begin{axis}[%
width=0.411\figurewidth,
height=0.265\figureheight,
at={(0\figurewidth,0.735\figureheight)},
scale only axis,
xmin=0.00,
xmax=1.00,
xtick={0.00,0.50,1.00},
xticklabels={\empty},
ymin=0.00,
ymax=1.00,
ylabel style={font=\color{mycolor1}},
ylabel={Häufigkeit},
axis background/.style={fill=white},
title style={font=\bfseries\color{mycolor1}},
title={\textbf{$\sym{beta}_0$}},
xmajorgrids,
ymajorgrids,
grid style={dotted, opacity=0.25},
legend style={legend cell align=left, align=left}
]
\addplot[const plot, color=white!10!black, line width=1.2pt] table[row sep=crcr] {%
0.00	1.00\\
0.03	0.00\\
1.00	0.00\\
};
\addlegendentry{Basis \ac{CCD}}

\addplot[const plot, color=white!20!black, dashed, line width=1.2pt] table[row sep=crcr] {%
0.00	1.00\\
0.03	0.00\\
1.00	0.00\\
};
\addlegendentry{Fall I ($-20\,\%$)}

\addplot[const plot, color=white!30!black, dotted, line width=1.2pt] table[row sep=crcr] {%
0.00	1.00\\
0.03	0.00\\
0.05	0.00\\
0.08	0.00\\
1.00	0.00\\
};
\addlegendentry{Fall I ($+20\,\%$)}

\addplot[const plot, color=white!40!black, dashdotted, line width=1.2pt] table[row sep=crcr] {%
0.00	1.00\\
0.03	0.00\\
1.00	0.00\\
};
\addlegendentry{Fall I (\ac{FC})}

\addplot[const plot, color=gray, line width=1.2pt] table[row sep=crcr] {%
0.00	1.00\\
0.03	0.00\\
0.05	0.00\\
0.08	0.00\\
0.10	0.00\\
1.00	0.00\\
};
\addlegendentry{Fall II (leicht)}

\addplot[const plot, color=white!60!black, dashed, line width=1.2pt] table[row sep=crcr] {%
0.00	1.00\\
0.03	0.00\\
0.08	0.00\\
0.10	0.00\\
1.00	0.00\\
};
\addlegendentry{Fall II (stark)}

\addplot[const plot, color=white!70!black, dotted, line width=1.2pt] table[row sep=crcr] {%
0.00	1.00\\
0.03	0.00\\
0.05	0.00\\
0.64	0.00\\
0.67	0.00\\
1.00	0.00\\
};
\addlegendentry{Fall III}

\addplot[const plot, color=white!80!black, dashdotted, line width=1.2pt] table[row sep=crcr] {%
0.00	1.00\\
0.03	0.00\\
0.05	0.00\\
0.08	0.00\\
0.10	0.00\\
1.00	0.00\\
};
\addlegendentry{Fall IV: \ac{SP}}

\addplot[const plot, color=white!90!black, line width=1.2pt] table[row sep=crcr] {%
0.00	1.00\\
0.03	0.00\\
0.05	0.00\\
0.08	0.00\\
0.51	0.00\\
0.54	0.00\\
1.00	0.00\\
};
\addlegendentry{Fall IV: \ac{ZP}}

\addplot [color=black, line width=1.2pt]
  table[row sep=crcr]{%
0.05	0.00\\
0.05	1.00\\
};
\addlegendentry{Signifikanz $\sym{alpha}=0.05$}

\end{axis}

\begin{axis}[%
width=0.411\figurewidth,
height=0.265\figureheight,
at={(0.54\figurewidth,0.735\figureheight)},
scale only axis,
xmin=0.00,
xmax=1.00,
xtick={0.00,0.50,1.00},
xticklabels={\empty},
ymin=0.00,
ymax=1.00,
ytick={0.00,0.20,0.40,0.60,0.80,1.00},
yticklabels={\empty},
axis background/.style={fill=white},
title style={font=\bfseries\color{mycolor1}},
title={\textbf{$\sym{beta}_1$}},
xmajorgrids,
ymajorgrids,
grid style={dotted, opacity=0.25}
]
\addplot[const plot, color=white!10!black, line width=1.2pt, forget plot] table[row sep=crcr] {%
0.00	0.92\\
0.03	0.02\\
0.05	0.01\\
0.08	0.01\\
0.10	0.00\\
0.13	0.00\\
0.15	0.00\\
0.18	0.00\\
0.21	0.00\\
0.23	0.00\\
0.26	0.00\\
0.28	0.00\\
0.31	0.00\\
0.33	0.00\\
0.36	0.00\\
0.38	0.00\\
0.41	0.00\\
0.44	0.00\\
0.46	0.00\\
0.49	0.00\\
0.54	0.00\\
0.56	0.00\\
0.59	0.00\\
0.62	0.00\\
0.64	0.00\\
0.67	0.00\\
0.69	0.00\\
0.72	0.00\\
0.74	0.00\\
0.79	0.00\\
0.82	0.00\\
0.85	0.00\\
0.87	0.00\\
0.90	0.00\\
0.92	0.00\\
0.97	0.00\\
1.00	0.00\\
};
\addplot[const plot, color=white!20!black, dashed, line width=1.2pt, forget plot] table[row sep=crcr] {%
0.00	0.91\\
0.03	0.02\\
0.05	0.01\\
0.08	0.01\\
0.10	0.01\\
0.13	0.00\\
0.15	0.00\\
0.18	0.00\\
0.21	0.00\\
0.23	0.00\\
0.26	0.00\\
0.28	0.00\\
0.31	0.00\\
0.33	0.00\\
0.36	0.00\\
0.38	0.00\\
0.41	0.00\\
0.44	0.00\\
0.46	0.00\\
0.49	0.00\\
0.51	0.00\\
0.54	0.00\\
0.56	0.00\\
0.59	0.00\\
0.62	0.00\\
0.67	0.00\\
0.69	0.00\\
0.72	0.00\\
0.74	0.00\\
0.77	0.00\\
0.79	0.00\\
0.82	0.00\\
0.85	0.00\\
0.87	0.00\\
0.90	0.00\\
0.92	0.00\\
0.95	0.00\\
0.97	0.00\\
1.00	0.00\\
};
\addplot[const plot, color=white!30!black, dotted, line width=1.2pt, forget plot] table[row sep=crcr] {%
0.00	0.92\\
0.03	0.02\\
0.05	0.01\\
0.08	0.01\\
0.10	0.00\\
0.13	0.00\\
0.15	0.00\\
0.18	0.00\\
0.21	0.00\\
0.23	0.00\\
0.26	0.00\\
0.28	0.00\\
0.31	0.00\\
0.33	0.00\\
0.36	0.00\\
0.38	0.00\\
0.41	0.00\\
0.44	0.00\\
0.46	0.00\\
0.49	0.00\\
0.51	0.00\\
0.56	0.00\\
0.59	0.00\\
0.62	0.00\\
0.64	0.00\\
0.67	0.00\\
0.69	0.00\\
0.72	0.00\\
0.74	0.00\\
0.77	0.00\\
0.79	0.00\\
0.82	0.00\\
0.85	0.00\\
0.87	0.00\\
0.90	0.00\\
0.92	0.00\\
0.95	0.00\\
1.00	0.00\\
};
\addplot[const plot, color=white!40!black, dashdotted, line width=1.2pt, forget plot] table[row sep=crcr] {%
0.00	0.90\\
0.03	0.02\\
0.05	0.01\\
0.08	0.01\\
0.10	0.01\\
0.13	0.00\\
0.15	0.00\\
0.21	0.00\\
0.23	0.00\\
0.26	0.00\\
0.28	0.00\\
0.31	0.00\\
0.33	0.00\\
0.36	0.00\\
0.38	0.00\\
0.41	0.00\\
0.44	0.00\\
0.46	0.00\\
0.49	0.00\\
0.51	0.00\\
0.54	0.00\\
0.56	0.00\\
0.59	0.00\\
0.62	0.00\\
0.64	0.00\\
0.67	0.00\\
0.69	0.00\\
0.72	0.00\\
0.74	0.00\\
0.77	0.00\\
0.79	0.00\\
0.82	0.00\\
0.85	0.00\\
0.87	0.00\\
0.90	0.00\\
0.92	0.00\\
0.95	0.00\\
0.97	0.00\\
1.00	0.00\\
};
\addplot[const plot, color=gray, line width=1.2pt, forget plot] table[row sep=crcr] {%
0.00	0.92\\
0.03	0.02\\
0.05	0.01\\
0.08	0.01\\
0.10	0.01\\
0.13	0.00\\
0.15	0.00\\
0.18	0.00\\
0.21	0.00\\
0.23	0.00\\
0.26	0.00\\
0.31	0.00\\
0.33	0.00\\
0.38	0.00\\
0.41	0.00\\
0.44	0.00\\
0.46	0.00\\
0.49	0.00\\
0.51	0.00\\
0.54	0.00\\
0.56	0.00\\
0.59	0.00\\
0.62	0.00\\
0.64	0.00\\
0.69	0.00\\
0.72	0.00\\
0.74	0.00\\
0.77	0.00\\
0.79	0.00\\
0.82	0.00\\
0.85	0.00\\
0.87	0.00\\
0.90	0.00\\
0.92	0.00\\
0.95	0.00\\
0.97	0.00\\
1.00	0.00\\
};
\addplot[const plot, color=white!60!black, dashed, line width=1.2pt, forget plot] table[row sep=crcr] {%
0.00	0.91\\
0.03	0.02\\
0.05	0.01\\
0.08	0.01\\
0.10	0.01\\
0.13	0.00\\
0.15	0.00\\
0.18	0.00\\
0.21	0.00\\
0.26	0.00\\
0.28	0.00\\
0.31	0.00\\
0.33	0.00\\
0.36	0.00\\
0.38	0.00\\
0.41	0.00\\
0.44	0.00\\
0.46	0.00\\
0.49	0.00\\
0.51	0.00\\
0.54	0.00\\
0.56	0.00\\
0.59	0.00\\
0.62	0.00\\
0.64	0.00\\
0.67	0.00\\
0.69	0.00\\
0.72	0.00\\
0.74	0.00\\
0.77	0.00\\
0.79	0.00\\
0.82	0.00\\
0.85	0.00\\
0.87	0.00\\
0.90	0.00\\
0.92	0.00\\
0.95	0.00\\
0.97	0.00\\
1.00	0.00\\
};
\addplot[const plot, color=white!70!black, dotted, line width=1.2pt, forget plot] table[row sep=crcr] {%
0.00	0.92\\
0.03	0.02\\
0.05	0.01\\
0.08	0.01\\
0.10	0.01\\
0.13	0.00\\
0.15	0.00\\
0.18	0.00\\
0.21	0.00\\
0.23	0.00\\
0.26	0.00\\
0.28	0.00\\
0.33	0.00\\
0.36	0.00\\
0.38	0.00\\
0.41	0.00\\
0.44	0.00\\
0.46	0.00\\
0.49	0.00\\
0.51	0.00\\
0.54	0.00\\
0.56	0.00\\
0.59	0.00\\
0.62	0.00\\
0.64	0.00\\
0.67	0.00\\
0.69	0.00\\
0.72	0.00\\
0.74	0.00\\
0.77	0.00\\
0.79	0.00\\
0.85	0.00\\
0.87	0.00\\
0.92	0.00\\
0.95	0.00\\
0.97	0.00\\
1.00	0.00\\
};
\addplot[const plot, color=white!80!black, dashdotted, line width=1.2pt, forget plot] table[row sep=crcr] {%
0.00	0.86\\
0.03	0.04\\
0.05	0.02\\
0.08	0.01\\
0.10	0.01\\
0.13	0.01\\
0.15	0.01\\
0.18	0.01\\
0.21	0.00\\
0.23	0.00\\
0.26	0.00\\
0.28	0.00\\
0.31	0.00\\
0.33	0.00\\
0.36	0.00\\
0.38	0.00\\
0.41	0.00\\
0.44	0.00\\
0.46	0.00\\
0.49	0.00\\
0.51	0.00\\
0.54	0.00\\
0.56	0.00\\
0.59	0.00\\
0.62	0.00\\
0.64	0.00\\
0.67	0.00\\
0.69	0.00\\
0.72	0.00\\
0.79	0.00\\
0.82	0.00\\
0.85	0.00\\
0.87	0.00\\
0.90	0.00\\
0.92	0.00\\
0.97	0.00\\
1.00	0.00\\
};
\addplot[const plot, color=white!90!black, line width=1.2pt, forget plot] table[row sep=crcr] {%
0.00	0.91\\
0.03	0.02\\
0.05	0.01\\
0.08	0.01\\
0.10	0.01\\
0.13	0.00\\
0.15	0.00\\
0.18	0.00\\
0.21	0.00\\
0.23	0.00\\
0.26	0.00\\
0.28	0.00\\
0.31	0.00\\
0.33	0.00\\
0.36	0.00\\
0.38	0.00\\
0.41	0.00\\
0.46	0.00\\
0.49	0.00\\
0.51	0.00\\
0.54	0.00\\
0.59	0.00\\
0.62	0.00\\
0.64	0.00\\
0.69	0.00\\
0.72	0.00\\
0.74	0.00\\
0.77	0.00\\
0.79	0.00\\
0.82	0.00\\
0.85	0.00\\
0.87	0.00\\
0.90	0.00\\
0.92	0.00\\
0.95	0.00\\
0.97	0.00\\
1.00	0.00\\
};
\addplot [color=black, line width=1.2pt, forget plot]
  table[row sep=crcr]{%
0.05	0.00\\
0.05	1.00\\
};
\end{axis}

\begin{axis}[%
width=0.411\figurewidth,
height=0.265\figureheight,
at={(0\figurewidth,0.368\figureheight)},
scale only axis,
xmin=0.00,
xmax=1.00,
xtick={0.00,0.50,1.00},
xticklabels={\empty},
ymin=0.00,
ymax=1.00,
ylabel style={font=\color{mycolor1}},
ylabel={Häufigkeit},
axis background/.style={fill=white},
title style={font=\bfseries\color{mycolor1}},
title={\textbf{$\sym{beta}_2$}},
xmajorgrids,
ymajorgrids,
grid style={dotted, opacity=0.25}
]
\addplot[const plot, color=white!10!black, line width=1.2pt, forget plot] table[row sep=crcr] {%
0.00	0.92\\
0.03	0.02\\
0.05	0.01\\
0.08	0.01\\
0.10	0.00\\
0.13	0.00\\
0.18	0.00\\
0.21	0.00\\
0.23	0.00\\
0.26	0.00\\
0.28	0.00\\
0.31	0.00\\
0.36	0.00\\
0.38	0.00\\
0.41	0.00\\
0.46	0.00\\
0.51	0.00\\
0.54	0.00\\
0.56	0.00\\
0.59	0.00\\
0.62	0.00\\
0.64	0.00\\
0.67	0.00\\
0.69	0.00\\
0.72	0.00\\
0.74	0.00\\
0.77	0.00\\
0.79	0.00\\
0.82	0.00\\
0.85	0.00\\
0.87	0.00\\
0.90	0.00\\
0.92	0.00\\
0.95	0.00\\
0.97	0.00\\
1.00	0.00\\
};
\addplot[const plot, color=white!20!black, dashed, line width=1.2pt, forget plot] table[row sep=crcr] {%
0.00	0.92\\
0.03	0.02\\
0.05	0.01\\
0.08	0.01\\
0.10	0.00\\
0.13	0.00\\
0.15	0.00\\
0.18	0.00\\
0.21	0.00\\
0.26	0.00\\
0.28	0.00\\
0.31	0.00\\
0.33	0.00\\
0.36	0.00\\
0.38	0.00\\
0.41	0.00\\
0.44	0.00\\
0.46	0.00\\
0.49	0.00\\
0.51	0.00\\
0.54	0.00\\
0.56	0.00\\
0.59	0.00\\
0.62	0.00\\
0.64	0.00\\
0.67	0.00\\
0.69	0.00\\
0.72	0.00\\
0.74	0.00\\
0.79	0.00\\
0.82	0.00\\
0.87	0.00\\
0.90	0.00\\
0.92	0.00\\
0.95	0.00\\
1.00	0.00\\
};
\addplot[const plot, color=white!30!black, dotted, line width=1.2pt, forget plot] table[row sep=crcr] {%
0.00	0.92\\
0.03	0.02\\
0.05	0.01\\
0.08	0.01\\
0.10	0.00\\
0.13	0.00\\
0.15	0.00\\
0.18	0.00\\
0.21	0.00\\
0.23	0.00\\
0.28	0.00\\
0.31	0.00\\
0.33	0.00\\
0.36	0.00\\
0.38	0.00\\
0.41	0.00\\
0.44	0.00\\
0.46	0.00\\
0.49	0.00\\
0.51	0.00\\
0.54	0.00\\
0.56	0.00\\
0.59	0.00\\
0.62	0.00\\
0.67	0.00\\
0.69	0.00\\
0.72	0.00\\
0.77	0.00\\
0.79	0.00\\
0.82	0.00\\
0.85	0.00\\
0.87	0.00\\
0.90	0.00\\
0.92	0.00\\
0.97	0.00\\
1.00	0.00\\
};
\addplot[const plot, color=white!40!black, dashdotted, line width=1.2pt, forget plot] table[row sep=crcr] {%
0.00	0.92\\
0.03	0.02\\
0.05	0.01\\
0.08	0.01\\
0.10	0.00\\
0.13	0.00\\
0.15	0.00\\
0.18	0.00\\
0.23	0.00\\
0.26	0.00\\
0.28	0.00\\
0.31	0.00\\
0.33	0.00\\
0.36	0.00\\
0.38	0.00\\
0.41	0.00\\
0.44	0.00\\
0.46	0.00\\
0.49	0.00\\
0.51	0.00\\
0.54	0.00\\
0.56	0.00\\
0.59	0.00\\
0.62	0.00\\
0.64	0.00\\
0.67	0.00\\
0.69	0.00\\
0.72	0.00\\
0.74	0.00\\
0.77	0.00\\
0.79	0.00\\
0.82	0.00\\
0.85	0.00\\
0.90	0.00\\
0.92	0.00\\
0.95	0.00\\
0.97	0.00\\
1.00	0.00\\
};
\addplot[const plot, color=gray, line width=1.2pt, forget plot] table[row sep=crcr] {%
0.00	0.91\\
0.03	0.02\\
0.05	0.01\\
0.08	0.01\\
0.10	0.01\\
0.13	0.00\\
0.15	0.00\\
0.18	0.00\\
0.21	0.00\\
0.23	0.00\\
0.26	0.00\\
0.28	0.00\\
0.31	0.00\\
0.33	0.00\\
0.36	0.00\\
0.38	0.00\\
0.41	0.00\\
0.44	0.00\\
0.46	0.00\\
0.49	0.00\\
0.51	0.00\\
0.56	0.00\\
0.59	0.00\\
0.64	0.00\\
0.67	0.00\\
0.69	0.00\\
0.72	0.00\\
0.77	0.00\\
0.79	0.00\\
0.82	0.00\\
0.85	0.00\\
0.87	0.00\\
0.92	0.00\\
0.95	0.00\\
0.97	0.00\\
1.00	0.00\\
};
\addplot[const plot, color=white!60!black, dashed, line width=1.2pt, forget plot] table[row sep=crcr] {%
0.00	0.92\\
0.03	0.02\\
0.05	0.01\\
0.08	0.01\\
0.10	0.00\\
0.13	0.00\\
0.15	0.00\\
0.18	0.00\\
0.21	0.00\\
0.23	0.00\\
0.26	0.00\\
0.28	0.00\\
0.31	0.00\\
0.33	0.00\\
0.36	0.00\\
0.38	0.00\\
0.41	0.00\\
0.44	0.00\\
0.46	0.00\\
0.49	0.00\\
0.51	0.00\\
0.54	0.00\\
0.56	0.00\\
0.59	0.00\\
0.62	0.00\\
0.67	0.00\\
0.69	0.00\\
0.72	0.00\\
0.74	0.00\\
0.77	0.00\\
0.79	0.00\\
0.82	0.00\\
0.85	0.00\\
0.87	0.00\\
0.92	0.00\\
0.95	0.00\\
0.97	0.00\\
1.00	0.00\\
};
\addplot[const plot, color=white!70!black, dotted, line width=1.2pt, forget plot] table[row sep=crcr] {%
0.00	0.93\\
0.03	0.02\\
0.05	0.01\\
0.08	0.01\\
0.10	0.00\\
0.13	0.00\\
0.15	0.00\\
0.18	0.00\\
0.21	0.00\\
0.23	0.00\\
0.26	0.00\\
0.28	0.00\\
0.31	0.00\\
0.36	0.00\\
0.41	0.00\\
0.44	0.00\\
0.49	0.00\\
0.51	0.00\\
0.54	0.00\\
0.56	0.00\\
0.59	0.00\\
0.62	0.00\\
0.64	0.00\\
0.67	0.00\\
0.69	0.00\\
0.72	0.00\\
0.74	0.00\\
0.77	0.00\\
0.79	0.00\\
0.82	0.00\\
0.85	0.00\\
0.87	0.00\\
0.90	0.00\\
0.92	0.00\\
0.95	0.00\\
0.97	0.00\\
1.00	0.00\\
};
\addplot[const plot, color=white!80!black, dashdotted, line width=1.2pt, forget plot] table[row sep=crcr] {%
0.00	0.92\\
0.03	0.02\\
0.05	0.01\\
0.08	0.01\\
0.10	0.00\\
0.13	0.00\\
0.15	0.00\\
0.18	0.00\\
0.21	0.00\\
0.23	0.00\\
0.26	0.00\\
0.28	0.00\\
0.31	0.00\\
0.33	0.00\\
0.36	0.00\\
0.38	0.00\\
0.41	0.00\\
0.44	0.00\\
0.46	0.00\\
0.51	0.00\\
0.56	0.00\\
0.59	0.00\\
0.62	0.00\\
0.64	0.00\\
0.67	0.00\\
0.69	0.00\\
0.72	0.00\\
0.74	0.00\\
0.77	0.00\\
0.82	0.00\\
0.85	0.00\\
0.87	0.00\\
0.90	0.00\\
0.92	0.00\\
0.95	0.00\\
0.97	0.00\\
1.00	0.00\\
};
\addplot[const plot, color=white!90!black, line width=1.2pt, forget plot] table[row sep=crcr] {%
0.00	0.92\\
0.03	0.02\\
0.05	0.01\\
0.08	0.01\\
0.10	0.00\\
0.13	0.00\\
0.15	0.00\\
0.18	0.00\\
0.21	0.00\\
0.23	0.00\\
0.26	0.00\\
0.28	0.00\\
0.31	0.00\\
0.33	0.00\\
0.38	0.00\\
0.41	0.00\\
0.44	0.00\\
0.46	0.00\\
0.49	0.00\\
0.51	0.00\\
0.54	0.00\\
0.56	0.00\\
0.59	0.00\\
0.62	0.00\\
0.64	0.00\\
0.69	0.00\\
0.72	0.00\\
0.74	0.00\\
0.77	0.00\\
0.79	0.00\\
0.82	0.00\\
0.85	0.00\\
0.87	0.00\\
0.90	0.00\\
0.92	0.00\\
0.95	0.00\\
0.97	0.00\\
1.00	0.00\\
};
\addplot [color=black, line width=1.2pt, forget plot]
  table[row sep=crcr]{%
0.05	0.00\\
0.05	1.00\\
};
\end{axis}

\begin{axis}[%
width=0.411\figurewidth,
height=0.265\figureheight,
at={(0.54\figurewidth,0.368\figureheight)},
scale only axis,
xmin=0.00,
xmax=1.00,
xtick={0.00,0.50,1.00},
xticklabels={\empty},
ymin=0.00,
ymax=1.00,
ytick={0.00,0.20,0.40,0.60,0.80,1.00},
yticklabels={\empty},
axis background/.style={fill=white},
title style={font=\bfseries\color{mycolor1}},
title={\textbf{$\sym{beta}_{12}$}},
xmajorgrids,
ymajorgrids,
grid style={dotted, opacity=0.25}
]
\addplot[const plot, color=white!10!black, line width=1.2pt, forget plot] table[row sep=crcr] {%
0.00	0.49\\
0.03	0.08\\
0.05	0.05\\
0.08	0.04\\
0.10	0.02\\
0.13	0.02\\
0.15	0.02\\
0.18	0.02\\
0.21	0.02\\
0.23	0.01\\
0.26	0.01\\
0.28	0.01\\
0.31	0.01\\
0.33	0.01\\
0.36	0.01\\
0.38	0.01\\
0.41	0.01\\
0.44	0.01\\
0.46	0.01\\
0.49	0.01\\
0.51	0.01\\
0.54	0.01\\
0.56	0.01\\
0.59	0.01\\
0.62	0.01\\
0.64	0.01\\
0.67	0.01\\
0.69	0.01\\
0.72	0.01\\
0.74	0.01\\
0.77	0.01\\
0.79	0.01\\
0.82	0.01\\
0.85	0.01\\
0.87	0.01\\
0.90	0.01\\
0.92	0.01\\
0.95	0.01\\
0.97	0.01\\
1.00	0.01\\
};
\addplot[const plot, color=white!20!black, dashed, line width=1.2pt, forget plot] table[row sep=crcr] {%
0.00	0.49\\
0.03	0.08\\
0.05	0.05\\
0.08	0.03\\
0.10	0.02\\
0.13	0.02\\
0.15	0.02\\
0.18	0.02\\
0.21	0.01\\
0.23	0.01\\
0.26	0.01\\
0.28	0.01\\
0.31	0.01\\
0.33	0.01\\
0.36	0.01\\
0.38	0.01\\
0.41	0.01\\
0.44	0.01\\
0.46	0.01\\
0.49	0.01\\
0.51	0.01\\
0.54	0.01\\
0.56	0.01\\
0.59	0.01\\
0.62	0.01\\
0.64	0.01\\
0.67	0.01\\
0.69	0.00\\
0.72	0.01\\
0.74	0.01\\
0.77	0.01\\
0.79	0.01\\
0.82	0.01\\
0.85	0.01\\
0.87	0.01\\
0.90	0.01\\
0.92	0.01\\
0.95	0.01\\
0.97	0.01\\
1.00	0.01\\
};
\addplot[const plot, color=white!30!black, dotted, line width=1.2pt, forget plot] table[row sep=crcr] {%
0.00	0.50\\
0.03	0.08\\
0.05	0.04\\
0.08	0.03\\
0.10	0.02\\
0.13	0.02\\
0.15	0.02\\
0.18	0.02\\
0.21	0.02\\
0.23	0.01\\
0.26	0.01\\
0.28	0.01\\
0.31	0.01\\
0.33	0.01\\
0.36	0.01\\
0.38	0.01\\
0.41	0.01\\
0.44	0.01\\
0.46	0.01\\
0.49	0.01\\
0.51	0.01\\
0.54	0.01\\
0.56	0.01\\
0.62	0.01\\
0.64	0.01\\
0.67	0.01\\
0.69	0.01\\
0.74	0.01\\
0.77	0.01\\
0.79	0.01\\
0.82	0.01\\
0.85	0.01\\
0.87	0.01\\
0.90	0.01\\
0.92	0.01\\
0.95	0.01\\
0.97	0.01\\
1.00	0.01\\
};
\addplot[const plot, color=white!40!black, dashdotted, line width=1.2pt, forget plot] table[row sep=crcr] {%
0.00	0.49\\
0.03	0.08\\
0.05	0.05\\
0.08	0.03\\
0.10	0.03\\
0.13	0.02\\
0.15	0.02\\
0.18	0.02\\
0.21	0.02\\
0.23	0.01\\
0.26	0.01\\
0.28	0.01\\
0.31	0.01\\
0.33	0.01\\
0.36	0.01\\
0.38	0.01\\
0.41	0.01\\
0.44	0.01\\
0.46	0.01\\
0.49	0.01\\
0.51	0.01\\
0.54	0.01\\
0.56	0.01\\
0.59	0.01\\
0.62	0.01\\
0.64	0.01\\
0.67	0.01\\
0.69	0.01\\
0.72	0.01\\
0.77	0.01\\
0.79	0.01\\
0.82	0.01\\
0.85	0.01\\
0.87	0.01\\
0.90	0.01\\
0.92	0.00\\
0.95	0.01\\
0.97	0.01\\
1.00	0.01\\
};
\addplot[const plot, color=gray, line width=1.2pt, forget plot] table[row sep=crcr] {%
0.00	0.49\\
0.03	0.08\\
0.05	0.04\\
0.08	0.03\\
0.10	0.03\\
0.13	0.03\\
0.15	0.02\\
0.18	0.02\\
0.21	0.02\\
0.23	0.01\\
0.26	0.01\\
0.28	0.01\\
0.31	0.01\\
0.33	0.01\\
0.36	0.01\\
0.38	0.01\\
0.41	0.01\\
0.44	0.01\\
0.46	0.01\\
0.49	0.01\\
0.51	0.01\\
0.54	0.01\\
0.56	0.01\\
0.59	0.01\\
0.62	0.01\\
0.64	0.01\\
0.67	0.01\\
0.69	0.01\\
0.72	0.01\\
0.74	0.01\\
0.77	0.01\\
0.79	0.01\\
0.82	0.01\\
0.85	0.01\\
0.87	0.01\\
0.90	0.01\\
0.95	0.00\\
0.97	0.01\\
1.00	0.01\\
};
\addplot[const plot, color=white!60!black, dashed, line width=1.2pt, forget plot] table[row sep=crcr] {%
0.00	0.54\\
0.03	0.07\\
0.05	0.05\\
0.08	0.03\\
0.10	0.02\\
0.13	0.02\\
0.15	0.02\\
0.18	0.01\\
0.21	0.02\\
0.23	0.01\\
0.26	0.01\\
0.28	0.01\\
0.31	0.01\\
0.33	0.01\\
0.36	0.01\\
0.38	0.01\\
0.41	0.01\\
0.44	0.01\\
0.46	0.01\\
0.49	0.01\\
0.54	0.01\\
0.56	0.01\\
0.59	0.01\\
0.62	0.01\\
0.64	0.01\\
0.67	0.01\\
0.69	0.01\\
0.72	0.01\\
0.74	0.01\\
0.77	0.01\\
0.79	0.01\\
0.82	0.01\\
0.85	0.01\\
0.87	0.01\\
0.90	0.01\\
0.92	0.01\\
0.95	0.01\\
0.97	0.01\\
1.00	0.01\\
};
\addplot[const plot, color=white!70!black, dotted, line width=1.2pt, forget plot] table[row sep=crcr] {%
0.00	0.50\\
0.03	0.07\\
0.05	0.05\\
0.08	0.04\\
0.10	0.03\\
0.13	0.02\\
0.15	0.02\\
0.18	0.02\\
0.21	0.01\\
0.23	0.01\\
0.26	0.01\\
0.28	0.01\\
0.31	0.01\\
0.33	0.01\\
0.36	0.01\\
0.38	0.01\\
0.41	0.01\\
0.44	0.01\\
0.46	0.01\\
0.49	0.01\\
0.51	0.01\\
0.54	0.01\\
0.56	0.01\\
0.59	0.01\\
0.62	0.01\\
0.64	0.01\\
0.67	0.01\\
0.69	0.01\\
0.72	0.01\\
0.74	0.01\\
0.77	0.01\\
0.79	0.01\\
0.82	0.01\\
0.85	0.01\\
0.87	0.01\\
0.90	0.01\\
0.92	0.01\\
0.95	0.01\\
0.97	0.01\\
1.00	0.01\\
};
\addplot[const plot, color=white!80!black, dashdotted, line width=1.2pt, forget plot] table[row sep=crcr] {%
0.00	0.49\\
0.03	0.08\\
0.05	0.05\\
0.08	0.03\\
0.10	0.03\\
0.13	0.02\\
0.15	0.02\\
0.18	0.02\\
0.21	0.02\\
0.23	0.01\\
0.26	0.01\\
0.28	0.01\\
0.31	0.01\\
0.33	0.01\\
0.36	0.01\\
0.38	0.01\\
0.41	0.01\\
0.44	0.01\\
0.46	0.01\\
0.49	0.01\\
0.51	0.01\\
0.54	0.01\\
0.56	0.01\\
0.59	0.01\\
0.62	0.01\\
0.64	0.01\\
0.67	0.01\\
0.69	0.01\\
0.72	0.01\\
0.74	0.01\\
0.77	0.01\\
0.79	0.01\\
0.82	0.01\\
0.85	0.01\\
0.87	0.01\\
0.90	0.01\\
0.92	0.01\\
0.95	0.01\\
0.97	0.01\\
1.00	0.01\\
};
\addplot[const plot, color=white!90!black, line width=1.2pt, forget plot] table[row sep=crcr] {%
0.00	0.49\\
0.03	0.08\\
0.05	0.05\\
0.08	0.03\\
0.10	0.03\\
0.13	0.02\\
0.15	0.02\\
0.18	0.02\\
0.21	0.02\\
0.23	0.01\\
0.26	0.01\\
0.28	0.01\\
0.31	0.01\\
0.33	0.01\\
0.36	0.01\\
0.38	0.01\\
0.41	0.01\\
0.44	0.01\\
0.46	0.01\\
0.49	0.01\\
0.51	0.01\\
0.54	0.01\\
0.56	0.01\\
0.59	0.01\\
0.62	0.01\\
0.64	0.01\\
0.67	0.01\\
0.69	0.01\\
0.72	0.01\\
0.74	0.01\\
0.79	0.01\\
0.82	0.01\\
0.85	0.01\\
0.87	0.01\\
0.90	0.01\\
0.92	0.01\\
0.95	0.01\\
0.97	0.01\\
1.00	0.01\\
};
\addplot [color=black, line width=1.2pt, forget plot]
  table[row sep=crcr]{%
0.05	0.00\\
0.05	1.00\\
};
\end{axis}

\begin{axis}[%
width=0.411\figurewidth,
height=0.265\figureheight,
at={(0\figurewidth,0\figureheight)},
scale only axis,
xmin=0.00,
xmax=1.00,
xlabel style={font=\color{mycolor1}},
xlabel={$\sym{p-Wert}$},
ymin=0.00,
ymax=1.00,
ylabel style={font=\color{mycolor1}},
ylabel={Häufigkeit},
axis background/.style={fill=white},
title style={font=\bfseries\color{mycolor1}},
title={\textbf{$\sym{beta}_{11}$}},
xmajorgrids,
ymajorgrids,
grid style={dotted, opacity=0.25}
]
\addplot[const plot, color=white!10!black, line width=1.2pt, forget plot] table[row sep=crcr] {%
0.00	0.63\\
0.03	0.06\\
0.05	0.03\\
0.08	0.02\\
0.10	0.02\\
0.13	0.02\\
0.15	0.01\\
0.18	0.01\\
0.21	0.01\\
0.23	0.01\\
0.26	0.01\\
0.28	0.01\\
0.31	0.01\\
0.33	0.01\\
0.36	0.01\\
0.38	0.01\\
0.41	0.01\\
0.44	0.01\\
0.46	0.01\\
0.49	0.01\\
0.51	0.01\\
0.54	0.01\\
0.56	0.01\\
0.59	0.01\\
0.62	0.01\\
0.64	0.01\\
0.67	0.01\\
0.69	0.01\\
0.72	0.00\\
0.74	0.00\\
0.77	0.01\\
0.79	0.00\\
0.82	0.01\\
0.85	0.00\\
0.87	0.00\\
0.90	0.00\\
0.92	0.00\\
0.97	0.00\\
1.00	0.00\\
};
\addplot[const plot, color=white!20!black, dashed, line width=1.2pt, forget plot] table[row sep=crcr] {%
0.00	0.58\\
0.03	0.06\\
0.05	0.04\\
0.08	0.03\\
0.10	0.02\\
0.13	0.02\\
0.15	0.02\\
0.18	0.01\\
0.21	0.01\\
0.23	0.01\\
0.26	0.01\\
0.28	0.01\\
0.31	0.01\\
0.33	0.01\\
0.36	0.01\\
0.38	0.01\\
0.41	0.01\\
0.44	0.01\\
0.49	0.01\\
0.51	0.01\\
0.56	0.01\\
0.59	0.01\\
0.62	0.01\\
0.64	0.01\\
0.67	0.00\\
0.69	0.01\\
0.72	0.01\\
0.74	0.00\\
0.77	0.01\\
0.79	0.01\\
0.82	0.01\\
0.85	0.01\\
0.87	0.01\\
0.90	0.00\\
0.92	0.01\\
0.95	0.00\\
0.97	0.00\\
1.00	0.00\\
};
\addplot[const plot, color=white!30!black, dotted, line width=1.2pt, forget plot] table[row sep=crcr] {%
0.00	0.70\\
0.03	0.05\\
0.05	0.03\\
0.08	0.02\\
0.10	0.02\\
0.13	0.02\\
0.15	0.01\\
0.18	0.01\\
0.21	0.01\\
0.23	0.01\\
0.26	0.01\\
0.28	0.01\\
0.31	0.01\\
0.33	0.01\\
0.36	0.01\\
0.38	0.01\\
0.41	0.01\\
0.44	0.01\\
0.49	0.00\\
0.51	0.00\\
0.54	0.01\\
0.56	0.00\\
0.59	0.00\\
0.62	0.00\\
0.64	0.01\\
0.67	0.00\\
0.69	0.00\\
0.72	0.00\\
0.74	0.00\\
0.77	0.00\\
0.79	0.00\\
0.82	0.00\\
0.85	0.00\\
0.87	0.00\\
0.90	0.00\\
0.92	0.00\\
0.95	0.00\\
0.97	0.00\\
1.00	0.00\\
};
\addplot[const plot, color=white!40!black, dashdotted, line width=1.2pt, forget plot] table[row sep=crcr] {%
0.00	0.55\\
0.03	0.07\\
0.05	0.04\\
0.08	0.03\\
0.10	0.02\\
0.13	0.02\\
0.15	0.02\\
0.18	0.01\\
0.21	0.01\\
0.23	0.01\\
0.26	0.01\\
0.28	0.01\\
0.31	0.01\\
0.33	0.01\\
0.36	0.01\\
0.38	0.01\\
0.41	0.01\\
0.44	0.01\\
0.46	0.01\\
0.49	0.01\\
0.51	0.01\\
0.56	0.01\\
0.59	0.01\\
0.62	0.01\\
0.64	0.01\\
0.67	0.01\\
0.69	0.01\\
0.72	0.01\\
0.74	0.01\\
0.77	0.01\\
0.79	0.01\\
0.82	0.01\\
0.85	0.00\\
0.87	0.01\\
0.90	0.01\\
0.92	0.01\\
0.95	0.01\\
0.97	0.01\\
1.00	0.01\\
};
\addplot[const plot, color=gray, line width=1.2pt, forget plot] table[row sep=crcr] {%
0.00	0.64\\
0.03	0.06\\
0.05	0.04\\
0.08	0.03\\
0.10	0.02\\
0.13	0.02\\
0.15	0.01\\
0.18	0.01\\
0.21	0.01\\
0.23	0.01\\
0.26	0.01\\
0.28	0.01\\
0.31	0.01\\
0.33	0.01\\
0.36	0.01\\
0.38	0.01\\
0.41	0.00\\
0.44	0.01\\
0.46	0.01\\
0.49	0.01\\
0.51	0.01\\
0.54	0.01\\
0.56	0.01\\
0.59	0.01\\
0.62	0.01\\
0.64	0.00\\
0.67	0.01\\
0.69	0.00\\
0.72	0.01\\
0.74	0.01\\
0.77	0.00\\
0.79	0.00\\
0.82	0.00\\
0.87	0.00\\
0.90	0.01\\
0.92	0.00\\
0.95	0.00\\
0.97	0.00\\
1.00	0.00\\
};
\addplot[const plot, color=white!60!black, dashed, line width=1.2pt, forget plot] table[row sep=crcr] {%
0.00	0.61\\
0.03	0.06\\
0.05	0.04\\
0.08	0.03\\
0.10	0.02\\
0.13	0.02\\
0.15	0.01\\
0.18	0.01\\
0.21	0.01\\
0.23	0.01\\
0.26	0.01\\
0.28	0.01\\
0.31	0.01\\
0.33	0.01\\
0.36	0.01\\
0.38	0.01\\
0.41	0.01\\
0.44	0.01\\
0.46	0.01\\
0.49	0.01\\
0.51	0.01\\
0.54	0.00\\
0.56	0.01\\
0.59	0.01\\
0.62	0.00\\
0.64	0.01\\
0.67	0.00\\
0.69	0.00\\
0.72	0.00\\
0.74	0.01\\
0.77	0.01\\
0.79	0.01\\
0.82	0.00\\
0.85	0.01\\
0.87	0.01\\
0.90	0.00\\
0.92	0.01\\
0.95	0.00\\
0.97	0.00\\
1.00	0.00\\
};
\addplot[const plot, color=white!70!black, dotted, line width=1.2pt, forget plot] table[row sep=crcr] {%
0.00	0.62\\
0.03	0.06\\
0.05	0.04\\
0.08	0.03\\
0.10	0.02\\
0.13	0.02\\
0.15	0.01\\
0.18	0.01\\
0.21	0.01\\
0.23	0.01\\
0.26	0.01\\
0.28	0.01\\
0.31	0.01\\
0.33	0.01\\
0.36	0.01\\
0.38	0.01\\
0.41	0.01\\
0.44	0.01\\
0.46	0.01\\
0.49	0.01\\
0.54	0.01\\
0.56	0.01\\
0.59	0.01\\
0.62	0.01\\
0.64	0.01\\
0.67	0.00\\
0.69	0.00\\
0.72	0.00\\
0.74	0.00\\
0.77	0.00\\
0.82	0.00\\
0.85	0.01\\
0.87	0.01\\
0.90	0.00\\
0.92	0.01\\
0.95	0.01\\
0.97	0.00\\
1.00	0.00\\
};
\addplot[const plot, color=white!80!black, dashdotted, line width=1.2pt, forget plot] table[row sep=crcr] {%
0.00	0.51\\
0.03	0.08\\
0.05	0.05\\
0.08	0.04\\
0.10	0.02\\
0.13	0.02\\
0.15	0.02\\
0.18	0.02\\
0.21	0.01\\
0.23	0.01\\
0.26	0.01\\
0.28	0.01\\
0.31	0.01\\
0.33	0.01\\
0.36	0.01\\
0.38	0.01\\
0.44	0.01\\
0.46	0.01\\
0.49	0.01\\
0.51	0.01\\
0.54	0.01\\
0.56	0.01\\
0.59	0.01\\
0.62	0.01\\
0.64	0.01\\
0.67	0.01\\
0.69	0.01\\
0.72	0.01\\
0.74	0.01\\
0.77	0.00\\
0.79	0.01\\
0.82	0.01\\
0.85	0.01\\
0.87	0.01\\
0.90	0.01\\
0.92	0.01\\
0.95	0.01\\
0.97	0.01\\
1.00	0.01\\
};
\addplot[const plot, color=white!90!black, line width=1.2pt, forget plot] table[row sep=crcr] {%
0.00	0.62\\
0.03	0.06\\
0.05	0.04\\
0.08	0.02\\
0.10	0.02\\
0.13	0.01\\
0.15	0.01\\
0.18	0.01\\
0.21	0.01\\
0.23	0.01\\
0.26	0.01\\
0.28	0.01\\
0.31	0.01\\
0.36	0.01\\
0.38	0.01\\
0.41	0.01\\
0.44	0.01\\
0.46	0.00\\
0.49	0.01\\
0.51	0.01\\
0.54	0.01\\
0.56	0.01\\
0.59	0.01\\
0.62	0.01\\
0.64	0.01\\
0.67	0.01\\
0.69	0.00\\
0.72	0.01\\
0.74	0.01\\
0.77	0.01\\
0.79	0.01\\
0.82	0.01\\
0.85	0.00\\
0.87	0.00\\
0.90	0.00\\
0.92	0.00\\
0.95	0.00\\
0.97	0.00\\
1.00	0.00\\
};
\addplot [color=black, line width=1.2pt, forget plot]
  table[row sep=crcr]{%
0.05	0.00\\
0.05	1.00\\
};
\end{axis}

\begin{axis}[%
width=0.411\figurewidth,
height=0.265\figureheight,
at={(0.54\figurewidth,0\figureheight)},
scale only axis,
xmin=0.00,
xmax=1.00,
xlabel style={font=\color{mycolor1}},
xlabel={$\sym{p-Wert}$},
ymin=0.00,
ymax=1.00,
ytick={0.00,0.20,0.40,0.60,0.80,1.00},
yticklabels={\empty},
axis background/.style={fill=white},
title style={font=\bfseries\color{mycolor1}},
title={\textbf{$\sym{beta}_{22}$}},
xmajorgrids,
ymajorgrids,
grid style={dotted, opacity=0.25}
]
\addplot[const plot, color=white!10!black, line width=1.2pt, forget plot] table[row sep=crcr] {%
0.00	0.64\\
0.03	0.06\\
0.05	0.03\\
0.08	0.02\\
0.10	0.02\\
0.13	0.02\\
0.15	0.02\\
0.18	0.01\\
0.21	0.01\\
0.23	0.01\\
0.28	0.01\\
0.31	0.01\\
0.33	0.01\\
0.36	0.01\\
0.38	0.01\\
0.41	0.00\\
0.44	0.01\\
0.46	0.01\\
0.49	0.01\\
0.51	0.00\\
0.54	0.01\\
0.56	0.01\\
0.59	0.01\\
0.62	0.00\\
0.64	0.01\\
0.67	0.00\\
0.69	0.00\\
0.72	0.00\\
0.74	0.00\\
0.77	0.01\\
0.79	0.00\\
0.85	0.01\\
0.87	0.00\\
0.90	0.01\\
0.92	0.01\\
0.95	0.01\\
0.97	0.00\\
1.00	0.00\\
};
\addplot[const plot, color=white!20!black, dashed, line width=1.2pt, forget plot] table[row sep=crcr] {%
0.00	0.63\\
0.03	0.06\\
0.05	0.03\\
0.08	0.03\\
0.10	0.02\\
0.13	0.02\\
0.15	0.01\\
0.18	0.01\\
0.21	0.01\\
0.23	0.01\\
0.26	0.01\\
0.28	0.01\\
0.31	0.01\\
0.33	0.01\\
0.36	0.01\\
0.38	0.01\\
0.41	0.01\\
0.46	0.01\\
0.51	0.00\\
0.54	0.00\\
0.56	0.01\\
0.59	0.01\\
0.62	0.00\\
0.64	0.01\\
0.67	0.00\\
0.69	0.00\\
0.72	0.01\\
0.74	0.00\\
0.77	0.00\\
0.79	0.00\\
0.82	0.00\\
0.85	0.01\\
0.87	0.01\\
0.90	0.00\\
0.92	0.00\\
0.95	0.01\\
0.97	0.00\\
1.00	0.00\\
};
\addplot[const plot, color=white!30!black, dotted, line width=1.2pt, forget plot] table[row sep=crcr] {%
0.00	0.63\\
0.03	0.06\\
0.05	0.03\\
0.08	0.02\\
0.10	0.02\\
0.13	0.02\\
0.15	0.01\\
0.18	0.01\\
0.21	0.01\\
0.23	0.01\\
0.26	0.01\\
0.28	0.01\\
0.31	0.01\\
0.33	0.01\\
0.36	0.01\\
0.38	0.01\\
0.41	0.01\\
0.44	0.01\\
0.46	0.01\\
0.49	0.00\\
0.51	0.01\\
0.54	0.00\\
0.56	0.00\\
0.59	0.00\\
0.62	0.01\\
0.64	0.01\\
0.67	0.00\\
0.69	0.01\\
0.72	0.00\\
0.77	0.00\\
0.79	0.00\\
0.82	0.01\\
0.85	0.00\\
0.87	0.01\\
0.90	0.01\\
0.92	0.00\\
0.95	0.01\\
0.97	0.00\\
1.00	0.00\\
};
\addplot[const plot, color=white!40!black, dashdotted, line width=1.2pt, forget plot] table[row sep=crcr] {%
0.00	0.63\\
0.03	0.06\\
0.05	0.03\\
0.08	0.02\\
0.10	0.02\\
0.13	0.02\\
0.15	0.01\\
0.18	0.01\\
0.21	0.01\\
0.23	0.01\\
0.26	0.01\\
0.28	0.01\\
0.31	0.01\\
0.33	0.01\\
0.36	0.01\\
0.38	0.01\\
0.41	0.01\\
0.44	0.01\\
0.46	0.01\\
0.49	0.01\\
0.51	0.00\\
0.54	0.01\\
0.56	0.01\\
0.59	0.00\\
0.62	0.00\\
0.64	0.01\\
0.67	0.01\\
0.69	0.01\\
0.72	0.00\\
0.74	0.01\\
0.77	0.00\\
0.79	0.01\\
0.82	0.00\\
0.85	0.01\\
0.87	0.01\\
0.90	0.01\\
0.92	0.00\\
0.95	0.00\\
0.97	0.00\\
1.00	0.00\\
};
\addplot[const plot, color=gray, line width=1.2pt, forget plot] table[row sep=crcr] {%
0.00	0.63\\
0.03	0.06\\
0.05	0.04\\
0.08	0.03\\
0.10	0.02\\
0.13	0.02\\
0.15	0.01\\
0.18	0.01\\
0.21	0.01\\
0.23	0.01\\
0.26	0.01\\
0.28	0.01\\
0.31	0.01\\
0.33	0.01\\
0.36	0.01\\
0.38	0.01\\
0.41	0.01\\
0.49	0.01\\
0.51	0.00\\
0.56	0.01\\
0.59	0.01\\
0.62	0.01\\
0.64	0.00\\
0.67	0.00\\
0.69	0.01\\
0.72	0.00\\
0.74	0.01\\
0.77	0.00\\
0.79	0.00\\
0.82	0.01\\
0.85	0.00\\
0.87	0.01\\
0.90	0.00\\
0.92	0.00\\
0.97	0.00\\
1.00	0.00\\
};
\addplot[const plot, color=white!60!black, dashed, line width=1.2pt, forget plot] table[row sep=crcr] {%
0.00	0.62\\
0.03	0.06\\
0.05	0.04\\
0.08	0.02\\
0.10	0.02\\
0.13	0.02\\
0.15	0.02\\
0.18	0.01\\
0.21	0.01\\
0.23	0.01\\
0.26	0.01\\
0.28	0.01\\
0.31	0.01\\
0.33	0.01\\
0.36	0.01\\
0.38	0.01\\
0.41	0.01\\
0.44	0.01\\
0.46	0.01\\
0.49	0.01\\
0.51	0.01\\
0.54	0.01\\
0.56	0.01\\
0.59	0.00\\
0.62	0.01\\
0.67	0.00\\
0.69	0.01\\
0.72	0.00\\
0.77	0.00\\
0.79	0.00\\
0.82	0.00\\
0.85	0.00\\
0.87	0.00\\
0.92	0.00\\
0.97	0.01\\
1.00	0.01\\
};
\addplot[const plot, color=white!70!black, dotted, line width=1.2pt, forget plot] table[row sep=crcr] {%
0.00	0.68\\
0.03	0.05\\
0.05	0.03\\
0.08	0.02\\
0.10	0.02\\
0.13	0.01\\
0.15	0.01\\
0.18	0.01\\
0.21	0.01\\
0.23	0.01\\
0.26	0.01\\
0.28	0.01\\
0.31	0.00\\
0.33	0.01\\
0.36	0.01\\
0.38	0.01\\
0.41	0.01\\
0.44	0.01\\
0.46	0.01\\
0.51	0.01\\
0.54	0.01\\
0.56	0.00\\
0.59	0.00\\
0.62	0.00\\
0.67	0.00\\
0.69	0.01\\
0.72	0.00\\
0.74	0.00\\
0.77	0.00\\
0.79	0.00\\
0.82	0.00\\
0.85	0.00\\
0.87	0.00\\
0.90	0.00\\
0.92	0.00\\
0.95	0.00\\
1.00	0.00\\
};
\addplot[const plot, color=white!80!black, dashdotted, line width=1.2pt, forget plot] table[row sep=crcr] {%
0.00	0.61\\
0.03	0.07\\
0.05	0.04\\
0.08	0.02\\
0.10	0.02\\
0.13	0.02\\
0.15	0.02\\
0.18	0.01\\
0.23	0.01\\
0.26	0.01\\
0.28	0.01\\
0.31	0.01\\
0.33	0.01\\
0.36	0.01\\
0.38	0.01\\
0.41	0.01\\
0.44	0.01\\
0.46	0.01\\
0.49	0.01\\
0.51	0.01\\
0.54	0.01\\
0.56	0.01\\
0.59	0.01\\
0.62	0.01\\
0.67	0.01\\
0.69	0.00\\
0.72	0.01\\
0.74	0.00\\
0.77	0.00\\
0.79	0.00\\
0.82	0.00\\
0.85	0.00\\
0.87	0.01\\
0.90	0.01\\
0.92	0.00\\
0.95	0.00\\
0.97	0.00\\
1.00	0.00\\
};
\addplot[const plot, color=white!90!black, line width=1.2pt, forget plot] table[row sep=crcr] {%
0.00	0.63\\
0.03	0.06\\
0.05	0.03\\
0.08	0.03\\
0.10	0.02\\
0.13	0.02\\
0.15	0.01\\
0.18	0.01\\
0.21	0.01\\
0.23	0.01\\
0.26	0.01\\
0.28	0.01\\
0.31	0.01\\
0.33	0.01\\
0.36	0.01\\
0.38	0.01\\
0.41	0.01\\
0.44	0.01\\
0.46	0.01\\
0.49	0.00\\
0.51	0.01\\
0.54	0.00\\
0.56	0.01\\
0.59	0.01\\
0.62	0.01\\
0.64	0.00\\
0.67	0.00\\
0.69	0.01\\
0.72	0.00\\
0.74	0.01\\
0.77	0.01\\
0.79	0.00\\
0.82	0.01\\
0.85	0.00\\
0.87	0.00\\
0.90	0.00\\
0.92	0.00\\
0.95	0.00\\
0.97	0.00\\
1.00	0.00\\
};
\addplot [color=black, line width=1.2pt, forget plot]
  table[row sep=crcr]{%
0.05	0.00\\
0.05	1.00\\
};
\addplot [color=black, line width=1.2pt, forget plot]
  table[row sep=crcr]{%
0.05	0.00\\
0.05	1.00\\
};
\addplot [color=black, line width=1.2pt, forget plot]
  table[row sep=crcr]{%
0.05	0.00\\
0.05	1.00\\
};
\addplot [color=black, line width=1.2pt, forget plot]
  table[row sep=crcr]{%
0.05	0.00\\
0.05	1.00\\
};
\addplot [color=black, line width=1.2pt, forget plot]
  table[row sep=crcr]{%
0.05	0.00\\
0.05	1.00\\
};
\addplot [color=black, line width=1.2pt, forget plot]
  table[row sep=crcr]{%
0.05	0.00\\
0.05	1.00\\
};
\end{axis}
\end{tikzpicture}%
    \caption{Relative Häufigkeiten der in der Simulation ermittelten $\sym{p-Wert}$e der Signifikanztests.}
    \label{fig:app.c_pval_hist}
\end{figure}

\begin{figure}[htbp]
    \centering
    \setlength\figurewidth{0.9\textwidth}
    \setlength\figureheight{19cm}
    % This file was created by matlab2tikz.
%
\definecolor{mycolor1}{rgb}{0.12941,0.12941,0.12941}%
%
\begin{tikzpicture}

\begin{axis}[%
width=0.254\figurewidth,
height=0.126\figureheight,
at={(0\figurewidth,0.874\figureheight)},
scale only axis,
clip=false,
xmin=0.00,
xmax=0.40,
xtick={0.00,0.20,0.40},
xticklabels={\empty},
ymin=-0.02,
ymax=0.06,
ylabel style={font=\color{mycolor1}},
ylabel={$\sym{beta}_0$},
axis background/.style={fill=white},
xmajorgrids,
ymajorgrids,
grid style={dotted, opacity=0.25}
]
\addplot [color=black, only marks, mark size=2.0pt, mark=*, mark options={solid, fill=black, black}, forget plot]
  table[row sep=crcr]{%
0.00	0.00\\
0.04	0.01\\
0.04	0.00\\
0.07	0.01\\
0.00	0.00\\
0.35	0.00\\
0.01	-0.02\\
0.18	0.03\\
0.00	0.06\\
};
\addplot [color=black, dashed, forget plot]
  table[row sep=crcr]{%
0.00	0.01\\
0.35	0.01\\
};
\node[left, align=right, inner sep=0, font=\color{mycolor1}]
at (rel axis cs:0.95,0.9) {$\sym{R2}=0,00$};
\end{axis}

\begin{axis}[%
width=0.254\figurewidth,
height=0.126\figureheight,
at={(0.345\figurewidth,0.874\figureheight)},
scale only axis,
clip=false,
xmin=68.98,
xmax=85.65,
xtick={70.00,80.00},
xticklabels={\empty},
ymin=-0.02,
ymax=0.06,
ytick={0.00,0.02,0.04},
yticklabels={\empty},
axis background/.style={fill=white},
xmajorgrids,
ymajorgrids,
grid style={dotted, opacity=0.25}
]
\addplot [color=black, only marks, mark size=2.0pt, mark=*, mark options={solid, fill=black, black}, forget plot]
  table[row sep=crcr]{%
81.19	0.00\\
76.00	0.01\\
85.05	0.00\\
73.48	0.01\\
80.32	0.00\\
82.97	-0.02\\
68.98	0.03\\
85.65	0.06\\
};
\addplot [color=black, dashed, forget plot]
  table[row sep=crcr]{%
68.98	0.01\\
85.65	0.01\\
};
\node[left, align=right, inner sep=0, font=\color{mycolor1}]
at (rel axis cs:0.95,0.9) {$\sym{R2}=0,01$};
\end{axis}

\begin{axis}[%
width=0.254\figurewidth,
height=0.126\figureheight,
at={(0.689\figurewidth,0.874\figureheight)},
scale only axis,
clip=false,
xmin=6.84,
xmax=8.00,
xtick={7.00,7.50,8.00},
xticklabels={\empty},
ymin=-0.02,
ymax=0.06,
ytick={0.00,0.02,0.04,0.06},
yticklabels={\empty},
axis background/.style={fill=white},
xmajorgrids,
ymajorgrids,
grid style={dotted, opacity=0.25}
]
\addplot [color=black, only marks, mark size=2.0pt, mark=*, mark options={solid, fill=black, black}, forget plot]
  table[row sep=crcr]{%
6.95	0.00\\
6.87	0.01\\
6.90	0.00\\
6.88	0.01\\
6.88	0.00\\
6.97	0.00\\
6.84	-0.02\\
7.16	0.03\\
7.95	0.06\\
};
\addplot [color=black, dashed, forget plot]
  table[row sep=crcr]{%
6.84	-0.00\\
7.95	0.06\\
};
\node[left, align=right, inner sep=0, font=\color{mycolor1}]
at (rel axis cs:0.95,0.9) {$\sym{R2}=0,82$};
\end{axis}

\begin{axis}[%
width=0.254\figurewidth,
height=0.126\figureheight,
at={(0\figurewidth,0.699\figureheight)},
scale only axis,
clip=false,
xmin=0.00,
xmax=0.40,
xtick={0.00,0.20,0.40},
xticklabels={\empty},
ymin=-0.89,
ymax=10.00,
ylabel style={font=\color{mycolor1}},
ylabel={$\sym{beta}_1$},
axis background/.style={fill=white},
xmajorgrids,
ymajorgrids,
grid style={dotted, opacity=0.25}
]
\addplot [color=black, only marks, mark size=2.0pt, mark=*, mark options={solid, fill=black, black}, forget plot]
  table[row sep=crcr]{%
0.00	0.00\\
0.04	2.04\\
0.04	-0.89\\
0.07	2.81\\
0.00	0.04\\
0.35	2.38\\
0.01	0.31\\
0.18	8.40\\
0.00	0.25\\
};
\addplot [color=black, dashed, forget plot]
  table[row sep=crcr]{%
0.00	0.71\\
0.35	5.24\\
};
\node[left, align=right, inner sep=0, font=\color{mycolor1}]
at (rel axis cs:0.95,0.9) {$\sym{R2}=0,29$};
\end{axis}

\begin{axis}[%
width=0.254\figurewidth,
height=0.126\figureheight,
at={(0.345\figurewidth,0.699\figureheight)},
scale only axis,
clip=false,
xmin=68.98,
xmax=85.65,
xtick={70.00,80.00},
xticklabels={\empty},
ymin=-1.12,
ymax=10.00,
ytick={0.00,5.00,10.00},
yticklabels={\empty},
axis background/.style={fill=white},
xmajorgrids,
ymajorgrids,
grid style={dotted, opacity=0.25}
]
\addplot [color=black, only marks, mark size=2.0pt, mark=*, mark options={solid, fill=black, black}, forget plot]
  table[row sep=crcr]{%
81.19	0.00\\
76.00	2.04\\
85.05	-0.89\\
73.48	2.81\\
81.18	0.04\\
80.32	2.38\\
82.97	0.31\\
68.98	8.40\\
85.65	0.25\\
};
\addplot [color=black, dashed, forget plot]
  table[row sep=crcr]{%
68.98	6.44\\
85.65	-1.12\\
};
\node[left, align=right, inner sep=0, font=\color{mycolor1}]
at (rel axis cs:0.95,0.9) {$\sym{R2}=0,80$};
\end{axis}

\begin{axis}[%
width=0.254\figurewidth,
height=0.126\figureheight,
at={(0.689\figurewidth,0.699\figureheight)},
scale only axis,
clip=false,
xmin=6.61,
xmax=11.91,
xtick={7.00,8.00,9.00,10.00,11.00},
xticklabels={\empty},
ymin=-0.89,
ymax=10.00,
ytick={0.00,5.00,10.00},
yticklabels={\empty},
axis background/.style={fill=white},
xmajorgrids,
ymajorgrids,
grid style={dotted, opacity=0.25}
]
\addplot [color=black, only marks, mark size=2.0pt, mark=*, mark options={solid, fill=black, black}, forget plot]
  table[row sep=crcr]{%
6.97	0.00\\
8.03	2.04\\
6.61	-0.89\\
8.61	2.81\\
7.13	0.04\\
8.08	2.38\\
7.15	0.31\\
11.91	8.40\\
6.94	0.25\\
};
\addplot [color=black, dashed, forget plot]
  table[row sep=crcr]{%
6.61	-0.57\\
11.91	8.53\\
};
\node[left, align=right, inner sep=0, font=\color{mycolor1}]
at (rel axis cs:0.95,0.9) {$\sym{R2}=0,99$};
\end{axis}

\begin{axis}[%
width=0.254\figurewidth,
height=0.126\figureheight,
at={(0\figurewidth,0.524\figureheight)},
scale only axis,
clip=false,
xmin=0.00,
xmax=0.40,
xtick={0.00,0.20,0.40},
xticklabels={\empty},
ymin=-1.01,
ymax=0.43,
ylabel style={font=\color{mycolor1}},
ylabel={$\sym{beta}_2$},
axis background/.style={fill=white},
xmajorgrids,
ymajorgrids,
grid style={dotted, opacity=0.25}
]
\addplot [color=black, only marks, mark size=2.0pt, mark=*, mark options={solid, fill=black, black}, forget plot]
  table[row sep=crcr]{%
0.00	0.00\\
0.04	0.06\\
0.04	-0.01\\
0.07	0.05\\
0.00	-0.01\\
0.35	-0.72\\
0.01	-1.01\\
0.18	0.43\\
0.00	0.35\\
};
\addplot [color=black, dashed, forget plot]
  table[row sep=crcr]{%
0.00	-0.02\\
0.35	-0.36\\
};
\node[left, align=right, inner sep=0, font=\color{mycolor1}]
at (rel axis cs:0.95,0.9) {$\sym{R2}=0,06$};
\end{axis}

\begin{axis}[%
width=0.254\figurewidth,
height=0.126\figureheight,
at={(0.345\figurewidth,0.524\figureheight)},
scale only axis,
clip=false,
xmin=68.98,
xmax=85.65,
xtick={70.00,80.00},
xticklabels={\empty},
ymin=-1.01,
ymax=0.43,
ytick={-1.00,-0.50,0.00},
yticklabels={\empty},
axis background/.style={fill=white},
xmajorgrids,
ymajorgrids,
grid style={dotted, opacity=0.25}
]
\addplot [color=black, only marks, mark size=2.0pt, mark=*, mark options={solid, fill=black, black}, forget plot]
  table[row sep=crcr]{%
81.19	0.00\\
76.00	0.06\\
85.05	-0.01\\
73.48	0.05\\
81.18	-0.01\\
80.32	-0.72\\
82.97	-1.01\\
68.98	0.43\\
85.65	0.35\\
};
\addplot [color=black, dashed, forget plot]
  table[row sep=crcr]{%
68.98	0.21\\
85.65	-0.28\\
};
\node[left, align=right, inner sep=0, font=\color{mycolor1}]
at (rel axis cs:0.95,0.9) {$\sym{R2}=0,12$};
\end{axis}

\begin{axis}[%
width=0.254\figurewidth,
height=0.126\figureheight,
at={(0.689\figurewidth,0.524\figureheight)},
scale only axis,
clip=false,
xmin=6.71,
xmax=7.12,
xtick={6.80,7.00},
xticklabels={\empty},
ymin=-1.01,
ymax=0.43,
ytick={-1.00,-0.50,0.00},
yticklabels={\empty},
axis background/.style={fill=white},
xmajorgrids,
ymajorgrids,
grid style={dotted, opacity=0.25}
]
\addplot [color=black, only marks, mark size=2.0pt, mark=*, mark options={solid, fill=black, black}, forget plot]
  table[row sep=crcr]{%
7.12	0.00\\
7.08	0.06\\
7.07	-0.01\\
6.97	0.05\\
6.91	-0.01\\
6.78	-0.72\\
6.71	-1.01\\
6.91	0.43\\
7.05	0.35\\
};
\addplot [color=black, dashed, forget plot]
  table[row sep=crcr]{%
6.71	-0.68\\
7.12	0.32\\
};
\node[left, align=right, inner sep=0, font=\color{mycolor1}]
at (rel axis cs:0.95,0.9) {$\sym{R2}=0,53$};
\end{axis}

\begin{axis}[%
width=0.254\figurewidth,
height=0.126\figureheight,
at={(0\figurewidth,0.35\figureheight)},
scale only axis,
clip=false,
xmin=0.00,
xmax=0.40,
xtick={0.00,0.20,0.40},
xticklabels={\empty},
ymin=-4.00,
ymax=0.62,
ylabel style={font=\color{mycolor1}},
ylabel={$\sym{beta}_{12}$},
axis background/.style={fill=white},
xmajorgrids,
ymajorgrids,
grid style={dotted, opacity=0.25}
]
\addplot [color=black, only marks, mark size=2.0pt, mark=*, mark options={solid, fill=black, black}, forget plot]
  table[row sep=crcr]{%
0.00	0.00\\
0.04	-0.09\\
0.04	0.06\\
0.07	-0.07\\
0.00	0.02\\
0.35	-3.94\\
0.01	-0.16\\
0.18	0.62\\
0.00	0.57\\
};
\addplot [color=black, dashed, forget plot]
  table[row sep=crcr]{%
0.00	0.39\\
0.35	-2.91\\
};
\node[left, align=right, inner sep=0, font=\color{mycolor1}]
at (rel axis cs:0.95,0.9) {$\sym{R2}=0,64$};
\end{axis}

\begin{axis}[%
width=0.254\figurewidth,
height=0.126\figureheight,
at={(0.345\figurewidth,0.35\figureheight)},
scale only axis,
clip=false,
xmin=68.98,
xmax=85.65,
xtick={70.00,80.00},
xticklabels={\empty},
ymin=-4.00,
ymax=0.62,
ytick={-4.00,-2.00,0.00},
yticklabels={\empty},
axis background/.style={fill=white},
xmajorgrids,
ymajorgrids,
grid style={dotted, opacity=0.25}
]
\addplot [color=black, only marks, mark size=2.0pt, mark=*, mark options={solid, fill=black, black}, forget plot]
  table[row sep=crcr]{%
81.19	0.00\\
76.00	-0.09\\
85.05	0.06\\
73.48	-0.07\\
81.18	0.02\\
80.32	-3.94\\
82.97	-0.16\\
68.98	0.62\\
85.65	0.57\\
};
\addplot [color=black, dashed, forget plot]
  table[row sep=crcr]{%
68.98	-0.08\\
85.65	-0.48\\
};
\node[left, align=right, inner sep=0, font=\color{mycolor1}]
at (rel axis cs:0.95,0.9) {$\sym{R2}=0,01$};
\end{axis}

\begin{axis}[%
width=0.254\figurewidth,
height=0.126\figureheight,
at={(0.689\figurewidth,0.35\figureheight)},
scale only axis,
clip=false,
xmin=11.73,
xmax=14.33,
xtick={12.00,13.00,14.00},
xticklabels={\empty},
ymin=-4.00,
ymax=0.62,
ytick={-4.00,-2.00,0.00},
yticklabels={\empty},
axis background/.style={fill=white},
xmajorgrids,
ymajorgrids,
grid style={dotted, opacity=0.25}
]
\addplot [color=black, only marks, mark size=2.0pt, mark=*, mark options={solid, fill=black, black}, forget plot]
  table[row sep=crcr]{%
13.89	0.00\\
13.96	-0.09\\
14.33	0.06\\
14.18	-0.07\\
14.00	0.02\\
11.73	-3.94\\
13.95	-0.16\\
13.87	0.62\\
14.15	0.57\\
};
\addplot [color=black, dashed, forget plot]
  table[row sep=crcr]{%
11.73	-3.80\\
14.33	0.58\\
};
\node[left, align=right, inner sep=0, font=\color{mycolor1}]
at (rel axis cs:0.95,0.9) {$\sym{R2}=0,92$};
\end{axis}

\begin{axis}[%
width=0.254\figurewidth,
height=0.126\figureheight,
at={(0\figurewidth,0.175\figureheight)},
scale only axis,
clip=false,
xmin=0.00,
xmax=0.40,
xtick={0.00,0.20,0.40},
xticklabels={\empty},
ymin=-4.57,
ymax=10.52,
ylabel style={font=\color{mycolor1}},
ylabel={$\sym{beta}_{11}$},
axis background/.style={fill=white},
xmajorgrids,
ymajorgrids,
grid style={dotted, opacity=0.25}
]
\addplot [color=black, only marks, mark size=2.0pt, mark=*, mark options={solid, fill=black, black}, forget plot]
  table[row sep=crcr]{%
0.00	0.00\\
0.04	4.34\\
0.04	-4.57\\
0.07	6.25\\
0.00	-0.05\\
0.35	2.74\\
0.01	1.19\\
0.18	10.52\\
0.00	0.84\\
};
\addplot [color=black, dashed, forget plot]
  table[row sep=crcr]{%
0.00	1.20\\
0.35	6.52\\
};
\node[left, align=right, inner sep=0, font=\color{mycolor1}]
at (rel axis cs:0.95,0.9) {$\sym{R2}=0,17$};
\end{axis}

\begin{axis}[%
width=0.254\figurewidth,
height=0.126\figureheight,
at={(0.345\figurewidth,0.175\figureheight)},
scale only axis,
clip=false,
xmin=68.98,
xmax=85.65,
xtick={70.00,80.00},
xticklabels={\empty},
ymin=-4.57,
ymax=10.52,
ytick={0.00,5.00,10.00},
yticklabels={\empty},
axis background/.style={fill=white},
xmajorgrids,
ymajorgrids,
grid style={dotted, opacity=0.25}
]
\addplot [color=black, only marks, mark size=2.0pt, mark=*, mark options={solid, fill=black, black}, forget plot]
  table[row sep=crcr]{%
81.19	0.00\\
76.00	4.34\\
85.05	-4.57\\
73.48	6.25\\
81.18	-0.05\\
80.32	2.74\\
82.97	1.19\\
68.98	10.52\\
85.65	0.84\\
};
\addplot [color=black, dashed, forget plot]
  table[row sep=crcr]{%
68.98	9.82\\
85.65	-2.09\\
};
\node[left, align=right, inner sep=0, font=\color{mycolor1}]
at (rel axis cs:0.95,0.9) {$\sym{R2}=0,84$};
\end{axis}

\begin{axis}[%
width=0.254\figurewidth,
height=0.126\figureheight,
at={(0.689\figurewidth,0.175\figureheight)},
scale only axis,
clip=false,
xmin=4.98,
xmax=12.29,
xtick={5.00,10.00},
xticklabels={\empty},
ymin=-4.57,
ymax=10.52,
ytick={0.00,5.00,10.00},
yticklabels={\empty},
axis background/.style={fill=white},
xmajorgrids,
ymajorgrids,
grid style={dotted, opacity=0.25}
]
\addplot [color=black, only marks, mark size=2.0pt, mark=*, mark options={solid, fill=black, black}, forget plot]
  table[row sep=crcr]{%
7.04	0.00\\
9.33	4.34\\
4.98	-4.57\\
10.37	6.25\\
7.05	-0.05\\
8.62	2.74\\
7.36	1.19\\
12.29	10.52\\
7.28	0.84\\
};
\addplot [color=black, dashed, forget plot]
  table[row sep=crcr]{%
4.98	-4.14\\
12.29	10.38\\
};
\node[left, align=right, inner sep=0, font=\color{mycolor1}]
at (rel axis cs:0.95,0.9) {$\sym{R2}=0,99$};
\end{axis}

\begin{axis}[%
width=0.254\figurewidth,
height=0.126\figureheight,
at={(0\figurewidth,0\figureheight)},
scale only axis,
clip=false,
xmin=0.00,
xmax=0.40,
xlabel style={font=\color{mycolor1}},
xlabel={$\sym{A2}$},
ymin=-4.00,
ymax=2.00,
ylabel style={font=\color{mycolor1}},
ylabel={$\sym{beta}_{22}$},
axis background/.style={fill=white},
xmajorgrids,
ymajorgrids,
grid style={dotted, opacity=0.25}
]
\addplot [color=black, only marks, mark size=2.0pt, mark=*, mark options={solid, fill=black, black}, forget plot]
  table[row sep=crcr]{%
0.00	0.00\\
0.04	0.08\\
0.04	0.05\\
0.07	-0.08\\
0.00	-0.07\\
0.35	1.35\\
0.01	-2.98\\
0.18	1.75\\
0.00	0.93\\
};
\addplot [color=black, dashed, forget plot]
  table[row sep=crcr]{%
0.00	-0.36\\
0.35	1.82\\
};
\node[left, align=right, inner sep=0, font=\color{mycolor1}]
at (rel axis cs:0.95,0.9) {$\sym{R2}=0,30$};
\end{axis}

\begin{axis}[%
width=0.254\figurewidth,
height=0.126\figureheight,
at={(0.345\figurewidth,0\figureheight)},
scale only axis,
clip=false,
xmin=68.98,
xmax=85.65,
xlabel style={font=\color{mycolor1}},
xlabel={$\sym{eff_I}$},
ymin=-4.00,
ymax=2.00,
ytick={-4.00,-2.00,0.00,2.00},
yticklabels={\empty},
axis background/.style={fill=white},
xmajorgrids,
ymajorgrids,
grid style={dotted, opacity=0.25}
]
\addplot [color=black, only marks, mark size=2.0pt, mark=*, mark options={solid, fill=black, black}, forget plot]
  table[row sep=crcr]{%
81.19	0.00\\
76.00	0.08\\
85.05	0.05\\
73.48	-0.08\\
81.18	-0.07\\
80.32	1.35\\
82.97	-2.98\\
68.98	1.75\\
85.65	0.93\\
};
\addplot [color=black, dashed, forget plot]
  table[row sep=crcr]{%
68.98	1.03\\
85.65	-0.43\\
};
\node[left, align=right, inner sep=0, font=\color{mycolor1}]
at (rel axis cs:0.95,0.9) {$\sym{R2}=0,13$};
\end{axis}

\begin{axis}[%
width=0.254\figurewidth,
height=0.126\figureheight,
at={(0.689\figurewidth,0\figureheight)},
scale only axis,
clip=false,
xmin=5.89,
xmax=8.00,
xlabel style={font=\color{mycolor1}},
xlabel={$\sym{mse}$},
ymin=-4.00,
ymax=2.00,
ytick={-4.00,-2.00,0.00,2.00},
yticklabels={\empty},
axis background/.style={fill=white},
xmajorgrids,
ymajorgrids,
grid style={dotted, opacity=0.25}
]
\addplot [color=black, only marks, mark size=2.0pt, mark=*, mark options={solid, fill=black, black}, forget plot]
  table[row sep=crcr]{%
6.99	0.00\\
7.13	0.08\\
6.88	0.05\\
7.22	-0.08\\
7.02	-0.07\\
7.73	1.35\\
5.89	-2.98\\
7.81	1.75\\
7.31	0.93\\
};
\addplot [color=black, dashed, forget plot]
  table[row sep=crcr]{%
5.89	-2.76\\
7.81	1.77\\
};
\node[left, align=right, inner sep=0, font=\color{mycolor1}]
at (rel axis cs:0.95,0.9) {$\sym{R2}=0,95$};
\end{axis}
\end{tikzpicture}%
    \caption{Korrelierende Metriken zur prozentualen Trennschärfe-Änderung ($\Delta \sym{power}$ [\%]) über alle Koeffizienten und Szenarien ergänzend zu Abbildung~\ref{fig:ma_abb5.1.4_korrelation}.}
    \label{fig:app.c_corr_all}
\end{figure}

\clearpage
\subsection{Szenarien bei Arbeitspunkt SNR$=1$} \label{app:c_snr1}

\begin{figure}[htbp]
    \centering
    \setlength\figurewidth{0.9\textwidth}
    \setlength\figureheight{4cm}
    % This file was created by matlab2tikz.
%
\definecolor{mycolor1}{rgb}{0.12941,0.12941,0.12941}%
%
\begin{tikzpicture}

\begin{axis}[%
width=0.751\figurewidth,
height=\figureheight,
at={(0\figurewidth,0\figureheight)},
scale only axis,
bar shift auto,
xmin=0.54,
xmax=6.46,
xtick={1.00,2.00,3.00,4.00,5.00,6.00},
xticklabels={{$\sym{beta}_0$},{$\sym{beta}_1$},{$\sym{beta}_2$},{$\sym{beta}_{12}$},{$\sym{beta}_{11}$},{$\sym{beta}_{22}$}},
ymin=0.00,
ymax=10.00,
ylabel style={font=\color{mycolor1}},
ylabel={\sym{mse}},
axis background/.style={fill=white},
xmajorgrids,
ymajorgrids,
grid style={dotted, opacity=0.25},
legend style={at={(1.03,1)}, anchor=north west, legend cell align=left, align=left}
]
\addplot[ybar, bar width=0.071, fill=white!10!black, draw=mycolor1, area legend] table[row sep=crcr] {%
1.00	4.45\\
2.00	4.46\\
3.00	4.56\\
4.00	8.89\\
5.00	4.50\\
6.00	4.47\\
};
\addplot[forget plot, color=mycolor1] table[row sep=crcr] {%
0.54	0.00\\
6.46	0.00\\
};
\addlegendentry{Basis \ac{CCD}}

\addplot[ybar, bar width=0.071, fill=black!25!darkgray, draw=mycolor1, area legend] table[row sep=crcr] {%
1.00	4.39\\
2.00	5.14\\
3.00	4.53\\
4.00	8.93\\
5.00	5.97\\
6.00	4.56\\
};
\addplot[forget plot, color=mycolor1] table[row sep=crcr] {%
0.54	0.00\\
6.46	0.00\\
};
\addlegendentry{Fall I ($-20\,\%$)}

\addplot[ybar, bar width=0.071, fill=black!45!gray, draw=mycolor1, area legend] table[row sep=crcr] {%
1.00	4.42\\
2.00	4.23\\
3.00	4.52\\
4.00	9.17\\
5.00	3.19\\
6.00	4.41\\
};
\addplot[forget plot, color=mycolor1] table[row sep=crcr] {%
0.54	0.00\\
6.46	0.00\\
};
\addlegendentry{Fall I ($+20\,\%$)}

\addplot[ybar, bar width=0.071, fill=white!15!darkgray, draw=mycolor1, area legend] table[row sep=crcr] {%
1.00	4.40\\
2.00	5.51\\
3.00	4.46\\
4.00	9.08\\
5.00	6.64\\
6.00	4.62\\
};
\addplot[forget plot, color=mycolor1] table[row sep=crcr] {%
0.54	0.00\\
6.46	0.00\\
};
\addlegendentry{Fall I (\ac{FC})}

\addplot[ybar, bar width=0.071, fill=white!45!black, draw=mycolor1, area legend] table[row sep=crcr] {%
1.00	4.41\\
2.00	4.56\\
3.00	4.42\\
4.00	8.96\\
5.00	4.51\\
6.00	4.49\\
};
\addplot[forget plot, color=mycolor1] table[row sep=crcr] {%
0.54	0.00\\
6.46	0.00\\
};
\addlegendentry{Fall II (leicht)}

\addplot[ybar, bar width=0.071, fill=gray!85!lightgray, draw=mycolor1, area legend] table[row sep=crcr] {%
1.00	4.46\\
2.00	5.17\\
3.00	4.34\\
4.00	7.51\\
5.00	5.52\\
6.00	4.95\\
};
\addplot[forget plot, color=mycolor1] table[row sep=crcr] {%
0.54	0.00\\
6.46	0.00\\
};
\addlegendentry{Fall II (stark)}

\addplot[ybar, bar width=0.071, fill=white!50!darkgray, draw=mycolor1, area legend] table[row sep=crcr] {%
1.00	4.37\\
2.00	4.58\\
3.00	4.30\\
4.00	8.93\\
5.00	4.71\\
6.00	3.77\\
};
\addplot[forget plot, color=mycolor1] table[row sep=crcr] {%
0.54	0.00\\
6.46	0.00\\
};
\addlegendentry{Fall III}

\addplot[ybar, bar width=0.071, fill=black!5!lightgray, draw=mycolor1, area legend] table[row sep=crcr] {%
1.00	4.58\\
2.00	7.62\\
3.00	4.42\\
4.00	8.88\\
5.00	7.87\\
6.00	5.00\\
};
\addplot[forget plot, color=mycolor1] table[row sep=crcr] {%
0.54	0.00\\
6.46	0.00\\
};
\addlegendentry{Fall IV: \ac{SP}}

\addplot[ybar, bar width=0.071, fill=white!80!black, draw=mycolor1, area legend] table[row sep=crcr] {%
1.00	5.09\\
2.00	4.44\\
3.00	4.51\\
4.00	9.05\\
5.00	4.66\\
6.00	4.68\\
};
\addplot[forget plot, color=mycolor1] table[row sep=crcr] {%
0.54	0.00\\
6.46	0.00\\
};
\addlegendentry{Fall IV: \ac{ZP}}

\end{axis}
\end{tikzpicture}%
    \caption{Güte der Schätzer (\ac{MSE}) für $\text{SNR}=1$.}
    \label{fig:app.c_msedev_snr1}
\end{figure}

\begin{figure}[htbp]
    \centering
    \setlength\figurewidth{0.9\textwidth}
    \setlength\figureheight{4cm}
    % This file was created by matlab2tikz.
%
\definecolor{mycolor1}{rgb}{0.12941,0.12941,0.12941}%
%
\begin{tikzpicture}

\begin{axis}[%
width=0.751\figurewidth,
height=\figureheight,
at={(0\figurewidth,0\figureheight)},
scale only axis,
bar shift auto,
xmin=0.54,
xmax=6.46,
xtick={1.00,2.00,3.00,4.00,5.00,6.00},
xticklabels={{$\sym{beta}_0$},{$\sym{beta}_1$},{$\sym{beta}_2$},{$\sym{beta}_{12}$},{$\sym{beta}_{11}$},{$\sym{beta}_{22}$}},
ymin=0.00,
ymax=100.00,
ylabel style={font=\color{mycolor1}},
ylabel={$\sym{power} [\%]$},
axis background/.style={fill=white},
xmajorgrids,
ymajorgrids,
grid style={dotted, opacity=0.25},
legend style={at={(1.03,1)}, anchor=north west, legend cell align=left, align=left}
]
\addplot[ybar, bar width=0.071, fill=white!10!black, draw=mycolor1, area legend] table[row sep=crcr] {%
1.00	99.83\\
2.00	88.04\\
3.00	87.77\\
4.00	50.24\\
5.00	63.33\\
6.00	63.82\\
};
\addplot[forget plot, color=mycolor1] table[row sep=crcr] {%
0.54	0.00\\
6.46	0.00\\
};
\addlegendentry{Basis \ac{CCD}}

\addplot[ybar, bar width=0.071, fill=black!25!darkgray, draw=mycolor1, area legend] table[row sep=crcr] {%
1.00	99.88\\
2.00	86.75\\
3.00	87.46\\
4.00	50.56\\
5.00	58.63\\
6.00	63.87\\
};
\addplot[forget plot, color=mycolor1] table[row sep=crcr] {%
0.54	0.00\\
6.46	0.00\\
};
\addlegendentry{Fall I ($-20\,\%$)}

\addplot[ybar, bar width=0.071, fill=black!45!gray, draw=mycolor1, area legend] table[row sep=crcr] {%
1.00	99.88\\
2.00	88.25\\
3.00	88.44\\
4.00	50.81\\
5.00	69.08\\
6.00	63.13\\
};
\addplot[forget plot, color=mycolor1] table[row sep=crcr] {%
0.54	0.00\\
6.46	0.00\\
};
\addlegendentry{Fall I ($+20\,\%$)}

\addplot[ybar, bar width=0.071, fill=white!15!darkgray, draw=mycolor1, area legend] table[row sep=crcr] {%
1.00	99.84\\
2.00	85.70\\
3.00	88.29\\
4.00	49.87\\
5.00	55.60\\
6.00	63.42\\
};
\addplot[forget plot, color=mycolor1] table[row sep=crcr] {%
0.54	0.00\\
6.46	0.00\\
};
\addlegendentry{Fall I (\ac{FC})}

\addplot[ybar, bar width=0.071, fill=white!45!black, draw=mycolor1, area legend] table[row sep=crcr] {%
1.00	99.83\\
2.00	88.37\\
3.00	87.60\\
4.00	49.80\\
5.00	63.70\\
6.00	63.28\\
};
\addplot[forget plot, color=mycolor1] table[row sep=crcr] {%
0.54	0.00\\
6.46	0.00\\
};
\addlegendentry{Fall II (leicht)}

\addplot[ybar, bar width=0.071, fill=gray!85!lightgray, draw=mycolor1, area legend] table[row sep=crcr] {%
1.00	99.89\\
2.00	86.60\\
3.00	87.95\\
4.00	54.72\\
5.00	61.28\\
6.00	62.68\\
};
\addplot[forget plot, color=mycolor1] table[row sep=crcr] {%
0.54	0.00\\
6.46	0.00\\
};
\addlegendentry{Fall II (stark)}

\addplot[ybar, bar width=0.071, fill=white!50!darkgray, draw=mycolor1, area legend] table[row sep=crcr] {%
1.00	99.90\\
2.00	87.35\\
3.00	89.15\\
4.00	50.92\\
5.00	62.05\\
6.00	67.67\\
};
\addplot[forget plot, color=mycolor1] table[row sep=crcr] {%
0.54	0.00\\
6.46	0.00\\
};
\addlegendentry{Fall III}

\addplot[ybar, bar width=0.071, fill=black!5!lightgray, draw=mycolor1, area legend] table[row sep=crcr] {%
1.00	99.81\\
2.00	81.44\\
3.00	88.02\\
4.00	50.20\\
5.00	52.31\\
6.00	62.06\\
};
\addplot[forget plot, color=mycolor1] table[row sep=crcr] {%
0.54	0.00\\
6.46	0.00\\
};
\addlegendentry{Fall IV: \ac{SP}}

\addplot[ybar, bar width=0.071, fill=white!80!black, draw=mycolor1, area legend] table[row sep=crcr] {%
1.00	99.78\\
2.00	87.84\\
3.00	87.98\\
4.00	50.63\\
5.00	62.23\\
6.00	63.09\\
};
\addplot[forget plot, color=mycolor1] table[row sep=crcr] {%
0.54	0.00\\
6.46	0.00\\
};
\addlegendentry{Fall IV: \ac{ZP}}

\addplot [color=black, line width=0.4pt, only marks, forget plot]
 plot [error bars/.cd, y dir=both, y explicit, error bar style={line width=0.4pt}, error mark options={line width=0.4pt, mark size=6.0pt, rotate=90}]
 table[row sep=crcr, y error plus index=2, y error minus index=3]{%
0.64	99.83	0.08	0.08\\
1.64	88.04	0.64	0.64\\
2.64	87.77	0.64	0.64\\
3.64	50.24	0.98	0.98\\
4.64	63.33	0.94	0.94\\
5.64	63.82	0.94	0.94\\
};
\addplot [color=black, line width=0.4pt, only marks, forget plot]
 plot [error bars/.cd, y dir=both, y explicit, error bar style={line width=0.4pt}, error mark options={line width=0.4pt, mark size=6.0pt, rotate=90}]
 table[row sep=crcr, y error plus index=2, y error minus index=3]{%
0.73	99.88	0.07	0.07\\
1.73	86.75	0.66	0.66\\
2.73	87.46	0.65	0.65\\
3.73	50.56	0.98	0.98\\
4.73	58.63	0.97	0.97\\
5.73	63.87	0.94	0.94\\
};
\addplot [color=black, line width=0.4pt, only marks, forget plot]
 plot [error bars/.cd, y dir=both, y explicit, error bar style={line width=0.4pt}, error mark options={line width=0.4pt, mark size=6.0pt, rotate=90}]
 table[row sep=crcr, y error plus index=2, y error minus index=3]{%
0.82	99.88	0.07	0.07\\
1.82	88.25	0.63	0.63\\
2.82	88.44	0.63	0.63\\
3.82	50.81	0.98	0.98\\
4.82	69.08	0.91	0.91\\
5.82	63.13	0.95	0.95\\
};
\addplot [color=black, line width=0.4pt, only marks, forget plot]
 plot [error bars/.cd, y dir=both, y explicit, error bar style={line width=0.4pt}, error mark options={line width=0.4pt, mark size=6.0pt, rotate=90}]
 table[row sep=crcr, y error plus index=2, y error minus index=3]{%
0.91	99.84	0.08	0.08\\
1.91	85.70	0.69	0.69\\
2.91	88.29	0.63	0.63\\
3.91	49.87	0.98	0.98\\
4.91	55.60	0.97	0.97\\
5.91	63.42	0.94	0.94\\
};
\addplot [color=black, line width=0.4pt, only marks, forget plot]
 plot [error bars/.cd, y dir=both, y explicit, error bar style={line width=0.4pt}, error mark options={line width=0.4pt, mark size=6.0pt, rotate=90}]
 table[row sep=crcr, y error plus index=2, y error minus index=3]{%
1.00	99.83	0.08	0.08\\
2.00	88.37	0.63	0.63\\
3.00	87.60	0.65	0.65\\
4.00	49.80	0.98	0.98\\
5.00	63.70	0.94	0.94\\
6.00	63.28	0.94	0.94\\
};
\addplot [color=black, line width=0.4pt, only marks, forget plot]
 plot [error bars/.cd, y dir=both, y explicit, error bar style={line width=0.4pt}, error mark options={line width=0.4pt, mark size=6.0pt, rotate=90}]
 table[row sep=crcr, y error plus index=2, y error minus index=3]{%
1.09	99.89	0.06	0.06\\
2.09	86.60	0.67	0.67\\
3.09	87.95	0.64	0.64\\
4.09	54.72	0.98	0.98\\
5.09	61.28	0.95	0.95\\
6.09	62.68	0.95	0.95\\
};
\addplot [color=black, line width=0.4pt, only marks, forget plot]
 plot [error bars/.cd, y dir=both, y explicit, error bar style={line width=0.4pt}, error mark options={line width=0.4pt, mark size=6.0pt, rotate=90}]
 table[row sep=crcr, y error plus index=2, y error minus index=3]{%
1.18	99.90	0.06	0.06\\
2.18	87.35	0.65	0.65\\
3.18	89.15	0.61	0.61\\
4.18	50.92	0.98	0.98\\
5.18	62.05	0.95	0.95\\
6.18	67.67	0.92	0.92\\
};
\addplot [color=black, line width=0.4pt, only marks, forget plot]
 plot [error bars/.cd, y dir=both, y explicit, error bar style={line width=0.4pt}, error mark options={line width=0.4pt, mark size=6.0pt, rotate=90}]
 table[row sep=crcr, y error plus index=2, y error minus index=3]{%
1.27	99.81	0.09	0.09\\
2.27	81.44	0.76	0.76\\
3.27	88.02	0.64	0.64\\
4.27	50.20	0.98	0.98\\
5.27	52.31	0.98	0.98\\
6.27	62.06	0.95	0.95\\
};
\addplot [color=black, line width=0.4pt, only marks, forget plot]
 plot [error bars/.cd, y dir=both, y explicit, error bar style={line width=0.4pt}, error mark options={line width=0.4pt, mark size=6.0pt, rotate=90}]
 table[row sep=crcr, y error plus index=2, y error minus index=3]{%
1.36	99.78	0.09	0.09\\
2.36	87.84	0.64	0.64\\
3.36	87.98	0.64	0.64\\
4.36	50.63	0.98	0.98\\
5.36	62.23	0.95	0.95\\
6.36	63.09	0.95	0.95\\
};
\end{axis}
\end{tikzpicture}%
    \caption{Trennschärfe mit $95\,\%$-Konfidenzintervallen in Abhängigkeit der Störszenarien für $\text{SNR}=1$.}
    \label{fig:app.c_powerdev_snr1}
\end{figure}

\clearpage

\subsection{Effizienzanalysen aus Abschnitt~\ref{sec:ldoe}} \label{app:c_effizienz}

\begin{figure}[htbp]
    \centering
    \setlength\figurewidth{0.8\textwidth}
    \setlength\figureheight{16cm}
    % This file was created by matlab2tikz.
%
\definecolor{mycolor1}{rgb}{0.06600,0.44300,0.74500}%
\definecolor{mycolor2}{rgb}{0.86600,0.32900,0.00000}%
\definecolor{mycolor3}{rgb}{0.92900,0.69400,0.12500}%
\definecolor{mycolor4}{rgb}{0.12941,0.12941,0.12941}%
%
\begin{tikzpicture}

\begin{axis}[%
width=0.958\figurewidth,
height=0.146\figureheight,
at={(0\figurewidth,0.854\figureheight)},
scale only axis,
xmin=46.40,
xmax=72.89,
xtick={50.00,55.00,60.00,65.00,70.00},
xticklabels={\empty},
ymin=0.00,
ymax=6.44,
ylabel style={font=\color{mycolor4}},
ylabel={$\sym{epi}$},
axis background/.style={fill=white},
title style={font=\bfseries\color{mycolor4}},
title={\textbf{Effizienz für $\sym{beta}_1$}},
legend style={at={(0.03,0.97)}, anchor=north west, legend cell align=left, align=left},
grid=major,
grid style={dotted, gray!50}
]
\addplot [color=mycolor1, line width=1.0pt, only marks, mark size=3.2pt, mark=*, mark options={solid, fill=black, draw=black}]
  table[row sep=crcr]{%
56.90	0.62\\
};
\addlegendentry{Basis \ac{CCD}}

\addplot [color=mycolor2, line width=1.0pt, only marks, mark size=1.9pt, mark=square*, mark options={solid, fill=gray, draw=black}]
  table[row sep=crcr]{%
50.40	0.69\\
};
\addlegendentry{$\textnormal{\ac{CCD}}_{-\textnormal{\ac{AP}}}$}

\addplot [color=mycolor3, line width=1.0pt, only marks, mark size=4.7pt, mark=diamond*, mark options={solid, fill=white!20!black, draw=black}]
  table[row sep=crcr]{%
68.89	5.15\\
};
\addlegendentry{\ac{EmALT}}

\end{axis}

\begin{axis}[%
width=0.958\figurewidth,
height=0.146\figureheight,
at={(0\figurewidth,0.64\figureheight)},
scale only axis,
xmin=46.40,
xmax=72.89,
xtick={50.00,55.00,60.00,65.00,70.00},
xticklabels={\empty},
ymin=0.00,
ymax=6.08,
ylabel style={font=\color{mycolor4}},
ylabel={$\sym{epi}$},
axis background/.style={fill=white},
title style={font=\bfseries\color{mycolor4}},
title={\textbf{Effizienz für $\sym{beta}_2$}},
grid=major,
grid style={dotted, gray!50}
]
\addplot [color=mycolor1, line width=1.0pt, only marks, mark size=3.2pt, mark=*, mark options={solid, fill=black, draw=black}, forget plot]
  table[row sep=crcr]{%
56.90	0.62\\
};
\addplot [color=mycolor2, line width=1.0pt, only marks, mark size=1.9pt, mark=square*, mark options={solid, fill=gray, draw=black}, forget plot]
  table[row sep=crcr]{%
50.40	0.69\\
};
\addplot [color=mycolor3, line width=1.0pt, only marks, mark size=4.7pt, mark=diamond*, mark options={solid, fill=white!20!black, draw=black}, forget plot]
  table[row sep=crcr]{%
68.89	4.86\\
};
\end{axis}

\begin{axis}[%
width=0.958\figurewidth,
height=0.146\figureheight,
at={(0\figurewidth,0.427\figureheight)},
scale only axis,
xmin=46.40,
xmax=72.89,
xtick={50.00,55.00,60.00,65.00,70.00},
xticklabels={\empty},
ymin=0.00,
ymax=4.16,
ylabel style={font=\color{mycolor4}},
ylabel={$\sym{epi}$},
axis background/.style={fill=white},
title style={font=\bfseries\color{mycolor4}},
title={\textbf{Effizienz für $\sym{beta}_{12}$}},
grid=major,
grid style={dotted, gray!50}
]
\addplot [color=mycolor1, line width=1.0pt, only marks, mark size=3.2pt, mark=*, mark options={solid, fill=black, draw=black}, forget plot]
  table[row sep=crcr]{%
56.90	0.50\\
};
\addplot [color=mycolor2, line width=1.0pt, only marks, mark size=1.9pt, mark=square*, mark options={solid, fill=gray, draw=black}, forget plot]
  table[row sep=crcr]{%
50.40	0.60\\
};
\addplot [color=mycolor3, line width=1.0pt, only marks, mark size=4.7pt, mark=diamond*, mark options={solid, fill=white!20!black, draw=black}, forget plot]
  table[row sep=crcr]{%
68.89	3.33\\
};
\end{axis}

\begin{axis}[%
width=0.958\figurewidth,
height=0.146\figureheight,
at={(0\figurewidth,0.213\figureheight)},
scale only axis,
xmin=46.40,
xmax=72.89,
xtick={50.00,55.00,60.00,65.00,70.00},
xticklabels={\empty},
ymin=0.00,
ymax=6.38,
ylabel style={font=\color{mycolor4}},
ylabel={$\sym{epi}$},
axis background/.style={fill=white},
title style={font=\bfseries\color{mycolor4}},
title={\textbf{Effizienz für $\sym{beta}_{11}$}},
grid=major,
grid style={dotted, gray!50}
]
\addplot [color=mycolor1, line width=1.0pt, only marks, mark size=3.2pt, mark=*, mark options={solid, fill=black, draw=black}, forget plot]
  table[row sep=crcr]{%
56.90	0.60\\
};
\addplot [color=mycolor2, line width=1.0pt, only marks, mark size=1.9pt, mark=square*, mark options={solid, fill=gray, draw=black}, forget plot]
  table[row sep=crcr]{%
50.40	0.65\\
};
\addplot [color=mycolor3, line width=1.0pt, only marks, mark size=4.7pt, mark=diamond*, mark options={solid, fill=white!20!black, draw=black}, forget plot]
  table[row sep=crcr]{%
68.89	5.11\\
};
\end{axis}

\begin{axis}[%
width=0.958\figurewidth,
height=0.146\figureheight,
at={(0\figurewidth,0\figureheight)},
scale only axis,
xmin=46.40,
xmax=72.89,
xlabel style={font=\color{mycolor4}},
xlabel={$\sym{eff_D}$ [\%]},
ymin=0.00,
ymax=4.84,
ylabel style={font=\color{mycolor4}},
ylabel={$\sym{epi}$},
axis background/.style={fill=white},
title style={font=\bfseries\color{mycolor4}},
title={\textbf{Effizienz für $\sym{beta}_{22}$}},
grid=major,
grid style={dotted, gray!50}
]
\addplot [color=mycolor1, line width=1.0pt, only marks, mark size=3.2pt, mark=*, mark options={solid, fill=black, draw=black}, forget plot]
  table[row sep=crcr]{%
56.90	0.61\\
};
\addplot [color=mycolor2, line width=1.0pt, only marks, mark size=1.9pt, mark=square*, mark options={solid, fill=gray, draw=black}, forget plot]
  table[row sep=crcr]{%
50.40	0.65\\
};
\addplot [color=mycolor3, line width=1.0pt, only marks, mark size=4.7pt, mark=diamond*, mark options={solid, fill=white!20!black, draw=black}, forget plot]
  table[row sep=crcr]{%
68.89	3.87\\
};
\end{axis}

\begin{axis}[%
width=1.236\figurewidth,
height=1.236\figureheight,
at={(-0.161\figurewidth,-0.136\figureheight)},
scale only axis,
xmin=0.00,
xmax=1.00,
ymin=0.00,
ymax=1.00,
axis line style={draw=none},
ticks=none,
axis x line*=bottom,
axis y line*=left,
grid=major,
grid style={dotted, gray!50}
]
\end{axis}
\end{tikzpicture}%
    \caption{\ac{KPI} ggü. \textit{D}-Effizienz aller Koeffizienten.}
    \label{fig:app.c_kpi_all_coeffs}
\end{figure}

\begin{figure}[htbp]
    \centering
    \setlength\figurewidth{0.8\textwidth}
    \setlength\figureheight{5cm}
    % This file was created by matlab2tikz.
%
\definecolor{mycolor1}{rgb}{0.12941,0.12941,0.12941}%
%
\begin{tikzpicture}

\begin{axis}[%
width=\figurewidth,
height=0.969\figureheight,
at={(0\figurewidth,0\figureheight)},
scale only axis,
xmin=0.00,
xmax=1.00,
xlabel style={font=\color{mycolor1}},
xlabel={$\sym{p-Wert}$},
ymin=0.00,
ymax=1.00,
ylabel style={font=\color{mycolor1}},
ylabel={Relative Häufigkeit},
axis background/.style={fill=white},
xmajorgrids,
ymajorgrids,
legend style={legend cell align=left, align=left}
]
\addplot[const plot, color=black, line width=1.5pt] table[row sep=crcr] {%
0.00	0.94\\
0.03	0.01\\
0.05	0.01\\
0.08	0.00\\
0.10	0.00\\
0.13	0.00\\
0.15	0.00\\
0.18	0.00\\
0.21	0.00\\
0.23	0.00\\
0.26	0.00\\
0.28	0.00\\
0.31	0.00\\
0.33	0.00\\
0.36	0.00\\
0.38	0.00\\
0.41	0.00\\
0.44	0.00\\
0.46	0.00\\
0.51	0.00\\
0.54	0.00\\
0.56	0.00\\
0.62	0.00\\
0.64	0.00\\
0.67	0.00\\
0.69	0.00\\
0.72	0.00\\
0.74	0.00\\
0.77	0.00\\
0.79	0.00\\
0.82	0.00\\
0.85	0.00\\
0.87	0.00\\
0.90	0.00\\
0.92	0.00\\
0.95	0.00\\
0.97	0.00\\
1.00	0.00\\
};
\addlegendentry{Basis \ac{CCD}}

\addplot[const plot, color=white!60!black, dotted, line width=1.5pt] table[row sep=crcr] {%
0.00	0.94\\
0.03	0.01\\
0.05	0.01\\
0.08	0.01\\
0.10	0.00\\
0.13	0.00\\
0.15	0.00\\
0.18	0.00\\
0.21	0.00\\
0.23	0.00\\
0.26	0.00\\
0.28	0.00\\
0.31	0.00\\
0.33	0.00\\
0.36	0.00\\
0.38	0.00\\
0.41	0.00\\
0.44	0.00\\
0.46	0.00\\
0.49	0.00\\
0.51	0.00\\
0.54	0.00\\
0.56	0.00\\
0.59	0.00\\
0.62	0.00\\
0.64	0.00\\
0.67	0.00\\
0.69	0.00\\
0.72	0.00\\
0.77	0.00\\
0.82	0.00\\
0.85	0.00\\
0.87	0.00\\
0.90	0.00\\
0.92	0.00\\
0.95	0.00\\
0.97	0.00\\
1.00	0.00\\
};
\addlegendentry{$\textnormal{\ac{CCD}}_{-\textnormal{\ac{AP}}}$}

\addplot[const plot, color=white!30!black, dashed, line width=1.5pt] table[row sep=crcr] {%
0.00	0.85\\
0.03	0.01\\
0.05	0.01\\
0.08	0.01\\
0.10	0.00\\
0.13	0.00\\
0.15	0.00\\
0.18	0.00\\
0.21	0.00\\
0.23	0.00\\
0.26	0.00\\
0.28	0.00\\
0.31	0.00\\
0.36	0.00\\
0.38	0.00\\
0.41	0.00\\
0.44	0.00\\
0.46	0.00\\
0.49	0.00\\
0.54	0.00\\
0.56	0.00\\
0.59	0.00\\
0.62	0.00\\
0.64	0.00\\
0.67	0.00\\
0.72	0.00\\
0.74	0.00\\
0.77	0.00\\
0.79	0.00\\
0.82	0.00\\
0.85	0.00\\
0.87	0.00\\
0.90	0.00\\
0.92	0.00\\
0.95	0.00\\
0.97	0.06\\
1.00	0.06\\
};
\addlegendentry{\ac{EmALT}}

\addplot [color=black, line width=1.5pt, forget plot]
  table[row sep=crcr]{%
0.05	0.00\\
0.05	1.00\\
};
\addlegendentry{Signifikanz $\sym{alpha} = 0.05$}

\end{axis}
\end{tikzpicture}%
    \caption{Aggregierte $\sym{p-Wert}$ Histogramme der \ac{MLE}-Signifikanztests. Die stochastische Auswertung liefert artefaktfreie Verteilungen.}
    \label{fig:ma_abb5.4.3_phisto_all_wbl}
\end{figure}
\begin{table}[htbp]
    \centering
    \caption{Relativer systematischer Schätzfehler (Bias) in Prozent bei GLLW-Schätzung (Weibull-Modell).}
    \label{tab:results_emalt_bias_wbl}
    \renewcommand{\arraystretch}{1.2}
    \small
    \begin{tabular*}{\textwidth}{@{\extracolsep{\fill}}l rrrrrr@{}}
        \toprule
        \textbf{Szenario} & \textbf{$\hat{\sym{beta}}_{0}$} & \textbf{$\hat{\sym{beta}}_{1}$} & \textbf{$\hat{\sym{beta}}_{2}$} & \textbf{$\hat{\sym{beta}}_{12}$} & \textbf{$\hat{\sym{beta}}_{11}$} & \textbf{$\hat{\sym{beta}}_{22}$} \\
        \midrule
        Basis \ac{CCD} & $+0{,}15\,\%$ & $-0{,}08\,\%$ & $+0{,}10\,\%$ & $-0{,}24\,\%$ & $+13{,}73\,\%$ & $+12{,}79\,\%$ \\
        $\textnormal{\ac{CCD}}_{-\textnormal{\ac{AP}}}$ & $+0{,}24\,\%$ & $+3{,}72\,\%$ & $+3{,}39\,\%$ & $+6{,}82\,\%$ & $+23{,}17\,\%$ & $+22{,}10\,\%$ \\
        \ac{EmALT} & $-0{,}00\,\%$ & $+19{,}55\,\%$ & $-39{,}14\,\%$ & $-47{,}84\,\%$ & $+46{,}51\,\%$ & $-53{,}26\,\%$ \\
        \bottomrule
    \end{tabular*}
\end{table}
